\documentclass[11pt,a4paper]{report}
% Compilation : tectonic Cours_Aerodynamique.tex (ou xelatex, deux passes)
% Figures à données : python3 tools/make_figures.py (appelle le code aero du projet et régénère figures/*.tex)
\usepackage{iftex}
\ifPDFTeX
  \usepackage[utf8]{inputenc}
  \usepackage[T1]{fontenc}
  \usepackage{lmodern}
\else
  \usepackage{fontspec}
\fi
\usepackage[french]{babel}
\usepackage[a4paper,margin=2.3cm]{geometry}
\usepackage{amsmath,amssymb,mathtools}
\usepackage{xcolor}
\usepackage{tikz}
\usetikzlibrary{arrows.meta,calc,positioning,shapes.geometric,decorations.pathreplacing,decorations.markings,patterns,babel}
\usepackage{pgfplots}
\pgfplotsset{compat=1.18}
\usepgfplotslibrary{groupplots}
\usepackage{tabularx,booktabs,array}
\usepackage{enumitem}
\setlist{nosep,leftmargin=*}
\usepackage{listings}
\usepackage[font=small,labelfont=bf]{caption}
\usepackage[raggedright]{titlesec}
\usepackage[most]{tcolorbox}
\usepackage[colorlinks=true,linkcolor=blue!50!black,urlcolor=blue!50!black,citecolor=blue!50!black]{hyperref}

\definecolor{cblue}{HTML}{2A78D6}
\definecolor{corange}{HTML}{EB6834}
\definecolor{cgreen}{HTML}{1BAF7A}
\definecolor{cink}{HTML}{1F2328}
\definecolor{cmuted}{HTML}{6B6F76}

% ---------------------------------------------------------------- notations
\newcommand{\dd}{\mathrm{d}}
\newcommand{\Ue}{U_e}
\renewcommand{\Re}{\mathit{Re}}
\newcommand{\vect}[1]{\mathbf{#1}}
\newcommand{\code}[1]{\texttt{#1}}

% ---------------------------------------------------------------- encadrés
\tcbset{coursbox/.style={enhanced,breakable,boxrule=0pt,leftrule=2.5pt,arc=0pt,
  left=6pt,right=6pt,top=4pt,bottom=4pt,fonttitle=\bfseries\small,coltitle=cink,
  attach title to upper={\ }}}
\newtcolorbox{definition}[1][]{coursbox,colback=cblue!5,colframe=cblue,title={Définition#1.}}
\newtcolorbox{resultat}[1][]{coursbox,colback=cgreen!6,colframe=cgreen,title={À retenir#1.}}
\newtcolorbox{projet}[1][]{coursbox,colback=corange!6,colframe=corange,title={Dans notre projet#1.}}
\newtcolorbox{piege}[1][]{coursbox,colback=cmuted!7,colframe=cmuted,title={Piège#1.}}
\newtcolorbox{codebox}[1][]{coursbox,colback=cmuted!3,colframe=cblue!60,title={Dans le code#1.}}

% ---------------------------------------------------------------- exercices
\newcounter{exo}[chapter]
\renewcommand{\theexo}{\thechapter.\arabic{exo}}
\newenvironment{exercice}[1][]{\refstepcounter{exo}\par\medskip\noindent
  \textbf{\color{cblue}Exercice~\theexo}\if\relax\detokenize{#1}\relax\else\ (#1)\fi.\ }{\par}
\newcommand{\exercices}{\section*{Exercices}\addcontentsline{toc}{section}{Exercices}
  \noindent\textit{Corrigés en annexe~\ref{annexe:corriges}. Chercher d'abord sans les regarder.}}

\tikzset{
  flow/.style={draw=cink,rounded corners=2pt,fill=white,align=center,font=\small,inner sep=4pt},
  arr/.style={-{Stealth[length=2.2mm]},cink,thick},
}

\lstset{basicstyle=\ttfamily\footnotesize,breaklines=true,columns=fullflexible,frame=single,rulecolor=\color{cmuted!50},
  backgroundcolor=\color{cmuted!6},keywordstyle=\color{cblue}\bfseries,commentstyle=\color{cmuted}\itshape,
  stringstyle=\color{corange},showstringspaces=false,language=Python,
  literate={é}{\'e}1{è}{\`e}1{à}{\`a}1{ê}{\^e}1{ç}{\c c}1{ô}{\^o}1{î}{\^i}1{ù}{\`u}1{É}{\'E}1{→}{$\to$}1}

\setcounter{secnumdepth}{2}
\setcounter{tocdepth}{1}
\setlength{\parindent}{0pt}
\setlength{\emergencystretch}{3em}
\setlength{\parskip}{0.45em}

% ============================================================================
% Itérations du projet HackathonQuantique : couleurs et marqueurs communs.
% Inclus par tous les cours et par la fiche projet, pour qu'une itération ait
% la même couleur partout. Pour ajouter une itération : une couleur, un nom,
% une date, une ligne dans \legendeiterations.
% ============================================================================
\definecolor{itV0}{HTML}{8D6E63}   % brun   : V0, modèle de la note d'intention
\definecolor{itV1}{HTML}{C2185B}   % framboise : V1, chaîne complète conforme à la note
\definecolor{itV2}{HTML}{00838F}   % cyan foncé : V2, problème reformulé
\definecolor{itV3}{HTML}{E09A00}   % ambre  : V3, concurrent équitable et problème plus grand
\definecolor{itV4}{HTML}{5E35B1}   % indigo : V4, réservé

\makeatletter
\newcommand{\itername}[1]{\@nameuse{itname@#1}}
\newcommand{\iterdate}[1]{\@nameuse{itdate@#1}}
\newcommand{\defiter}[3]{\@namedef{itname@#1}{#2}\@namedef{itdate@#1}{#3}}
\makeatother
\defiter{V0}{Modèle de la note d'intention}{24--30 oct. 2026}
\defiter{V1}{Chaîne complète conforme à la note}{4--5 nov. 2026}
\defiter{V2}{Problème reformulé}{5--6 nov. 2026}
\defiter{V3}{Concurrent équitable, problème plus grand}{à partir du 6 nov. 2026}

% Pastille : \iterbadge{V1}
\newcommand{\iterbadge}[1]{\tikz[baseline=(b.base)]\node[fill=it#1,text=white,rounded corners=2pt,
  inner xsep=3pt,inner ysep=1.5pt,font=\sffamily\bfseries\scriptsize] (b) {#1};}

% Bandeau de début de chapitre : \iterbanner{V0}
\newcommand{\iterbanner}[1]{\par\noindent\begin{tikzpicture}
  \node[fill=it#1!12,draw=it#1,line width=0.6pt,rounded corners=2pt,text width=\dimexpr\linewidth-12pt\relax,
    inner sep=4pt,anchor=west,font=\small] at (0,0)
    {\iterbadge{#1}\ \textbf{\color{it#1}Itération #1} --- \itername{#1} \hfill{\color{it#1}\iterdate{#1}}};
  \end{tikzpicture}\par\medskip}

% Encadré d'évolution : \begin{evolution}{V1}{titre} ... \end{evolution}
\newtcolorbox{evolution}[2]{enhanced,breakable,boxrule=0pt,leftrule=3pt,arc=0pt,
  left=6pt,right=6pt,top=4pt,bottom=4pt,colback=it#1!7,colframe=it#1,
  fonttitle=\bfseries\small,coltitle=it#1,attach title to upper={\ },
  title={\iterbadge{#1}\ Évolution #1 : #2.}}

% Légende des itérations, pour les avant-propos
\newcommand{\legendeiterations}{%
\begin{center}\small
\begin{tabular}{lll}
\iterbadge{V0} & \itername{V0} & \iterdate{V0}\\
\iterbadge{V1} & \itername{V1} & \iterdate{V1}\\
\iterbadge{V2} & \itername{V2} & \iterdate{V2}\\
\iterbadge{V3} & \itername{V3} & \iterdate{V3}\\
\end{tabular}
\end{center}}

% Encadré « pistes d'amélioration » : défini ici s'il manque (cours V0 sans cet encadré)
\providecolor{cviolet}{HTML}{7A4FB3}
\ifcsname pistes\endcsname\else
\newtcolorbox{pistes}[1][]{enhanced,breakable,boxrule=0pt,leftrule=2.5pt,arc=0pt,
  left=6pt,right=6pt,top=4pt,bottom=4pt,fonttitle=\bfseries\small,coltitle=cink,attach title to upper={\ },
  colback=cviolet!6,colframe=cviolet,title={Pistes d'amélioration#1.}}
\fi


\title{\textbf{Le modèle aérodynamique de la phase 1}\\[0.4em]\large Panneaux de Hess-Smith, couche limite intégrale et traînée de Squire-Young :\\ théorie et code du simulateur du hackathon quantique}
\author{Notes de travail, projet \emph{HackathonQuantique}}
\date{Version du 3 novembre 2026}

\begin{document}
\pagenumbering{Alph}
\maketitle
\pagenumbering{roman}

\chapter*{Avant-propos}
\addcontentsline{toc}{chapter}{Avant-propos}
\paragraph{Objet.} Ce cours explique, de la théorie jusqu'à la ligne de code, le modèle aérodynamique de la phase~1 du hackathon : la fonction qui, à un profil NACA à quatre chiffres $(m,p,t)$, associe sa portance, sa traînée, son moment et un diagnostic de validité. C'est l'« oracle » que les phases suivantes appellent : l'énumération exhaustive (front de Pareto exact), NSGA-II, et la boucle QUBO/QAOA décrite dans le cours compagnon \emph{Calcul quantique pour l'optimisation}.

\paragraph{Ce qu'on suppose connu.} Les équations de Navier-Stokes, la notion de couche limite et ses grandeurs intégrales, la turbulence de paroi. Ces bases sont rappelées en quelques lignes, sans démonstration. En revanche, tout ce qui est propre au modèle est démontré ou justifié : les intégrales de panneau, la condition de Kutta discrète, la quadrature de Thwaites, l'équation d'entraînement de Head, l'exposant de Squire-Young, et chaque correction d'ingénierie (\code{x\_stop}, \code{x\_sep\_crit}, interpolation de la transition).

\paragraph{Fil conducteur.} Le rôle du modèle et ses exigences (ch.~\ref{ch:intro}) $\to$ la géométrie NACA (ch.~\ref{ch:naca}) $\to$ l'écoulement potentiel et ses singularités (ch.~\ref{ch:potentiel}) $\to$ la méthode des panneaux de Hess-Smith (ch.~\ref{ch:panneaux}) $\to$ la couche limite intégrale : principe et découpage (ch.~\ref{ch:cl}), laminaire par Thwaites (ch.~\ref{ch:thwaites}), transition (ch.~\ref{ch:transition}), turbulent par Head (ch.~\ref{ch:head}) $\to$ la traînée par Squire-Young (ch.~\ref{ch:trainee}) $\to$ l'assemblage dans \code{evaluate} (ch.~\ref{ch:modele}) $\to$ l'espace de conception et les interfaces avec les autres phases (ch.~\ref{ch:espace}) $\to$ la validation et les limites (ch.~\ref{ch:validation}).

\paragraph{Figures et chiffres.} Toutes les courbes à données de ce cours sont \textbf{calculées par le code du projet} (dossier \texttt{Code/aero/} du projet \emph{HackathonQuantique}), via le script \texttt{tools/make\_figures.py} de ce dossier. Elles sont tracées en pgfplots, et les nombres cités dans le texte sont dans \texttt{figures/chiffres.json}. Si le code change, on relance le script et on recompile. Les schémas de principe sont en TikZ.

\paragraph{Encadrés.} Bleu : définitions. Vert : à retenir. Gris : pièges. Bleu clair \emph{Dans le code} : où et comment la notion est implémentée. Orange \emph{Dans notre projet} : lien avec l'optimisation et les autres phases. Chaque chapitre se termine par un ou deux exercices, corrigés en annexe~\ref{annexe:corriges}. L'annexe~\ref{annexe:memo} résume formules et API sur une page.

\paragraph{Conventions.} Tout est adimensionné par la corde $c$ et la vitesse amont $U_\infty$ : $x,y,s$ sont en cordes, les vitesses sont des rapports $U/U_\infty$, et $\nu=1/\Re$ avec $\Re=U_\infty c/\nu$. Le contour est parcouru comme dans XFOIL : bord de fuite extrados $\to$ bord d'attaque $\to$ bord de fuite intrados. Un \emph{count} de traînée vaut $10^{-4}$.

\paragraph{Itérations du projet.} Le projet avance par itérations, chacune avec sa couleur dans tous les cours et dans la fiche projet (dossier \texttt{Fiche\_Projet/} du projet). Chaque chapitre commence par un bandeau qui indique l'itération de son contenu ; un ajout ou une révision faite lors d'une itération ultérieure apparaît dans un encadré « Évolution » de la couleur correspondante.
\legendeiterations

\tableofcontents
\cleardoublepage
\pagenumbering{arabic}

\chapter{Le rôle du modèle et ses exigences}\label{ch:intro}
\iterbanner{V0}

\section{Une fonction objectif, pas une simulation de référence}
Le hackathon cherche des profils qui réalisent le meilleur compromis entre traînée aérodynamique et surface équivalente radar (SER). Il compare pour cela trois façons d'explorer l'espace de conception : l'énumération exhaustive (le front de Pareto exact), un algorithme génétique (NSGA-II) et une boucle quantique (modèle quadratique, QUBO, QAOA ou recuit). Les trois ont besoin de la même chose : une fonction
\[
(m,p,t)\ \longmapsto\ \big(C_L,\ C_D,\ C_M,\ x_{tr},\ \text{validité}\big)\ +\ \text{géométrie},
\]
qu'on appellera des centaines à des milliers de fois. La phase~1 construit cette fonction.

Ce contexte change complètement la façon de juger un modèle aérodynamique. La question n'est pas « quel est le $C_D$ le plus juste ? », mais « quel modèle permet de \emph{classer} correctement des centaines de profils, vite, sans jamais planter, et avec une sortie assez régulière pour qu'un modèle quadratique puisse l'approcher ? ».

\begin{center}
\renewcommand{\arraystretch}{1.2}
\begin{tabularx}{\linewidth}{lXX}
\toprule
Exigence & Pourquoi & Conséquence de conception \\
\midrule
\textbf{Rapide} & 512 profils pour l'énumération, des milliers d'appels pour NSGA-II avec plusieurs graines & Pas de RANS : panneaux et couche limite intégrale, environ 8~ms par profil \\
\textbf{Robuste} & Un optimiseur génétique propose des formes extrêmes (7~\% de cambrure, 20~\% d'épaisseur) & Pas de couplage itératif qui peut diverger ; des drapeaux de validité plutôt que des exceptions \\
\textbf{Déterministe et régulier} & La boucle quantique ajuste un polynôme quadratique sur les sorties & Aucun bruit de convergence ; les seules irrégularités sont physiques (transition, décollement) \\
\textbf{Cohérent} & Les trois méthodes doivent voir \emph{le même} problème & Une seule grille, un seul point de vol, un seul compteur de simulations \\
\bottomrule
\end{tabularx}
\end{center}

\section{La décomposition de Prandtl, mise en algorithme}
À grand nombre de Reynolds et sans décollement étendu, l'écoulement autour d'un profil se sépare en deux régions (figure~\ref{fig:prandtl}) :
\begin{itemize}
\item un \textbf{écoulement extérieur non visqueux}, irrotationnel, qui fixe la distribution de pression, donc la portance et le moment ;
\item une \textbf{couche limite} mince, plus le sillage qui la prolonge, où la viscosité agit, et qui fixe la traînée.
\end{itemize}
Le lien entre les deux est la vitesse à la frontière de la couche limite, $U_e(s)$. On calcule d'abord l'écoulement extérieur comme si le fluide était parfait, on en tire $U_e(s)$, puis on résout la couche limite sur chaque face avec cette vitesse imposée. C'est un \textbf{couplage faible} : la couche limite ne rétroagit pas sur l'écoulement extérieur. Le chapitre~\ref{ch:validation} discute ce que ce choix fait perdre.

\begin{figure}[htbp]
\centering
\begin{tikzpicture}[x=9cm,y=9cm]
% profil (NACA 2412 approx.)
\fill[cmuted!25,draw=cink,thick] plot[smooth,tension=0.6] coordinates {(1,0) (0.8,0.035) (0.5,0.07) (0.3,0.08) (0.12,0.065) (0.03,0.035) (0,0) (0.03,-0.022) (0.12,-0.035) (0.3,-0.04) (0.5,-0.033) (0.8,-0.014) (1,0)};
% couche limite (épaisseur exagérée)
\draw[cblue,thick,dashed] plot[smooth,tension=0.6] coordinates {(0.0,0.003) (0.03,0.04) (0.12,0.073) (0.3,0.093) (0.5,0.088) (0.8,0.065) (1,0.035)};
\draw[cblue,thick,dashed] plot[smooth,tension=0.6] coordinates {(0.0,-0.003) (0.03,-0.027) (0.12,-0.043) (0.3,-0.052) (0.5,-0.049) (0.8,-0.035) (1,-0.022)};
% sillage
\fill[cblue!10] (1,0.035) .. controls (1.15,0.03) and (1.3,0.025) .. (1.45,0.022) -- (1.45,-0.018) .. controls (1.3,-0.02) and (1.15,-0.021) .. (1,-0.022) -- cycle;
\draw[cblue,thick,dashed] (1,0.035) .. controls (1.15,0.03) and (1.3,0.025) .. (1.45,0.022);
\draw[cblue,thick,dashed] (1,-0.022) .. controls (1.15,-0.021) and (1.3,-0.02) .. (1.45,-0.018);
% lignes de courant extérieures
\foreach \y in {0.2,0.14,-0.12,-0.18}{
  \draw[cmuted,-{Stealth[length=1.6mm]}] (-0.25,\y) .. controls (0.2,\y+0.25*\y) and (0.8,\y+0.1*\y) .. (1.45,\y);}
\draw[-{Stealth[length=2.5mm]},cink,very thick] (-0.33,0) -- (-0.12,0) node[midway,above,font=\small] {$U_\infty$};
\node[font=\small,align=center,cmuted] at (0.55,0.27) {écoulement extérieur non visqueux : $\nabla^2\phi=0$\\ $\Rightarrow$ pression, $C_L$, $C_M$, $U_e(s)$};
\node[font=\small,cblue,align=left,anchor=west] at (0.62,0.13) {couche limite :\\ $\theta(s)$, $H(s)$, $c_f(s)$};
\node[font=\small,cblue,align=center] at (1.27,-0.08) {sillage\\ $\Rightarrow C_D$};
\end{tikzpicture}
\caption{La décomposition de Prandtl qui structure le modèle. Épaisseurs de couche limite très exagérées : à $\Re=3\cdot10^6$, $\delta$ fait environ 1~\% de la corde au bord de fuite.}\label{fig:prandtl}
\end{figure}

\section{La chaîne de calcul}
La figure~\ref{fig:chaine} montre les étapes, et le tableau qui suit associe chacune à son module Python.

\begin{figure}[htbp]
\centering
\begin{tikzpicture}[node distance=0.55cm]
\node[flow] (in) {$(m,p,t)$, $\alpha$ ou $C_L^\star$, $\Re$};
\node[flow,below=of in] (geo) {\textbf{Géométrie NACA} (ch.~\ref{ch:naca})\\ $2n+1$ points, distribution en cosinus};
\node[flow,below=of geo] (pan) {\textbf{Panneaux de Hess-Smith} (ch.~\ref{ch:panneaux})\\ $N$ sources $\sigma_j$, un tourbillon $\gamma$, condition de Kutta};
\node[flow,below=of pan] (stag) {\textbf{Point d'arrêt, deux faces} (ch.~\ref{ch:cl})\\ $U_e(s)=|V_t|$ sur l'extrados et l'intrados};
\node[flow,below=of stag] (bl) {\textbf{Couche limite intégrale}, sur chaque face\\
Thwaites (ch.~\ref{ch:thwaites}) $\to$ transition (ch.~\ref{ch:transition}) $\to$ Head (ch.~\ref{ch:head}), arrêt en $x_{stop}$};
\node[flow,below=of bl] (sy) {\textbf{Squire-Young} (ch.~\ref{ch:trainee})\\ $C_D=\sum_{\text{faces}}2\,\theta_{TE}\,U_{e,TE}^{(H_{TE}+5)/2}$};
\node[flow,below=of sy] (out) {\code{AeroResult} (ch.~\ref{ch:modele})\\ $C_L$, $C_D$, $C_{D,f}$, $C_M$, $x_{tr}$, $x_{sep}$, $r_{BA}$, validité};
\draw[arr] (in) -- (geo); \draw[arr] (geo) -- (pan);
\draw[arr] (pan) -- node[right,font=\scriptsize] {$V_t$, $C_p$, $C_L$, $C_M$} (stag);
\draw[arr] (stag) -- (bl);
\draw[arr] (bl) -- node[right,font=\scriptsize] {$\theta_{TE}$, $H_{TE}$, $U_{e,TE}$} (sy);
\draw[arr] (sy) -- (out);
\node[flow,right=0.6cm of geo,text width=2.6cm,fill=corange!8] (stl) {Export STL\\ vers la SER (phase~2)};
\draw[arr,corange] (geo) -- (stl);
\end{tikzpicture}
\caption{La chaîne de calcul du modèle aérodynamique, telle qu'enchaînée par \code{aero.model.evaluate}.}\label{fig:chaine}
\end{figure}

\begin{center}
\begin{tabularx}{\linewidth}{lXl}
\toprule
Fichier (\texttt{Code/aero/}) & Contenu & Chapitre\\
\midrule
\code{naca4.py} & ligne moyenne, épaisseur, points du contour, rayon de bord d'attaque & \ref{ch:naca}\\
\code{panel.py} & panneaux, matrices d'influence, système de Hess-Smith, $C_p$, $C_L$, $C_M$ & \ref{ch:panneaux}\\
\code{boundary\_layer.py} & point d'arrêt, Thwaites, transition, Head, Squire-Young & \ref{ch:cl}--\ref{ch:trainee}\\
\code{model.py} & \code{evaluate}, \code{AeroResult}, \code{alpha\_for\_cl} & \ref{ch:modele}\\
\code{design\_space.py} & grille $8\times8\times8$, encodage binaire, \code{AeroSimulator} & \ref{ch:espace}\\
\code{geometry\_export.py} & extrusion 2,5D et STL pour la SER & \ref{ch:espace}\\
\code{xfoil\_runner.py} & pilotage d'XFOIL pour la validation & \ref{ch:validation}\\
\bottomrule
\end{tabularx}
\end{center}
Trois scripts les utilisent : \code{demo\_profil.py} (un profil en détail), \code{validation.py} (tests V1 à V5) et \code{generer\_base.py} (les 512 profils de la grille).

\section{Pourquoi pas RANS, pourquoi pas XFOIL}
\paragraph{RANS.} Un calcul RANS 2D d'un profil coûte de quelques minutes à une heure. L'énumération exhaustive prendrait des jours, et NSGA-II avec 10 graines des semaines. C'est l'outil de l'étape industrielle, pas du démonstrateur. La note d'intention le dit explicitement : seuls le simulateur et le nombre de variables changent pour passer à l'échelle.

\paragraph{XFOIL.} XFOIL résout la même physique que nous, mieux : couplage fort, transition $e^N$, sillage. Mais comme boîte noire d'optimisation, il a deux défauts. Il ne converge pas toujours, surtout aux formes extrêmes et à forte incidence, et ses échecs sont silencieux : une incidence non convergée disparaît simplement de la polaire. Il sert donc ici de \textbf{référence de validation}, pas d'oracle (ch.~\ref{ch:validation}).

\begin{projet}
Sur la machine de développement, \code{evaluate} prend environ 8~ms par profil (NACA 2412, $\alpha=2^\circ$, 200 panneaux). Les 512 profils de la grille sont donc calculés en quelques secondes. C'est ce qui rend possible le calcul du front de Pareto \emph{exact}, auquel on compare ensuite NSGA-II et la boucle quantique en nombre de simulations.
\end{projet}

\exercices
\begin{exercice}
NSGA-II est lancé avec une population de 40 individus pendant 50 générations, pour 10 graines aléatoires. (a) Majorer le nombre d'appels au modèle. (b) Combien de temps cela prend-il avec notre modèle (8~ms) ? Avec un RANS à 30~min par profil ? (c) Pourquoi le nombre de simulations \emph{distinctes} est-il en réalité bien plus petit ?
\end{exercice}
\begin{exercice}
Citer deux raisons pour lesquelles un modèle plus précis mais qui échoue de temps en temps est un mauvais choix pour une énumération exhaustive suivie d'un ajustement quadratique.
\end{exercice}

\chapter{La géométrie : profils NACA à quatre chiffres}\label{ch:naca}
\iterbanner{V0}

\section{Trois paramètres, deux lois}
Un profil NACA $MPXX$ superpose une \textbf{ligne moyenne} (la cambrure) et une \textbf{loi d'épaisseur} symétrique. Trois nombres le définissent, et ce sont les trois variables de conception du hackathon :
\begin{itemize}
\item $m=M/100$ : la cambrure maximale, en fraction de corde ;
\item $p=P/10$ : l'abscisse de cette cambrure maximale ;
\item $t=XX/100$ : l'épaisseur relative maximale.
\end{itemize}
Le NACA 2412 a donc $m=0{,}02$, $p=0{,}4$ et $t=0{,}12$.

\paragraph{La ligne moyenne.} Elle est formée de deux arcs de parabole qui se raccordent en $x=p$, au sommet de la cambrure :
\[
y_c(x)=\begin{cases}\dfrac{m}{p^2}\big(2px-x^2\big) & 0\le x<p,\\[8pt]
\dfrac{m}{(1-p)^2}\big((1-2p)+2px-x^2\big) & p\le x\le1,\end{cases}
\qquad
\frac{\dd y_c}{\dd x}=\begin{cases}\dfrac{2m}{p^2}(p-x) & x<p,\\[8pt] \dfrac{2m}{(1-p)^2}(p-x) & x\ge p.\end{cases}
\]
Les deux arcs valent $m$ en $x=p$, avec une pente nulle : la ligne est continue et dérivable. Elle s'annule aux deux bords ($y_c(0)=y_c(1)=0$). Sa courbure, en revanche, saute en $x=p$, de $-2m/p^2$ à $-2m/(1-p)^2$.

\paragraph{La loi d'épaisseur.} C'est un polynôme calé sur des profils de bonne qualité dans les années 1930 :
\[
y_t(x)=5t\big(0{,}2969\sqrt x-0{,}1260\,x-0{,}3516\,x^2+0{,}2843\,x^3+a_4\,x^4\big).
\]
Le coefficient d'origine $a_4=-0{,}1015$ laisse un bord de fuite d'épaisseur finie ($y_t(1)=0{,}0021\cdot5t$). Le code utilise $a_4=-0{,}1036$, qui donne $y_t(1)=0$ : un \textbf{bord de fuite fermé}. C'est nécessaire pour un contour de panneaux fermé et pour un STL étanche. Le terme en $\sqrt x$ donne un bord d'attaque arrondi, de rayon
\[
r_{BA}=1{,}1019\,t^2 .
\]
L'épaisseur maximale vaut $t$ et se trouve vers $x\approx0{,}30$, \emph{quelle que soit} la cambrure.

\paragraph{L'assemblage.} L'épaisseur est appliquée \emph{perpendiculairement} à la ligne moyenne (figure~\ref{fig:naca-construction}). Avec $\vartheta=\arctan(\dd y_c/\dd x)$ :
\begin{align*}
x_{u}&=x-y_t\sin\vartheta, & y_{u}&=y_c+y_t\cos\vartheta & &\text{(extrados)},\\
x_{l}&=x+y_t\sin\vartheta, & y_{l}&=y_c-y_t\cos\vartheta & &\text{(intrados)}.
\end{align*}

\begin{figure}[htbp]
\centering
\begin{tikzpicture}[x=10.5cm,y=10.5cm]
\draw[cmuted,thin] (0,0) -- (1,0) node[right,font=\small] {corde};
\draw[cblue,thick,dashed] plot[domain=0:0.4,samples=30] (\x,{0.05/0.16*(0.8*\x-\x*\x)}) -- plot[domain=0.4:1,samples=40] (\x,{0.05/0.36*(0.2+0.8*\x-\x*\x)});
\node[cblue,font=\small] at (0.75,0.055) {ligne moyenne $y_c$};
% point de construction
\def\xx{0.22}
\pgfmathsetmacro{\yc}{0.05/0.16*(0.8*\xx-\xx*\xx)}
\pgfmathsetmacro{\sl}{2*0.05/0.16*(0.4-\xx)}
\pgfmathsetmacro{\th}{atan(\sl)}
\pgfmathsetmacro{\yt}{5*0.12*(0.2969*sqrt(\xx)-0.1260*\xx-0.3516*\xx*\xx+0.2843*\xx*\xx*\xx-0.1036*\xx*\xx*\xx*\xx)}
\fill[cink] (\xx,\yc) circle (0.2pt);
\draw[corange,very thick,-{Stealth[length=2mm]}] (\xx,\yc) -- ++(90+\th:\yt) node[above left,font=\small] {$y_t$};
\draw[corange,very thick,-{Stealth[length=2mm]}] (\xx,\yc) -- ++(-90+\th:\yt) node[below right,font=\small] {$y_t$};
\draw[cink,thin] (\xx,\yc) -- ++(\th:0.12);
\draw[cink,thin] (\xx,\yc) -- ++(0:0.12);
\draw[cink] (\xx+0.07,\yc) arc (0:\th:0.07);
\node[font=\small] at (\xx+0.105,\yc+0.012) {$\vartheta$};
% profil complet
\draw[cink,thick] plot[smooth,tension=0.5] coordinates {(1,0) (0.8,0.032) (0.6,0.058) (0.4,0.077) (0.25,0.08) (0.12,0.07) (0.04,0.042) (0,0) (0.04,-0.026) (0.12,-0.036) (0.25,-0.039) (0.4,-0.034) (0.6,-0.022) (0.8,-0.011) (1,0)};
\draw[cmuted,{Stealth[length=1.5mm]}-{Stealth[length=1.5mm]}] (0,-0.075) -- (0.4,-0.075) node[midway,below,font=\small] {$p$};
\draw[cmuted,{Stealth[length=1.5mm]}-{Stealth[length=1.5mm]}] (0.4,0) -- (0.4,0.05) node[midway,right,font=\small] {$m$};
\end{tikzpicture}
\caption{Construction d'un profil NACA (cambrure exagérée) : l'épaisseur $y_t$ est reportée de part et d'autre de la ligne moyenne, perpendiculairement à celle-ci.}\label{fig:naca-construction}
\end{figure}

\section{Ce que chaque paramètre fait physiquement}
\begin{figure}[htbp]
\centering
% figure générée par tools/make_figures.py : parametres
\begin{tikzpicture}
\begin{groupplot}[group style={group size=1 by 3,vertical sep=1.5cm,x descriptions at=edge bottom},width=\linewidth,height=3.9cm,
  axis equal image,xmin=-0.02,xmax=1.02,ymin=-0.12,ymax=0.17,xlabel={$x/c$},ylabel={$y/c$},
  title style={font=\small},legend columns=3,legend style={draw=none,font=\scriptsize,at={(0.99,0.97)},anchor=north east,fill=white,fill opacity=0.8,text opacity=1}]
\nextgroupplot[title={effet de $m$ ($p=0{,}4$, $t=12$~\%)}]
\addplot[cmuted,thick] coordinates {(1,-1.6653e-17) (0.99931,9.9554e-05) (0.99726,0.00039741) (0.99384,0.00089119) (0.98907,0.0015769) (0.98296,0.0024493) (0.97553,0.0035014) (0.96679,0.0047252) (0.95677,0.0061118) (0.9455,0.0076508) (0.93301,0.0093315) (0.91934,0.011142) (0.90451,0.013071) (0.88857,0.015105) (0.87157,0.017232) (0.85355,0.019438) (0.83457,0.021712) (0.81466,0.024038) (0.79389,0.026405) (0.77232,0.028798) (0.75,0.031204) (0.727,0.03361) (0.70337,0.036002) (0.67918,0.038366) (0.65451,0.040686) (0.62941,0.042949) (0.60396,0.04514) (0.57822,0.047242) (0.55226,0.049241) (0.52617,0.051119) (0.5,0.052862) (0.47383,0.054451) (0.44774,0.055872) (0.42178,0.057108) (0.39604,0.058144) (0.37059,0.058965) (0.34549,0.059557) (0.32082,0.059907) (0.29663,0.060005) (0.273,0.059841) (0.25,0.059407) (0.22768,0.058699) (0.20611,0.057712) (0.18534,0.056446) (0.16543,0.054901) (0.14645,0.053083) (0.12843,0.050995) (0.11143,0.048648) (0.095492,0.046049) (0.080665,0.04321) (0.066987,0.040145) (0.054497,0.036867) (0.043227,0.033389) (0.03321,0.029726) (0.024472,0.025893) (0.017037,0.021904) (0.010926,0.01777) (0.0061558,0.013503) (0.0027391,0.0091145) (0.00068523,0.0046113) (0,0) (0.00068523,-0.0046113) (0.0027391,-0.0091145) (0.0061558,-0.013503) (0.010926,-0.01777) (0.017037,-0.021904) (0.024472,-0.025893) (0.03321,-0.029726) (0.043227,-0.033389) (0.054497,-0.036867) (0.066987,-0.040145) (0.080665,-0.04321) (0.095492,-0.046049) (0.11143,-0.048648) (0.12843,-0.050995) (0.14645,-0.053083) (0.16543,-0.054901) (0.18534,-0.056446) (0.20611,-0.057712) (0.22768,-0.058699) (0.25,-0.059407) (0.273,-0.059841) (0.29663,-0.060005) (0.32082,-0.059907) (0.34549,-0.059557) (0.37059,-0.058965) (0.39604,-0.058144) (0.42178,-0.057108) (0.44774,-0.055872) (0.47383,-0.054451) (0.5,-0.052862) (0.52617,-0.051119) (0.55226,-0.049241) (0.57822,-0.047242) (0.60396,-0.04514) (0.62941,-0.042949) (0.65451,-0.040686) (0.67918,-0.038366) (0.70337,-0.036002) (0.727,-0.03361) (0.75,-0.031204) (0.77232,-0.028798) (0.79389,-0.026405) (0.81466,-0.024038) (0.83457,-0.021712) (0.85355,-0.019438) (0.87157,-0.017232) (0.88857,-0.015105) (0.90451,-0.013071) (0.91934,-0.011142) (0.93301,-0.0093315) (0.9455,-0.0076508) (0.95677,-0.0061118) (0.96679,-0.0047252) (0.97553,-0.0035014) (0.98296,-0.0024493) (0.98907,-0.0015769) (0.99384,-0.00089119) (0.99726,-0.00039741) (0.99931,-9.9554e-05) (1,1.6653e-17)};\addlegendentry{$m=0$}
\addplot[cmuted,dashed,thin,forget plot] coordinates {(0,0) (0.016949,0) (0.033898,0) (0.050847,0) (0.067797,0) (0.084746,0) (0.10169,0) (0.11864,0) (0.13559,0) (0.15254,0) (0.16949,0) (0.18644,0) (0.20339,0) (0.22034,0) (0.23729,0) (0.25424,0) (0.27119,0) (0.28814,0) (0.30508,0) (0.32203,0) (0.33898,0) (0.35593,0) (0.37288,0) (0.38983,0) (0.40678,0) (0.42373,0) (0.44068,0) (0.45763,0) (0.47458,0) (0.49153,0) (0.50847,0) (0.52542,0) (0.54237,0) (0.55932,0) (0.57627,0) (0.59322,0) (0.61017,0) (0.62712,0) (0.64407,0) (0.66102,0) (0.67797,0) (0.69492,0) (0.71186,0) (0.72881,0) (0.74576,0) (0.76271,0) (0.77966,0) (0.79661,0) (0.81356,0) (0.83051,0) (0.84746,0) (0.86441,0) (0.88136,0) (0.89831,0) (0.91525,0) (0.9322,0) (0.94915,0) (0.9661,0) (0.98305,0) (1,0)};
\addplot[cblue,very thick] coordinates {(1,-1.6571e-17) (0.99932,0.00016755) (0.9973,0.00066874) (0.99393,0.0014993) (0.98923,0.0026521) (0.9832,0.0041173) (0.97586,0.0058826) (0.96723,0.0079334) (0.95734,0.010253) (0.9462,0.012822) (0.93384,0.01562) (0.9203,0.018625) (0.9056,0.021814) (0.8898,0.025163) (0.87292,0.028647) (0.85502,0.032241) (0.83613,0.035918) (0.81632,0.039652) (0.79562,0.043419) (0.7741,0.047191) (0.75182,0.050943) (0.72882,0.05465) (0.70519,0.058287) (0.68097,0.061829) (0.65623,0.065252) (0.63105,0.068532) (0.60549,0.071647) (0.57962,0.074575) (0.55351,0.077293) (0.52724,0.079781) (0.50088,0.082021) (0.4745,0.083993) (0.44818,0.08568) (0.42199,0.087068) (0.39596,0.088141) (0.36994,0.088799) (0.34427,0.088987) (0.31904,0.088705) (0.29431,0.087957) (0.27016,0.086749) (0.24666,0.085095) (0.2239,0.083009) (0.20192,0.080511) (0.18081,0.077624) (0.16062,0.074374) (0.14142,0.07079) (0.12326,0.066905) (0.10619,0.062751) (0.090267,0.058365) (0.075527,0.053784) (0.062013,0.049042) (0.04976,0.044179) (0.0388,0.039228) (0.029159,0.034224) (0.020861,0.029199) (0.013923,0.024182) (0.0083607,0.0192) (0.0041829,0.014275) (0.0013961,0.0094244) (2.3425e-06,0.0046631) (0,0) (0.0013681,-0.0044577) (0.004082,-0.0086055) (0.0081288,-0.012442) (0.013492,-0.015967) (0.020151,-0.01918) (0.028082,-0.022082) (0.03726,-0.024674) (0.047655,-0.02696) (0.059234,-0.028943) (0.071962,-0.030629) (0.085802,-0.032024) (0.10072,-0.033138) (0.11666,-0.033979) (0.13359,-0.034561) (0.15147,-0.034899) (0.17025,-0.035007) (0.18987,-0.034903) (0.21029,-0.034609) (0.23147,-0.034144) (0.25334,-0.033532) (0.27585,-0.032797) (0.29896,-0.031964) (0.32259,-0.031057) (0.34671,-0.030101) (0.37124,-0.029123) (0.39613,-0.028147) (0.42158,-0.027147) (0.44729,-0.02606) (0.47316,-0.024901) (0.49912,-0.023687) (0.52509,-0.022434) (0.55102,-0.021157) (0.57681,-0.019868) (0.60242,-0.01858) (0.62777,-0.017304) (0.65278,-0.016048) (0.6774,-0.014819) (0.70155,-0.013626) (0.72517,-0.012471) (0.74818,-0.01136) (0.77054,-0.010294) (0.79216,-0.0092772) (0.813,-0.0083095) (0.833,-0.0073921) (0.85209,-0.0065257) (0.87022,-0.0057107) (0.88735,-0.0049472) (0.90341,-0.0042357) (0.91837,-0.0035765) (0.93219,-0.0029701) (0.94481,-0.0024172) (0.95621,-0.0019186) (0.96635,-0.0014752) (0.97519,-0.0010881) (0.98273,-0.00075825) (0.98892,-0.00048672) (0.99376,-0.00027443) (0.99722,-0.00012218) (0.9993,-3.0577e-05) (1,1.6571e-17)};\addlegendentry{$m=3~\%$}
\addplot[cblue,dashed,thin,forget plot] coordinates {(0,0) (0.016949,0.0024885) (0.033898,0.0048693) (0.050847,0.0071423) (0.067797,0.0093077) (0.084746,0.011365) (0.10169,0.013315) (0.11864,0.015157) (0.13559,0.016892) (0.15254,0.018518) (0.16949,0.020037) (0.18644,0.021449) (0.20339,0.022752) (0.22034,0.023948) (0.23729,0.025036) (0.25424,0.026016) (0.27119,0.026889) (0.28814,0.027654) (0.30508,0.028311) (0.32203,0.02886) (0.33898,0.029302) (0.35593,0.029636) (0.37288,0.029862) (0.38983,0.029981) (0.40678,0.029996) (0.42373,0.029953) (0.44068,0.029862) (0.45763,0.029723) (0.47458,0.029537) (0.49153,0.029302) (0.50847,0.029019) (0.52542,0.028689) (0.54237,0.028311) (0.55932,0.027885) (0.57627,0.027411) (0.59322,0.026889) (0.61017,0.026319) (0.62712,0.025701) (0.64407,0.025036) (0.66102,0.024323) (0.67797,0.023561) (0.69492,0.022752) (0.71186,0.021895) (0.72881,0.02099) (0.74576,0.020037) (0.76271,0.019037) (0.77966,0.017988) (0.79661,0.016892) (0.81356,0.015747) (0.83051,0.014555) (0.84746,0.013315) (0.86441,0.012027) (0.88136,0.010691) (0.89831,0.0093077) (0.91525,0.0078761) (0.9322,0.0063966) (0.94915,0.0048693) (0.9661,0.0032941) (0.98305,0.001671) (1,0)};
\addplot[corange,very thick] coordinates {(1,-1.633e-17) (0.99933,0.00023459) (0.99734,0.00093632) (0.99402,0.0020991) (0.98938,0.0037127) (0.98343,0.0057633) (0.97619,0.0082332) (0.96767,0.011101) (0.95789,0.014343) (0.94687,0.017932) (0.93465,0.021837) (0.92124,0.026027) (0.90668,0.030468) (0.891,0.035125) (0.87425,0.039959) (0.85646,0.044935) (0.83768,0.050013) (0.81795,0.055155) (0.79733,0.060321) (0.77587,0.065475) (0.75362,0.070578) (0.73064,0.075592) (0.70699,0.080481) (0.68274,0.08521) (0.65795,0.089745) (0.63268,0.094053) (0.60702,0.098103) (0.58102,0.10187) (0.55476,0.10531) (0.52832,0.10842) (0.50176,0.11117) (0.47517,0.11353) (0.44862,0.11549) (0.4222,0.11703) (0.39587,0.11814) (0.36929,0.11863) (0.34306,0.11839) (0.31726,0.11745) (0.29199,0.11582) (0.26733,0.11352) (0.24336,0.1106) (0.22016,0.10708) (0.1978,0.10301) (0.17637,0.098448) (0.15592,0.093438) (0.13653,0.08804) (0.11825,0.082313) (0.10114,0.076319) (0.085239,0.070121) (0.0706,0.063781) (0.057259,0.057362) (0.045249,0.050923) (0.034597,0.044522) (0.025325,0.038211) (0.017452,0.032041) (0.01099,0.026055) (0.0059485,0.020291) (0.0023306,0.014783) (0.00013649,0.0095539) (-0.00063772,0.0046228) (0,0) (0.0020082,-0.004212) (0.0053416,-0.0079161) (0.0099811,-0.011118) (0.015904,-0.013825) (0.023084,-0.01605) (0.031491,-0.017807) (0.041094,-0.019112) (0.051858,-0.019987) (0.063744,-0.020453) (0.076715,-0.020535) (0.090729,-0.020263) (0.10574,-0.019665) (0.12172,-0.018775) (0.13861,-0.017626) (0.15636,-0.016257) (0.17495,-0.014704) (0.19431,-0.013008) (0.21441,-0.011209) (0.2352,-0.0093499) (0.25664,-0.0074726) (0.27868,-0.0056194) (0.30127,-0.0038326) (0.32437,-0.0021532) (0.34792,-0.00062125) (0.37189,0.00072516) (0.39622,0.0018503) (0.42137,0.0028142) (0.44685,0.0037552) (0.47249,0.0046567) (0.49824,0.0055012) (0.52402,0.0062729) (0.54977,0.0069586) (0.57542,0.0075474) (0.60089,0.0080312) (0.62613,0.0084043) (0.65107,0.0086637) (0.67563,0.0088089) (0.69975,0.0088418) (0.72335,0.0087665) (0.74638,0.0085892) (0.76877,0.0083177) (0.79046,0.0079615) (0.81137,0.0075311) (0.83145,0.0070382) (0.85065,0.0064948) (0.8689,0.0059137) (0.88614,0.0053074) (0.90234,0.0046886) (0.91743,0.0040695) (0.93138,0.003462) (0.94413,0.002877) (0.95566,0.0023249) (0.96591,0.001815) (0.97487,0.0013559) (0.9825,0.00095476) (0.98877,0.00061796) (0.99367,0.00035063) (0.99718,0.0001568) (0.9993,3.9343e-05) (1,1.633e-17)};\addlegendentry{$m=6~\%$}
\addplot[corange,dashed,thin,forget plot] coordinates {(0,0) (0.016949,0.004977) (0.033898,0.0097386) (0.050847,0.014285) (0.067797,0.018615) (0.084746,0.022731) (0.10169,0.02663) (0.11864,0.030315) (0.13559,0.033783) (0.15254,0.037037) (0.16949,0.040075) (0.18644,0.042897) (0.20339,0.045504) (0.22034,0.047896) (0.23729,0.050072) (0.25424,0.052032) (0.27119,0.053778) (0.28814,0.055307) (0.30508,0.056622) (0.32203,0.05772) (0.33898,0.058604) (0.35593,0.059272) (0.37288,0.059724) (0.38983,0.059961) (0.40678,0.059992) (0.42373,0.059906) (0.44068,0.059724) (0.45763,0.059447) (0.47458,0.059073) (0.49153,0.058604) (0.50847,0.058039) (0.52542,0.057378) (0.54237,0.056622) (0.55932,0.055769) (0.57627,0.054821) (0.59322,0.053778) (0.61017,0.052638) (0.62712,0.051403) (0.64407,0.050072) (0.66102,0.048645) (0.67797,0.047122) (0.69492,0.045504) (0.71186,0.04379) (0.72881,0.04198) (0.74576,0.040075) (0.76271,0.038073) (0.77966,0.035976) (0.79661,0.033783) (0.81356,0.031495) (0.83051,0.02911) (0.84746,0.02663) (0.86441,0.024054) (0.88136,0.021383) (0.89831,0.018615) (0.91525,0.015752) (0.9322,0.012793) (0.94915,0.0097386) (0.9661,0.0065881) (0.98305,0.003342) (1,0)};
\nextgroupplot[title={effet de $p$ ($m=5$~\%, $t=12$~\%)}]
\addplot[cmuted,thick] coordinates {(1,-1.6525e-17) (0.99933,0.0001844) (0.99731,0.00073616) (0.99395,0.0016509) (0.98927,0.0029215) (0.98326,0.0045381) (0.97595,0.0064881) (0.96735,0.0087568) (0.95749,0.011327) (0.94639,0.01418) (0.93407,0.017294) (0.92058,0.020647) (0.90594,0.024217) (0.89019,0.027977) (0.87337,0.031903) (0.85553,0.035968) (0.83671,0.040147) (0.81696,0.044412) (0.79633,0.048736) (0.77488,0.053094) (0.75267,0.057457) (0.72975,0.0618) (0.70619,0.066096) (0.68205,0.07032) (0.65739,0.074445) (0.63228,0.078447) (0.6068,0.082302) (0.581,0.085984) (0.55497,0.089472) (0.52877,0.092742) (0.50248,0.095772) (0.47616,0.098543) (0.4499,0.10104) (0.42376,0.10323) (0.39782,0.10511) (0.37216,0.10667) (0.34685,0.10789) (0.32195,0.10876) (0.29754,0.10927) (0.27369,0.10942) (0.25046,0.10921) (0.22793,0.10864) (0.20616,0.10771) (0.18327,0.10614) (0.16071,0.1032) (0.1394,0.099028) (0.11945,0.093795) (0.10091,0.087691) (0.083851,0.080901) (0.068311,0.073606) (0.05432,0.065979) (0.041894,0.058182) (0.031044,0.050364) (0.021769,0.042663) (0.014067,0.035198) (0.0079261,0.028074) (0.0033323,0.021379) (0.00026701,0.015182) (-0.0012922,0.0095347) (-0.0013713,0.0044693) (0,0) (0.0027418,-0.0037852) (0.0067703,-0.0068144) (0.012045,-0.0091211) (0.01852,-0.010752) (0.026148,-0.011763) (0.034877,-0.012224) (0.04465,-0.01221) (0.055411,-0.011809) (0.067099,-0.01111) (0.079655,-0.01021) (0.093018,-0.0092081) (0.10713,-0.0082058) (0.12194,-0.0073036) (0.13741,-0.0066015) (0.15349,-0.0061982) (0.17016,-0.0061909) (0.18741,-0.0066763) (0.20605,-0.0077147) (0.22743,-0.0087581) (0.24954,-0.009601) (0.27232,-0.010254) (0.29573,-0.010728) (0.31969,-0.011037) (0.34414,-0.011195) (0.36902,-0.011217) (0.39426,-0.011119) (0.4198,-0.010917) (0.44557,-0.010625) (0.4715,-0.01026) (0.49752,-0.0098348) (0.52357,-0.0093643) (0.54956,-0.0088608) (0.57543,-0.0083355) (0.60111,-0.0077986) (0.62653,-0.0072586) (0.65163,-0.0067229) (0.67632,-0.0061973) (0.70055,-0.0056866) (0.72424,-0.0051942) (0.74733,-0.0047226) (0.76975,-0.0042733) (0.79145,-0.003847) (0.81236,-0.0034442) (0.83242,-0.0030645) (0.85158,-0.0027076) (0.86977,-0.002373) (0.88696,-0.0020601) (0.90308,-0.0017684) (0.91809,-0.0014978) (0.93195,-0.0012481) (0.94462,-0.0010194) (0.95606,-0.00081205) (0.96623,-0.00062663) (0.97511,-0.00046375) (0.98267,-0.00032417) (0.98888,-0.00020864) (0.99373,-0.00011789) (0.99721,-5.257e-05) (0.9993,-1.3169e-05) (1,1.6525e-17)};\addlegendentry{$p=0{,}20$}
\addplot[cmuted,dashed,thin,forget plot] coordinates {(0,0) (0.016949,0.0081155) (0.033898,0.015513) (0.050847,0.022192) (0.067797,0.028153) (0.084746,0.033396) (0.10169,0.03792) (0.11864,0.041727) (0.13559,0.044815) (0.15254,0.047185) (0.16949,0.048837) (0.18644,0.04977) (0.20339,0.049999) (0.22034,0.049968) (0.23729,0.049891) (0.25424,0.04977) (0.27119,0.049604) (0.28814,0.049393) (0.30508,0.049137) (0.32203,0.048837) (0.33898,0.048491) (0.35593,0.0481) (0.37288,0.047665) (0.38983,0.047185) (0.40678,0.04666) (0.42373,0.046089) (0.44068,0.045475) (0.45763,0.044815) (0.47458,0.04411) (0.49153,0.04336) (0.50847,0.042566) (0.52542,0.041727) (0.54237,0.040842) (0.55932,0.039913) (0.57627,0.038939) (0.59322,0.03792) (0.61017,0.036856) (0.62712,0.035748) (0.64407,0.034594) (0.66102,0.033396) (0.67797,0.032152) (0.69492,0.030864) (0.71186,0.029531) (0.72881,0.028153) (0.74576,0.02673) (0.76271,0.025262) (0.77966,0.023749) (0.79661,0.022192) (0.81356,0.020589) (0.83051,0.018942) (0.84746,0.01725) (0.86441,0.015513) (0.88136,0.013731) (0.89831,0.011904) (0.91525,0.010032) (0.9322,0.0081155) (0.94915,0.0061539) (0.9661,0.0041475) (0.98305,0.0020962) (1,0)};
\addplot[cblue,very thick] coordinates {(1,-1.6427e-17) (0.99933,0.00021234) (0.99733,0.00084752) (0.99399,0.0019) (0.98933,0.0033607) (0.98335,0.005217) (0.97608,0.0074529) (0.96753,0.01005) (0.95771,0.012985) (0.94665,0.016235) (0.93438,0.019772) (0.92093,0.023568) (0.90632,0.027593) (0.8906,0.031815) (0.87381,0.0362) (0.85598,0.040715) (0.83717,0.045326) (0.81741,0.049999) (0.79676,0.054699) (0.77528,0.059392) (0.75302,0.064044) (0.73004,0.068622) (0.70639,0.073093) (0.68215,0.077425) (0.65738,0.081588) (0.63214,0.085553) (0.60651,0.08929) (0.58055,0.092773) (0.55435,0.095977) (0.52796,0.098877) (0.50147,0.10145) (0.47495,0.10368) (0.44848,0.10555) (0.42213,0.10704) (0.3959,0.10814) (0.36951,0.10868) (0.34346,0.10859) (0.31785,0.10787) (0.29276,0.10654) (0.26827,0.10461) (0.24445,0.10212) (0.22139,0.099082) (0.19916,0.095544) (0.17783,0.091545) (0.15747,0.087127) (0.13814,0.082338) (0.11989,0.077229) (0.10279,0.071852) (0.086882,0.06626) (0.072207,0.060508) (0.058807,0.054647) (0.046715,0.048732) (0.035961,0.042811) (0.026568,0.036932) (0.018555,0.031139) (0.011938,0.02547) (0.0067275,0.019961) (0.0029283,0.014639) (0.00054272,0.0095283) (-0.00043136,0.0046452) (0,0) (0.0018018,-0.0043029) (0.0049354,-0.0081635) (0.0093834,-0.011585) (0.015125,-0.014572) (0.022136,-0.017133) (0.030388,-0.019278) (0.039852,-0.021017) (0.050494,-0.022366) (0.062278,-0.02334) (0.075168,-0.023958) (0.089122,-0.024242) (0.1041,-0.024214) (0.12006,-0.023898) (0.13696,-0.023324) (0.15475,-0.022519) (0.1734,-0.021515) (0.19285,-0.020344) (0.21305,-0.019041) (0.23397,-0.017641) (0.25555,-0.016179) (0.27774,-0.014693) (0.3005,-0.013219) (0.32378,-0.011793) (0.34752,-0.010451) (0.37167,-0.0092252) (0.39619,-0.0081488) (0.42144,-0.0071731) (0.44699,-0.0061837) (0.47272,-0.0051969) (0.49853,-0.00423) (0.52438,-0.0032987) (0.55018,-0.0024167) (0.57588,-0.0015957) (0.6014,-0.00084498) (0.62668,-0.00017188) (0.65164,0.00041864) (0.67622,0.00092376) (0.70035,0.0013427) (0.72395,0.0016766) (0.74698,0.0019282) (0.76936,0.0021018) (0.79102,0.0022032) (0.81191,0.0022388) (0.83196,0.0022161) (0.85112,0.0021431) (0.86933,0.0020278) (0.88654,0.0018788) (0.90269,0.0017043) (0.91774,0.0015123) (0.93165,0.0013104) (0.94436,0.0011058) (0.95584,0.00090506) (0.96606,0.00071404) (0.97498,0.00053799) (0.98257,0.00038144) (0.98882,0.00024822) (0.9937,0.00014141) (0.9972,6.3412e-05) (0.9993,1.5937e-05) (1,1.6427e-17)};\addlegendentry{$p=0{,}40$}
\addplot[cblue,dashed,thin,forget plot] coordinates {(0,0) (0.016949,0.0041475) (0.033898,0.0081155) (0.050847,0.011904) (0.067797,0.015513) (0.084746,0.018942) (0.10169,0.022192) (0.11864,0.025262) (0.13559,0.028153) (0.15254,0.030864) (0.16949,0.033396) (0.18644,0.035748) (0.20339,0.03792) (0.22034,0.039913) (0.23729,0.041727) (0.25424,0.04336) (0.27119,0.044815) (0.28814,0.046089) (0.30508,0.047185) (0.32203,0.0481) (0.33898,0.048837) (0.35593,0.049393) (0.37288,0.04977) (0.38983,0.049968) (0.40678,0.049994) (0.42373,0.049922) (0.44068,0.04977) (0.45763,0.049539) (0.47458,0.049228) (0.49153,0.048837) (0.50847,0.048366) (0.52542,0.047815) (0.54237,0.047185) (0.55932,0.046475) (0.57627,0.045685) (0.59322,0.044815) (0.61017,0.043865) (0.62712,0.042836) (0.64407,0.041727) (0.66102,0.040538) (0.67797,0.039269) (0.69492,0.03792) (0.71186,0.036492) (0.72881,0.034984) (0.74576,0.033396) (0.76271,0.031728) (0.77966,0.02998) (0.79661,0.028153) (0.81356,0.026246) (0.83051,0.024259) (0.84746,0.022192) (0.86441,0.020045) (0.88136,0.017819) (0.89831,0.015513) (0.91525,0.013127) (0.9322,0.010661) (0.94915,0.0081155) (0.9661,0.0054901) (0.98305,0.002785) (1,0)};
\addplot[corange,very thick] coordinates {(1,-1.6257e-17) (0.99934,0.00024935) (0.99735,0.00099489) (0.99403,0.0022291) (0.98941,0.0039397) (0.98348,0.0061095) (0.97625,0.0087168) (0.96774,0.011736) (0.95798,0.015137) (0.94697,0.018886) (0.93475,0.022947) (0.92133,0.02728) (0.90676,0.031844) (0.89106,0.036594) (0.87428,0.041486) (0.85644,0.046472) (0.83759,0.051506) (0.81778,0.05654) (0.79705,0.061528) (0.77546,0.066422) (0.75307,0.071177) (0.72992,0.075748) (0.70609,0.080092) (0.68163,0.084167) (0.65661,0.087935) (0.63109,0.091359) (0.60516,0.094405) (0.57888,0.097041) (0.55232,0.099239) (0.52577,0.10102) (0.49913,0.10244) (0.47246,0.10348) (0.44585,0.10411) (0.41936,0.10434) (0.39309,0.10415) (0.3671,0.10354) (0.34147,0.10251) (0.31629,0.10105) (0.29162,0.099185) (0.26755,0.09691) (0.24414,0.094241) (0.22146,0.091197) (0.19959,0.087795) (0.17858,0.08406) (0.15851,0.080018) (0.13943,0.075698) (0.12139,0.071132) (0.10445,0.066351) (0.088649,0.061393) (0.07404,0.05629) (0.060657,0.051081) (0.048537,0.045799) (0.037711,0.040481) (0.028204,0.035158) (0.02004,0.029862) (0.013237,0.024621) (0.0078086,0.019461) (0.0037665,0.014403) (0.0011165,0.0094657) (-0.00013866,0.0046616) (0,0) (0.0015091,-0.0044125) (0.0043616,-0.0084721) (0.0085452,-0.012177) (0.014044,-0.015527) (0.020838,-0.018522) (0.028904,-0.021161) (0.038216,-0.023446) (0.048744,-0.025379) (0.060456,-0.026964) (0.073317,-0.028205) (0.08729,-0.029109) (0.10233,-0.029683) (0.11841,-0.029937) (0.13547,-0.029883) (0.15347,-0.029535) (0.17236,-0.028908) (0.1921,-0.02802) (0.21263,-0.02689) (0.2339,-0.02554) (0.25586,-0.023994) (0.27846,-0.022274) (0.30164,-0.020407) (0.32534,-0.018418) (0.34951,-0.016334) (0.37408,-0.014182) (0.399,-0.011987) (0.4242,-0.0097743) (0.44962,-0.0075688) (0.4752,-0.0053929) (0.50087,-0.0032675) (0.52657,-0.0012115) (0.55221,0.00075811) (0.57756,0.0025658) (0.60275,0.0041574) (0.62773,0.0055267) (0.65241,0.0066711) (0.67674,0.0075916) (0.70065,0.0082928) (0.72407,0.0087821) (0.74693,0.0090702) (0.76918,0.0091701) (0.79074,0.0090975) (0.81154,0.0088696) (0.83154,0.0085052) (0.85067,0.0080245) (0.86887,0.0074483) (0.88608,0.0067976) (0.90225,0.0060937) (0.91734,0.0053575) (0.93128,0.0046091) (0.94404,0.0038681) (0.95557,0.0031526) (0.96584,0.0024794) (0.97481,0.0018638) (0.98245,0.0013192) (0.98874,0.00085744) (0.99365,0.00048808) (0.99718,0.00021876) (0.99929,5.4967e-05) (1,1.6257e-17)};\addlegendentry{$p=0{,}55$}
\addplot[corange,dashed,thin,forget plot] coordinates {(0,0) (0.016949,0.0030342) (0.033898,0.0059734) (0.050847,0.0088176) (0.067797,0.011567) (0.084746,0.014221) (0.10169,0.016781) (0.11864,0.019245) (0.13559,0.021614) (0.15254,0.023889) (0.16949,0.026068) (0.18644,0.028153) (0.20339,0.030142) (0.22034,0.032037) (0.23729,0.033837) (0.25424,0.035541) (0.27119,0.037151) (0.28814,0.038666) (0.30508,0.040085) (0.32203,0.04141) (0.33898,0.04264) (0.35593,0.043775) (0.37288,0.044815) (0.38983,0.04576) (0.40678,0.04661) (0.42373,0.047365) (0.44068,0.048025) (0.45763,0.04859) (0.47458,0.04906) (0.49153,0.049435) (0.50847,0.049715) (0.52542,0.0499) (0.54237,0.04999) (0.55932,0.049979) (0.57627,0.04983) (0.59322,0.049539) (0.61017,0.049106) (0.62712,0.048532) (0.64407,0.047815) (0.66102,0.046957) (0.67797,0.045957) (0.69492,0.044815) (0.71186,0.043531) (0.72881,0.042105) (0.74576,0.040538) (0.76271,0.038828) (0.77966,0.036977) (0.79661,0.034984) (0.81356,0.032849) (0.83051,0.030572) (0.84746,0.028153) (0.86441,0.025592) (0.88136,0.02289) (0.89831,0.020045) (0.91525,0.017059) (0.9322,0.013931) (0.94915,0.010661) (0.9661,0.0072492) (0.98305,0.0036955) (1,0)};
\nextgroupplot[title={effet de $t$ ($m=2$~\%, $p=0{,}4$)}]
\addplot[cmuted,thick] coordinates {(1,-8.3082e-18) (0.99932,9.5323e-05) (0.99727,0.00038046) (0.99387,0.00085291) (0.98913,0.0015086) (0.98304,0.0023417) (0.97564,0.0033453) (0.96694,0.0045107) (0.95696,0.0058281) (0.94573,0.0072865) (0.93329,0.0088741) (0.91966,0.010578) (0.90487,0.012385) (0.88898,0.01428) (0.87202,0.01625) (0.85404,0.018279) (0.83509,0.020352) (0.81521,0.022454) (0.79447,0.02457) (0.77291,0.026685) (0.75061,0.028785) (0.72761,0.030854) (0.70397,0.032878) (0.67978,0.034843) (0.65508,0.036736) (0.62996,0.038544) (0.60447,0.040253) (0.57868,0.041852) (0.55268,0.043329) (0.52653,0.044673) (0.50029,0.045874) (0.47406,0.046922) (0.44788,0.047809) (0.42185,0.048528) (0.39602,0.04907) (0.37037,0.049374) (0.34509,0.049404) (0.32022,0.049164) (0.29586,0.048657) (0.27206,0.04789) (0.24889,0.04687) (0.22642,0.04561) (0.20471,0.044123) (0.18383,0.042422) (0.16383,0.040526) (0.14477,0.038452) (0.1267,0.03622) (0.10968,0.033851) (0.093744,0.031367) (0.078945,0.02879) (0.065322,0.026141) (0.05291,0.023443) (0.041744,0.020718) (0.031853,0.017984) (0.023262,0.015262) (0.015993,0.012569) (0.010066,0.0099208) (0.0054943,0.00733) (0.0022887,0.0048079) (0.0004562,0.0023627) (0,0) (0.00091426,-0.0022258) (0.0031894,-0.004262) (0.0068174,-0.0061083) (0.011786,-0.0077654) (0.018081,-0.0092345) (0.025682,-0.010518) (0.034567,-0.011618) (0.04471,-0.012539) (0.056083,-0.013286) (0.068653,-0.013866) (0.082384,-0.014284) (0.097239,-0.014549) (0.11318,-0.01467) (0.13015,-0.014658) (0.14813,-0.014524) (0.16704,-0.014281) (0.18685,-0.013942) (0.2075,-0.013521) (0.22894,-0.013034) (0.25111,-0.012495) (0.27395,-0.011921) (0.29741,-0.011328) (0.32141,-0.010732) (0.3459,-0.010147) (0.37081,-0.0095897) (0.39607,-0.009074) (0.42171,-0.0085804) (0.44759,-0.0080623) (0.47361,-0.0075275) (0.49971,-0.0069847) (0.52581,-0.0064414) (0.55185,-0.0059048) (0.57775,-0.005381) (0.60344,-0.0048751) (0.62886,-0.0043915) (0.65393,-0.0039336) (0.67859,-0.0035038) (0.70276,-0.0031038) (0.72639,-0.0027345) (0.74939,-0.002396) (0.77172,-0.0020879) (0.79332,-0.0018092) (0.81411,-0.0015587) (0.83404,-0.0013347) (0.85306,-0.0011353) (0.87112,-0.00095864) (0.88816,-0.00080276) (0.90414,-0.00066572) (0.91901,-0.00054571) (0.93274,-0.00044107) (0.94527,-0.00035028) (0.95658,-0.00027204) (0.96664,-0.00020524) (0.97542,-0.00014893) (0.98288,-0.00010238) (0.98902,-6.5003e-05) (0.99381,-3.6343e-05) (0.99725,-1.6084e-05) (0.99931,-4.011e-06) (1,8.3082e-18)};\addlegendentry{$t=6~\%$}
\addplot[cblue,very thick] coordinates {(1,-1.6616e-17) (0.99932,0.00014499) (0.99729,0.00057873) (0.9939,0.0012975) (0.98918,0.0022953) (0.98312,0.0035638) (0.97575,0.0050924) (0.96709,0.0068686) (0.95715,0.0088781) (0.94597,0.011105) (0.93356,0.013532) (0.91998,0.01614) (0.90524,0.01891) (0.88939,0.021822) (0.87247,0.024854) (0.85453,0.027985) (0.83561,0.031195) (0.81577,0.03446) (0.79505,0.03776) (0.77351,0.041072) (0.75121,0.044375) (0.72822,0.047648) (0.70458,0.050869) (0.68037,0.054017) (0.65566,0.057071) (0.6305,0.060012) (0.60498,0.062817) (0.57915,0.065468) (0.5531,0.067946) (0.52688,0.07023) (0.50059,0.072303) (0.47428,0.074147) (0.44803,0.075745) (0.42192,0.077082) (0.39599,0.078142) (0.37016,0.078855) (0.34468,0.07918) (0.31963,0.079112) (0.29508,0.07865) (0.27111,0.077795) (0.24777,0.076553) (0.22515,0.074933) (0.20331,0.072945) (0.18232,0.070605) (0.16222,0.06793) (0.14309,0.06494) (0.12497,0.061659) (0.10793,0.058112) (0.091996,0.054325) (0.077226,0.050327) (0.063657,0.046145) (0.051324,0.041808) (0.040261,0.037346) (0.030495,0.032785) (0.022051,0.028152) (0.01495,0.023471) (0.0092059,0.018764) (0.0048327,0.014049) (0.0018383,0.0093428) (0.00022717,0.0046569) (0,0) (0.0011433,-0.00452) (0.0036398,-0.0087969) (0.007479,-0.012828) (0.012647,-0.016609) (0.019125,-0.020136) (0.026892,-0.023408) (0.035924,-0.026419) (0.046194,-0.029168) (0.057669,-0.031651) (0.070318,-0.033869) (0.084103,-0.03582) (0.098987,-0.037507) (0.11493,-0.038931) (0.13188,-0.040097) (0.1498,-0.041012) (0.16865,-0.041685) (0.18836,-0.042124) (0.2089,-0.042343) (0.23021,-0.042356) (0.25223,-0.042178) (0.2749,-0.041827) (0.29818,-0.041321) (0.322,-0.040679) (0.3463,-0.039923) (0.37102,-0.039071) (0.3961,-0.038146) (0.42164,-0.037134) (0.44744,-0.035998) (0.47339,-0.034752) (0.49941,-0.033414) (0.52545,-0.031999) (0.55143,-0.030522) (0.57728,-0.028997) (0.60293,-0.027439) (0.62832,-0.025859) (0.65336,-0.024269) (0.67799,-0.022677) (0.70216,-0.021095) (0.72577,-0.019529) (0.74879,-0.017986) (0.77113,-0.016475) (0.79274,-0.014999) (0.81355,-0.013565) (0.83352,-0.012178) (0.85258,-0.010842) (0.87067,-0.0095628) (0.88775,-0.0083442) (0.90378,-0.0071909) (0.91869,-0.0061076) (0.93246,-0.0050987) (0.94504,-0.0041687) (0.9564,-0.0033221) (0.96649,-0.0025632) (0.9753,-0.001896) (0.9828,-0.0013244) (0.98897,-0.00085179) (0.99379,-0.00048097) (0.99723,-0.00021436) (0.99931,-5.3678e-05) (1,1.6616e-17)};\addlegendentry{$t=12~\%$}
\addplot[corange,very thick] coordinates {(1,-2.7694e-17) (0.99933,0.00021121) (0.9973,0.00084309) (0.99394,0.0018904) (0.98925,0.0033444) (0.98323,0.0051932) (0.9759,0.0074219) (0.96729,0.010013) (0.9574,0.012945) (0.94627,0.016196) (0.93393,0.019742) (0.92041,0.023556) (0.90573,0.02761) (0.88994,0.031877) (0.87308,0.036326) (0.85518,0.040928) (0.83631,0.045652) (0.8165,0.050469) (0.79582,0.055346) (0.7743,0.060254) (0.75202,0.065162) (0.72903,0.07004) (0.70539,0.074857) (0.68117,0.079582) (0.65643,0.084185) (0.63123,0.088635) (0.60566,0.092903) (0.57978,0.096957) (0.55365,0.10077) (0.52736,0.10431) (0.50098,0.10754) (0.47458,0.11045) (0.44823,0.11299) (0.42201,0.11515) (0.39595,0.1169) (0.36987,0.11816) (0.34414,0.11888) (0.31884,0.11904) (0.29405,0.11864) (0.26984,0.11767) (0.24629,0.11613) (0.22347,0.11403) (0.20145,0.11137) (0.1803,0.10818) (0.16008,0.10447) (0.14085,0.10026) (0.12267,0.095578) (0.10559,0.09046) (0.089666,0.084936) (0.074934,0.079042) (0.061436,0.072816) (0.049209,0.066295) (0.038283,0.059517) (0.028686,0.05252) (0.020438,0.045339) (0.013558,0.038007) (0.008059,0.030555) (0.0039506,0.023008) (0.0012378,0.015389) (-7.8199e-05,0.0077159) (0,0) (0.0014487,-0.007579) (0.0042403,-0.014843) (0.0083611,-0.021786) (0.013793,-0.028399) (0.020516,-0.034672) (0.028506,-0.040594) (0.037734,-0.046154) (0.048171,-0.051339) (0.059784,-0.056138) (0.072538,-0.06054) (0.086396,-0.064536) (0.10132,-0.068117) (0.11726,-0.071278) (0.13418,-0.074016) (0.15204,-0.07633) (0.17079,-0.078223) (0.19038,-0.079701) (0.21076,-0.080773) (0.23189,-0.081452) (0.25371,-0.081755) (0.27617,-0.081701) (0.29922,-0.081311) (0.32279,-0.08061) (0.34684,-0.079623) (0.37131,-0.07838) (0.39614,-0.076909) (0.42155,-0.075206) (0.44724,-0.073245) (0.47309,-0.071052) (0.49902,-0.068653) (0.52497,-0.066075) (0.55088,-0.063344) (0.57666,-0.060486) (0.60225,-0.057525) (0.62759,-0.054483) (0.65259,-0.051382) (0.6772,-0.048242) (0.70135,-0.045082) (0.72496,-0.041921) (0.74798,-0.038774) (0.77034,-0.035657) (0.79197,-0.032585) (0.81282,-0.029573) (0.83282,-0.026635) (0.85192,-0.023785) (0.87007,-0.021035) (0.88721,-0.018399) (0.90329,-0.015891) (0.91827,-0.013523) (0.93209,-0.011309) (0.94473,-0.0092599) (0.95614,-0.0073888) (0.9663,-0.0057071) (0.97516,-0.0042255) (0.9827,-0.0029539) (0.9889,-0.0019008) (0.99375,-0.0010738) (0.99722,-0.00047872) (0.9993,-0.0001199) (1,2.7694e-17)};\addlegendentry{$t=20~\%$}
\end{groupplot}
\end{tikzpicture}

\caption{Effet de chaque paramètre (calculé par \code{naca4}). Tirets : ligne moyenne.}\label{fig:parametres}
\end{figure}
\begin{itemize}
\item \textbf{L'épaisseur $t$} pilote le rayon de bord d'attaque ($r_{BA}\propto t^2$), donc la tolérance en incidence et la forme du pic d'aspiration. Elle fixe aussi la survitesse moyenne sur le profil, donc le gradient de pression adverse en aval et la traînée de forme. Elle augmente légèrement la pente de portance, $\dd C_L/\dd\alpha\approx2\pi(1+0{,}77\,t)$.
\item \textbf{La cambrure $m$} décale la courbe de portance. En théorie des profils minces, $C_L=2\pi(\alpha-\alpha_{L0})$ avec un angle de portance nulle $\alpha_{L0}<0$ proportionnel à $m$. Elle crée un moment piqueur, $C_{M,c/4}<0$.
\item \textbf{La position $p$} répartit la charge le long de la corde. Une cambrure arrière ($p$ grand) charge le bord de fuite, augmente $|C_{M}|$ et crée un gradient adverse plus fort en fin d'extrados. Une cambrure avant renforce le pic d'aspiration au bord d'attaque.
\end{itemize}

\begin{figure}[htbp]
\centering
% figure générée par tools/make_figures.py : profils
\begin{tikzpicture}
\begin{axis}[width=\linewidth,height=4.6cm,axis equal image,xmin=-0.02,xmax=1.02,ymin=-0.12,ymax=0.16,
  xlabel={$x/c$},ylabel={$y/c$},legend columns=3,legend style={draw=none,at={(0.5,1.03)},anchor=south,font=\small}]
\addplot[cmuted,thick] coordinates {(1,-1.6653e-17) (0.99961,5.6016e-05) (0.99846,0.00022381) (0.99653,0.00050262) (0.99384,0.00089119) (0.99039,0.0013878) (0.98618,0.0019902) (0.98123,0.0026957) (0.97553,0.0035014) (0.9691,0.0044037) (0.96194,0.0053988) (0.95407,0.0064826) (0.9455,0.0076508) (0.93625,0.0088987) (0.92632,0.010221) (0.91573,0.011614) (0.90451,0.013071) (0.89266,0.014587) (0.8802,0.016158) (0.86716,0.017777) (0.85355,0.019438) (0.8394,0.021138) (0.82472,0.022869) (0.80955,0.024626) (0.79389,0.026405) (0.77779,0.028198) (0.76125,0.03) (0.74431,0.031806) (0.727,0.03361) (0.70933,0.035406) (0.69134,0.037188) (0.67306,0.03895) (0.65451,0.040686) (0.63572,0.04239) (0.61672,0.044055) (0.59755,0.045674) (0.57822,0.047242) (0.55877,0.048752) (0.53923,0.050196) (0.51963,0.051568) (0.5,0.052862) (0.48037,0.054069) (0.46077,0.055184) (0.44123,0.056199) (0.42178,0.057108) (0.40245,0.057905) (0.38328,0.058582) (0.36428,0.059135) (0.34549,0.059557) (0.32694,0.059843) (0.30866,0.059988) (0.29067,0.059989) (0.273,0.059841) (0.25569,0.059542) (0.23875,0.059088) (0.22221,0.058478) (0.20611,0.057712) (0.19045,0.056788) (0.17528,0.055708) (0.1606,0.054472) (0.14645,0.053083) (0.13284,0.051542) (0.1198,0.049854) (0.10734,0.048021) (0.095492,0.046049) (0.084265,0.043942) (0.07368,0.041705) (0.063752,0.039345) (0.054497,0.036867) (0.045928,0.034276) (0.03806,0.03158) (0.030904,0.028783) (0.024472,0.025893) (0.018772,0.022915) (0.013815,0.019854) (0.0096074,0.016715) (0.0061558,0.013503) (0.0034658,0.010223) (0.0015413,0.0068767) (0.00038548,0.0034684) (0,0) (0.00038548,-0.0034684) (0.0015413,-0.0068767) (0.0034658,-0.010223) (0.0061558,-0.013503) (0.0096074,-0.016715) (0.013815,-0.019854) (0.018772,-0.022915) (0.024472,-0.025893) (0.030904,-0.028783) (0.03806,-0.03158) (0.045928,-0.034276) (0.054497,-0.036867) (0.063752,-0.039345) (0.07368,-0.041705) (0.084265,-0.043942) (0.095492,-0.046049) (0.10734,-0.048021) (0.1198,-0.049854) (0.13284,-0.051542) (0.14645,-0.053083) (0.1606,-0.054472) (0.17528,-0.055708) (0.19045,-0.056788) (0.20611,-0.057712) (0.22221,-0.058478) (0.23875,-0.059088) (0.25569,-0.059542) (0.273,-0.059841) (0.29067,-0.059989) (0.30866,-0.059988) (0.32694,-0.059843) (0.34549,-0.059557) (0.36428,-0.059135) (0.38328,-0.058582) (0.40245,-0.057905) (0.42178,-0.057108) (0.44123,-0.056199) (0.46077,-0.055184) (0.48037,-0.054069) (0.5,-0.052862) (0.51963,-0.051568) (0.53923,-0.050196) (0.55877,-0.048752) (0.57822,-0.047242) (0.59755,-0.045674) (0.61672,-0.044055) (0.63572,-0.04239) (0.65451,-0.040686) (0.67306,-0.03895) (0.69134,-0.037188) (0.70933,-0.035406) (0.727,-0.03361) (0.74431,-0.031806) (0.76125,-0.03) (0.77779,-0.028198) (0.79389,-0.026405) (0.80955,-0.024626) (0.82472,-0.022869) (0.8394,-0.021138) (0.85355,-0.019438) (0.86716,-0.017777) (0.8802,-0.016158) (0.89266,-0.014587) (0.90451,-0.013071) (0.91573,-0.011614) (0.92632,-0.010221) (0.93625,-0.0088987) (0.9455,-0.0076508) (0.95407,-0.0064826) (0.96194,-0.0053988) (0.9691,-0.0044037) (0.97553,-0.0035014) (0.98123,-0.0026957) (0.98618,-0.0019902) (0.99039,-0.0013878) (0.99384,-0.00089119) (0.99653,-0.00050262) (0.99846,-0.00022381) (0.99961,-5.6016e-05) (1,1.6653e-17)};
\addlegendentry{NACA 0012}
\addplot[cblue,very thick] coordinates {(1,-1.6616e-17) (0.99962,8.1582e-05) (0.99847,0.00032594) (0.99657,0.0007319) (0.9939,0.0012975) (0.99048,0.0020202) (0.98631,0.0028964) (0.9814,0.003922) (0.97575,0.0050924) (0.96937,0.0064021) (0.96228,0.0078452) (0.95447,0.0094151) (0.94597,0.011105) (0.93678,0.012907) (0.92692,0.014814) (0.9164,0.016818) (0.90524,0.01891) (0.89346,0.021082) (0.88106,0.023324) (0.86808,0.025628) (0.85453,0.027985) (0.84043,0.030386) (0.8258,0.032822) (0.81067,0.035283) (0.79505,0.03776) (0.77897,0.040244) (0.76245,0.042726) (0.74553,0.045197) (0.72822,0.047648) (0.71055,0.05007) (0.69254,0.052453) (0.67424,0.05479) (0.65566,0.057071) (0.63683,0.059288) (0.61778,0.061432) (0.59855,0.063495) (0.57915,0.065468) (0.55963,0.067344) (0.54001,0.069113) (0.52032,0.070769) (0.50059,0.072303) (0.48085,0.073708) (0.46114,0.074977) (0.44149,0.076104) (0.42192,0.077082) (0.40247,0.077904) (0.38303,0.078547) (0.36375,0.078973) (0.34468,0.07918) (0.32585,0.079166) (0.30729,0.07893) (0.28903,0.078473) (0.27111,0.077795) (0.25354,0.0769) (0.23637,0.07579) (0.21962,0.07447) (0.20331,0.072945) (0.18748,0.071222) (0.17215,0.069308) (0.15735,0.067211) (0.14309,0.06494) (0.1294,0.062506) (0.11631,0.059917) (0.10384,0.057187) (0.091996,0.054325) (0.080807,0.051345) (0.070289,0.048257) (0.060456,0.045074) (0.051324,0.041808) (0.042906,0.038472) (0.035214,0.035076) (0.02826,0.031633) (0.022051,0.028152) (0.016598,0.024645) (0.011907,0.02112) (0.0079837,0.017585) (0.0048327,0.014049) (0.0024573,0.010518) (0.00085968,0.0069967) (4.0695e-05,0.0034897) (0,0) (0.00073027,-0.0034127) (0.002223,-0.006689) (0.0044742,-0.0098278) (0.007479,-0.012828) (0.011231,-0.015687) (0.015723,-0.018404) (0.020946,-0.020978) (0.026892,-0.023408) (0.033549,-0.025691) (0.040906,-0.027826) (0.048951,-0.029813) (0.057669,-0.031651) (0.067048,-0.03334) (0.077071,-0.034878) (0.087723,-0.036267) (0.098987,-0.037507) (0.11085,-0.038599) (0.12328,-0.039546) (0.13627,-0.040349) (0.1498,-0.041012) (0.16385,-0.041539) (0.1784,-0.041933) (0.19342,-0.042199) (0.2089,-0.042343) (0.22481,-0.042372) (0.24113,-0.04229) (0.25784,-0.042106) (0.2749,-0.041827) (0.29231,-0.041461) (0.31003,-0.041016) (0.32803,-0.0405) (0.3463,-0.039923) (0.36481,-0.039292) (0.38352,-0.038617) (0.40244,-0.037905) (0.42164,-0.037134) (0.44097,-0.036293) (0.4604,-0.035388) (0.47989,-0.034426) (0.49941,-0.033414) (0.51894,-0.032359) (0.53845,-0.031267) (0.55791,-0.030144) (0.57728,-0.028997) (0.59654,-0.027831) (0.61566,-0.026651) (0.63461,-0.025462) (0.65336,-0.024269) (0.67188,-0.023075) (0.69014,-0.021884) (0.70811,-0.020701) (0.72577,-0.019529) (0.74309,-0.018369) (0.76005,-0.017226) (0.7766,-0.016102) (0.79274,-0.014999) (0.80843,-0.013919) (0.82365,-0.012865) (0.83837,-0.011839) (0.85258,-0.010842) (0.86624,-0.0098771) (0.87934,-0.0089456) (0.89186,-0.0080496) (0.90378,-0.0071909) (0.91507,-0.0063716) (0.92572,-0.0055935) (0.93572,-0.0048586) (0.94504,-0.0041687) (0.95367,-0.0035257) (0.9616,-0.0029314) (0.96882,-0.0023876) (0.9753,-0.001896) (0.98105,-0.0014582) (0.98606,-0.0010756) (0.9903,-0.00074943) (0.99379,-0.00048097) (0.9965,-0.00027113) (0.99844,-0.00012069) (0.99961,-3.0201e-05) (1,1.6616e-17)};
\addlegendentry{NACA 2412}
\addplot[corange,very thick] coordinates {(1,-2.0634e-17) (0.99962,0.00012079) (0.9985,0.00048257) (0.99662,0.0010836) (0.99399,0.001921) (0.99062,0.0029907) (0.98651,0.0042877) (0.98166,0.0058057) (0.97608,0.0075377) (0.96979,0.0094755) (0.96278,0.01161) (0.95506,0.013932) (0.94665,0.01643) (0.93756,0.019094) (0.9278,0.021911) (0.91739,0.024869) (0.90633,0.027956) (0.89464,0.031158) (0.88235,0.034462) (0.86946,0.037853) (0.85599,0.041319) (0.84197,0.044845) (0.82741,0.048416) (0.81234,0.05202) (0.79677,0.055641) (0.78073,0.059266) (0.76425,0.06288) (0.74734,0.06647) (0.73004,0.070022) (0.71236,0.073522) (0.69435,0.076957) (0.67601,0.080314) (0.65738,0.083579) (0.63849,0.086741) (0.61937,0.089786) (0.60005,0.092702) (0.58055,0.095477) (0.56092,0.098101) (0.54117,0.10056) (0.52134,0.10285) (0.50147,0.10495) (0.48158,0.10686) (0.4617,0.10856) (0.44187,0.11006) (0.42213,0.11133) (0.40249,0.11238) (0.38267,0.11316) (0.36296,0.11359) (0.34346,0.11368) (0.32421,0.11342) (0.30524,0.11282) (0.28658,0.11189) (0.26826,0.11062) (0.25033,0.10903) (0.23282,0.10712) (0.21574,0.10491) (0.19915,0.1024) (0.18306,0.099621) (0.1675,0.096574) (0.15251,0.093279) (0.1381,0.089754) (0.12431,0.086017) (0.11115,0.082086) (0.09865,0.077982) (0.086827,0.073724) (0.0757,0.069333) (0.065285,0.06483) (0.055598,0.060235) (0.046652,0.055567) (0.038459,0.050847) (0.031031,0.046094) (0.024375,0.041324) (0.018499,0.036555) (0.013409,0.031803) (0.0091099,0.027083) (0.0056045,0.022405) (0.0028946,0.017783) (0.00098061,0.013225) (-0.00013821,0.0087379) (-0.00046398,0.0043285) (0,0) (0.0012349,-0.0041744) (0.0032209,-0.0081225) (0.0059509,-0.011844) (0.0094171,-0.015339) (0.01361,-0.018609) (0.01852,-0.021652) (0.024136,-0.024471) (0.030445,-0.027066) (0.037434,-0.02944) (0.04509,-0.031594) (0.053397,-0.033531) (0.062341,-0.035254) (0.071906,-0.036766) (0.082075,-0.038073) (0.09283,-0.039178) (0.10416,-0.040087) (0.11603,-0.040806) (0.12844,-0.041343) (0.14137,-0.041704) (0.15479,-0.041899) (0.16869,-0.041936) (0.18305,-0.041825) (0.19785,-0.041576) (0.21307,-0.041202) (0.22869,-0.040713) (0.24469,-0.040121) (0.26105,-0.03944) (0.27774,-0.038683) (0.29476,-0.037863) (0.31208,-0.036993) (0.32967,-0.036088) (0.34752,-0.035161) (0.3656,-0.034226) (0.38389,-0.033295) (0.40242,-0.032382) (0.42144,-0.031437) (0.44059,-0.030435) (0.45984,-0.029384) (0.47916,-0.028293) (0.49853,-0.027172) (0.51792,-0.026028) (0.53729,-0.024869) (0.55662,-0.023702) (0.57588,-0.022535) (0.59504,-0.021374) (0.61407,-0.020223) (0.63295,-0.019088) (0.65164,-0.017974) (0.67011,-0.016883) (0.68834,-0.015819) (0.70629,-0.014785) (0.72395,-0.013783) (0.74128,-0.012814) (0.75825,-0.01188) (0.77484,-0.010982) (0.79101,-0.010119) (0.80676,-0.0092929) (0.82204,-0.0085032) (0.83683,-0.0077497) (0.85112,-0.0070324) (0.86487,-0.0063508) (0.87806,-0.0057048) (0.89067,-0.005094) (0.90269,-0.0045179) (0.91408,-0.0039765) (0.92484,-0.0034694) (0.93493,-0.0029965) (0.94435,-0.0025578) (0.95308,-0.0021532) (0.9611,-0.0017827) (0.96841,-0.0014466) (0.97497,-0.001145) (0.9808,-0.00087806) (0.98586,-0.00064606) (0.99017,-0.00044925) (0.9937,-0.00028784) (0.99645,-0.00016206) (0.99842,-7.2072e-05) (0.99961,-1.8025e-05) (1,2.0634e-17)};
\addlegendentry{NACA 4415}
\end{axis}
\end{tikzpicture}

\caption{Trois profils de référence du cours, tracés à partir des points du code.}\label{fig:profils}
\end{figure}

\begin{piege}
Pour $m=0$, la position $p$ n'a aucun effet. Sur la grille $8\times8\times8$, les 8 valeurs de $p$ à $m=0$ donnent le même profil symétrique. Ce sont des doublons dans le front de Pareto, et des minima dégénérés dans le QUBO. Le code gère ce cas en renvoyant $y_c=0$ dès que $m=0$ ou $p=0$, ce qui évite aussi une division par zéro.
\end{piege}

\section{La discrétisation : distribution en cosinus}
Le contour est échantillonné en $2n+1$ points ($n=$\,\code{n\_half}, 100 par défaut, soit 200 panneaux). Les abscisses suivent une \textbf{distribution en cosinus} :
\[
x_k=\tfrac12\big(1-\cos\beta_k\big),\qquad \beta_k=\frac{k\pi}{n},\quad k=0,\dots,n .
\]
C'est la projection sur l'axe des $x$ de points également espacés sur un demi-cercle. Les points se resserrent aux deux bords, où $\dd x/\dd\beta=\tfrac12\sin\beta\to0$. Le premier intervalle vaut $\tfrac12(1-\cos\tfrac\pi n)\approx\pi^2/(4n^2)$, soit $2{,}5\cdot10^{-4}$ pour $n=100$, contre $10^{-2}$ en distribution uniforme : quarante fois plus fin.

C'est indispensable :
\begin{itemize}
\item au \textbf{bord d'attaque}, la courbure est maximale, la vitesse passe de 0 (point d'arrêt) à son maximum sur quelques millièmes de corde, et la couche limite démarre là ;
\item au \textbf{bord de fuite}, la condition de Kutta porte sur les deux derniers panneaux, et la couche limite y est la plus sensible.
\end{itemize}

\begin{figure}[htbp]
\centering
% figure générée par tools/make_figures.py : cosinus
\begin{tikzpicture}
\begin{axis}[width=0.62\linewidth,height=5cm,xmin=-0.005,xmax=0.2,ymin=0,ymax=0.065,scaled ticks=false,
  xtick={0,0.05,0.1,0.15,0.2},ytick={0,0.02,0.04,0.06},tick label style={/pgf/number format/fixed,/pgf/number format/use comma,font=\small},
  xlabel={$x/c$},ylabel={$y/c$},legend style={draw=none,font=\small,at={(0.98,0.04)},anchor=south east}]
\addplot[cmuted,thin,forget plot] coordinates {(0,0) (0.0025063,0.0087274) (0.0050125,0.012228) (0.0075188,0.014866) (0.010025,0.017057) (0.012531,0.018961) (0.015038,0.020661) (0.017544,0.022205) (0.02005,0.023625) (0.022556,0.024944) (0.025063,0.026177) (0.027569,0.027337) (0.030075,0.028433) (0.032581,0.029474) (0.035088,0.030464) (0.037594,0.031409) (0.0401,0.032313) (0.042607,0.033179) (0.045113,0.034012) (0.047619,0.034813) (0.050125,0.035585) (0.052632,0.036329) (0.055138,0.037048) (0.057644,0.037743) (0.06015,0.038416) (0.062657,0.039067) (0.065163,0.039698) (0.067669,0.04031) (0.070175,0.040904) (0.072682,0.04148) (0.075188,0.04204) (0.077694,0.042585) (0.080201,0.043114) (0.082707,0.043629) (0.085213,0.04413) (0.087719,0.044617) (0.090226,0.045092) (0.092732,0.045554) (0.095238,0.046004) (0.097744,0.046443) (0.10025,0.04687) (0.10276,0.047286) (0.10526,0.047692) (0.10777,0.048088) (0.11028,0.048474) (0.11278,0.04885) (0.11529,0.049217) (0.11779,0.049574) (0.1203,0.049923) (0.12281,0.050263) (0.12531,0.050595) (0.12782,0.050918) (0.13033,0.051234) (0.13283,0.051541) (0.13534,0.051841) (0.13784,0.052134) (0.14035,0.052419) (0.14286,0.052697) (0.14536,0.052968) (0.14787,0.053232) (0.15038,0.053489) (0.15288,0.05374) (0.15539,0.053984) (0.15789,0.054222) (0.1604,0.054454) (0.16291,0.05468) (0.16541,0.054899) (0.16792,0.055113) (0.17043,0.055321) (0.17293,0.055524) (0.17544,0.055721) (0.17794,0.055912) (0.18045,0.056098) (0.18296,0.056279) (0.18546,0.056454) (0.18797,0.056624) (0.19048,0.05679) (0.19298,0.05695) (0.19549,0.057106) (0.19799,0.057256) (0.2005,0.057402) (0.20301,0.057543) (0.20551,0.05768) (0.20802,0.057812) (0.21053,0.05794) (0.21303,0.058063) (0.21554,0.058182) (0.21805,0.058297) (0.22055,0.058407) (0.22306,0.058514) (0.22556,0.058616) (0.22807,0.058714) (0.23058,0.058808) (0.23308,0.058898) (0.23559,0.058984) (0.2381,0.059067) (0.2406,0.059146) (0.24311,0.05922) (0.24561,0.059292) (0.24812,0.059359) (0.25063,0.059423) (0.25313,0.059484) (0.25564,0.05954) (0.25815,0.059594) (0.26065,0.059644) (0.26316,0.05969) (0.26566,0.059734) (0.26817,0.059773) (0.27068,0.05981) (0.27318,0.059843) (0.27569,0.059874) (0.2782,0.059901) (0.2807,0.059924) (0.28321,0.059945) (0.28571,0.059963) (0.28822,0.059978) (0.29073,0.059989) (0.29323,0.059998) (0.29574,0.060004) (0.29825,0.060007) (0.30075,0.060007) (0.30326,0.060004) (0.30576,0.059998) (0.30827,0.05999) (0.31078,0.059979) (0.31328,0.059965) (0.31579,0.059948) (0.3183,0.059929) (0.3208,0.059907) (0.32331,0.059883) (0.32581,0.059856) (0.32832,0.059826) (0.33083,0.059794) (0.33333,0.05976) (0.33584,0.059722) (0.33835,0.059683) (0.34085,0.059641) (0.34336,0.059596) (0.34586,0.05955) (0.34837,0.0595) (0.35088,0.059449) (0.35338,0.059395) (0.35589,0.059339) (0.3584,0.059281) (0.3609,0.05922) (0.36341,0.059157) (0.36591,0.059092) (0.36842,0.059025) (0.37093,0.058955) (0.37343,0.058884) (0.37594,0.05881) (0.37845,0.058734) (0.38095,0.058656) (0.38346,0.058576) (0.38596,0.058494) (0.38847,0.05841) (0.39098,0.058324) (0.39348,0.058236) (0.39599,0.058146) (0.3985,0.058054) (0.401,0.05796) (0.40351,0.057864) (0.40602,0.057766) (0.40852,0.057667) (0.41103,0.057565) (0.41353,0.057462) (0.41604,0.057356) (0.41855,0.057249) (0.42105,0.05714) (0.42356,0.05703) (0.42607,0.056917) (0.42857,0.056803) (0.43108,0.056687) (0.43358,0.056569) (0.43609,0.05645) (0.4386,0.056328) (0.4411,0.056205) (0.44361,0.056081) (0.44612,0.055955) (0.44862,0.055827) (0.45113,0.055697) (0.45363,0.055566) (0.45614,0.055433) (0.45865,0.055299) (0.46115,0.055163) (0.46366,0.055025) (0.46617,0.054886) (0.46867,0.054746) (0.47118,0.054603) (0.47368,0.05446) (0.47619,0.054315) (0.4787,0.054168) (0.4812,0.05402) (0.48371,0.05387) (0.48622,0.053719) (0.48872,0.053566) (0.49123,0.053412) (0.49373,0.053256) (0.49624,0.0531) (0.49875,0.052941) (0.50125,0.052781) (0.50376,0.05262) (0.50627,0.052458) (0.50877,0.052294) (0.51128,0.052129) (0.51378,0.051962) (0.51629,0.051794) (0.5188,0.051625) (0.5213,0.051454) (0.52381,0.051282) (0.52632,0.051109) (0.52882,0.050934) (0.53133,0.050759) (0.53383,0.050581) (0.53634,0.050403) (0.53885,0.050223) (0.54135,0.050043) (0.54386,0.04986) (0.54637,0.049677) (0.54887,0.049492) (0.55138,0.049307) (0.55388,0.04912) (0.55639,0.048931) (0.5589,0.048742) (0.5614,0.048551) (0.56391,0.048359) (0.56642,0.048166) (0.56892,0.047972) (0.57143,0.047777) (0.57393,0.047581) (0.57644,0.047383) (0.57895,0.047184) (0.58145,0.046984) (0.58396,0.046783) (0.58647,0.046581) (0.58897,0.046378) (0.59148,0.046173) (0.59398,0.045968) (0.59649,0.045761) (0.599,0.045554) (0.6015,0.045345) (0.60401,0.045135) (0.60652,0.044924) (0.60902,0.044712) (0.61153,0.044499) (0.61404,0.044285) (0.61654,0.04407) (0.61905,0.043854) (0.62155,0.043637) (0.62406,0.043419) (0.62657,0.043199) (0.62907,0.042979) (0.63158,0.042758) (0.63409,0.042535) (0.63659,0.042312) (0.6391,0.042088) (0.6416,0.041862) (0.64411,0.041636) (0.64662,0.041409) (0.64912,0.04118) (0.65163,0.040951) (0.65414,0.040721) (0.65664,0.040489) (0.65915,0.040257) (0.66165,0.040024) (0.66416,0.03979) (0.66667,0.039554) (0.66917,0.039318) (0.67168,0.039081) (0.67419,0.038843) (0.67669,0.038604) (0.6792,0.038364) (0.6817,0.038123) (0.68421,0.037882) (0.68672,0.037639) (0.68922,0.037395) (0.69173,0.03715) (0.69424,0.036905) (0.69674,0.036658) (0.69925,0.036411) (0.70175,0.036163) (0.70426,0.035913) (0.70677,0.035663) (0.70927,0.035412) (0.71178,0.03516) (0.71429,0.034907) (0.71679,0.034653) (0.7193,0.034398) (0.7218,0.034143) (0.72431,0.033886) (0.72682,0.033629) (0.72932,0.03337) (0.73183,0.033111) (0.73434,0.032851) (0.73684,0.03259) (0.73935,0.032328) (0.74185,0.032065) (0.74436,0.031801) (0.74687,0.031536) (0.74937,0.031271) (0.75188,0.031004) (0.75439,0.030737) (0.75689,0.030469) (0.7594,0.0302) (0.7619,0.02993) (0.76441,0.029659) (0.76692,0.029387) (0.76942,0.029114) (0.77193,0.028841) (0.77444,0.028566) (0.77694,0.028291) (0.77945,0.028014) (0.78195,0.027737) (0.78446,0.027459) (0.78697,0.02718) (0.78947,0.0269) (0.79198,0.02662) (0.79449,0.026338) (0.79699,0.026055) (0.7995,0.025772) (0.80201,0.025488) (0.80451,0.025202) (0.80702,0.024916) (0.80952,0.024629) (0.81203,0.024341) (0.81454,0.024052) (0.81704,0.023763) (0.81955,0.023472) (0.82206,0.02318) (0.82456,0.022888) (0.82707,0.022595) (0.82957,0.0223) (0.83208,0.022005) (0.83459,0.021709) (0.83709,0.021412) (0.8396,0.021114) (0.84211,0.020815) (0.84461,0.020516) (0.84712,0.020215) (0.84962,0.019913) (0.85213,0.019611) (0.85464,0.019307) (0.85714,0.019003) (0.85965,0.018697) (0.86216,0.018391) (0.86466,0.018084) (0.86717,0.017776) (0.86967,0.017467) (0.87218,0.017157) (0.87469,0.016846) (0.87719,0.016534) (0.8797,0.016221) (0.88221,0.015907) (0.88471,0.015592) (0.88722,0.015276) (0.88972,0.014959) (0.89223,0.014642) (0.89474,0.014323) (0.89724,0.014003) (0.89975,0.013683) (0.90226,0.013361) (0.90476,0.013038) (0.90727,0.012715) (0.90977,0.01239) (0.91228,0.012064) (0.91479,0.011738) (0.91729,0.01141) (0.9198,0.011081) (0.92231,0.010752) (0.92481,0.010421) (0.92732,0.010089) (0.92982,0.0097563) (0.93233,0.0094225) (0.93484,0.0090877) (0.93734,0.0087517) (0.93985,0.0084148) (0.94236,0.0080768) (0.94486,0.0077377) (0.94737,0.0073976) (0.94987,0.0070564) (0.95238,0.0067141) (0.95489,0.0063708) (0.95739,0.0060264) (0.9599,0.0056808) (0.96241,0.0053342) (0.96491,0.0049865) (0.96742,0.0046377) (0.96992,0.0042878) (0.97243,0.0039367) (0.97494,0.0035846) (0.97744,0.0032313) (0.97995,0.0028768) (0.98246,0.0025213) (0.98496,0.0021646) (0.98747,0.0018067) (0.98997,0.0014477) (0.99248,0.0010875) (0.99499,0.00072621) (0.99749,0.00036369) (1,-1.6653e-17)};
\addplot[only marks,mark=*,mark size=1.6pt,cblue] coordinates {(0,0) (0.0027391,0.0091145) (0.010926,0.01777) (0.024472,0.025893) (0.043227,0.033389) (0.066987,0.040145) (0.095492,0.046049) (0.12843,0.050995) (0.16543,0.054901) (0.20611,0.057712) (0.25,0.059407) (0.29663,0.060005) (0.34549,0.059557) (0.39604,0.058144) (0.44774,0.055872) (0.5,0.052862) (0.55226,0.049241) (0.60396,0.04514) (0.65451,0.040686) (0.70337,0.036002) (0.75,0.031204) (0.79389,0.026405) (0.83457,0.021712) (0.87157,0.017232) (0.90451,0.013071) (0.93301,0.0093315) (0.95677,0.0061118) (0.97553,0.0035014) (0.98907,0.0015769) (0.99726,0.00039741) (1,-1.6653e-17)};
\addlegendentry{cosinus (60 panneaux)}
\addplot[only marks,mark=square*,mark size=1.6pt,corange] coordinates {(0,0) (0.033333,0.029776) (0.066667,0.040067) (0.1,0.046828) (0.13333,0.051602) (0.16667,0.055007) (0.2,0.057373) (0.23333,0.058907) (0.26667,0.05975) (0.3,0.060007) (0.33333,0.05976) (0.36667,0.059072) (0.4,0.057998) (0.43333,0.056581) (0.46667,0.054858) (0.5,0.052862) (0.53333,0.050617) (0.56667,0.048147) (0.6,0.04547) (0.63333,0.042602) (0.66667,0.039554) (0.7,0.036337) (0.73333,0.032955) (0.76667,0.029414) (0.8,0.025715) (0.83333,0.021857) (0.86667,0.017838) (0.9,0.01365) (0.93333,0.0092887) (0.96667,0.0047425) (1,-1.6653e-17)};
\addlegendentry{uniforme (60 panneaux)}
\end{axis}
\end{tikzpicture}

\caption{Zoom sur le bord d'attaque d'un NACA 0012 avec 60 panneaux. La distribution en cosinus place 5 points avant $x=0{,}05$, la distribution uniforme seulement 2 : elle rate complètement l'arrondi du nez.}\label{fig:cosinus}
\end{figure}

\paragraph{Ordre de parcours.} Les points sont rangés comme dans XFOIL : du bord de fuite par l'extrados jusqu'au bord d'attaque, puis par l'intrados jusqu'au bord de fuite. Le contour est donc parcouru dans le sens trigonométrique, et le point de bord de fuite apparaît deux fois (premier et dernier). Avec ce sens, la \textbf{normale sortante} d'un segment de tangente unitaire $(t_x,t_y)$ est $(t_y,-t_x)$ : elle est à droite quand on avance le long du contour. Le même fichier de coordonnées sert d'entrée à XFOIL et au maillage de la SER.

\begin{codebox}[ : \code{aero/naca4.py}]
\begin{lstlisting}
beta = np.linspace(0.0, np.pi, n_half + 1)
x = 0.5 * (1 - np.cos(beta))                 # distribution en cosinus
yt = thickness(x, t, closed_te)              # a4 = -0.1036 : bord de fuite ferme
yc, dyc = camber_line(x, m, p)
th = np.arctan(dyc)
xu, yu = x - yt * np.sin(th), yc + yt * np.cos(th)   # extrados
xl, yl = x + yt * np.sin(th), yc - yt * np.cos(th)   # intrados
X = np.concatenate([xu[::-1], xl[1:]])       # BF extrados -> BA -> BF intrados
Y = np.concatenate([yu[::-1], yl[1:]])
\end{lstlisting}
\code{xu[::-1]} parcourt l'extrados à l'envers (du bord de fuite vers le bord d'attaque), et \code{xl[1:]} évite de dupliquer le point de bord d'attaque. On obtient $X[0]=X[-1]=1$.
\end{codebox}

\begin{projet}
La paramétrisation NACA est pauvre (trois degrés de liberté, des formes toujours « raisonnables ») mais interprétable et énumérable : $8^3=512$ profils. Le rayon de bord d'attaque $r_{BA}=1{,}1019\,t^2$ est exporté dans \code{AeroResult.r\_le}, car c'est probablement le lien physique le plus direct entre les deux objectifs. Il pilote le pic d'aspiration côté aérodynamique, et vraisemblablement le retour radar de face côté SER, ce que l'équipe de la phase~2 doit confirmer. Pour un cas industriel, on passerait à des paramétrisations plus riches (CST de Kulfan, B-splines), donc à plus de variables et plus de qubits.
\end{projet}

\exercices
\begin{exercice}
Pour le NACA 2412 : (a) calculer $y_c(p)$ et la pente de la ligne moyenne en $x=p$ par les deux formules, et vérifier le raccord ; (b) calculer le rayon de bord d'attaque ; (c) que vaut la courbure de la ligne moyenne juste avant et juste après $x=p$ ?
\end{exercice}
\begin{exercice}
Avec $n=100$, comparer l'abscisse du premier point après le bord d'attaque (ou du dernier avant le bord de fuite) en distribution cosinus et en distribution uniforme. Pourquoi le point de bord d'attaque lui-même n'est-il pas un point de collocation ?
\end{exercice}

\chapter{L'écoulement potentiel et ses singularités}\label{ch:potentiel}
\iterbanner{V0}

Ce chapitre rassemble ce dont la méthode des panneaux a besoin : le problème de Laplace extérieur, les singularités élémentaires, l'idée de les poser sur la frontière, et ce qui rend la solution unique, la condition de Kutta.

\section{Le problème à résoudre}
Hors de la couche limite et du sillage, l'écoulement est incompressible et irrotationnel. Il dérive donc d'un potentiel, $\vect V=\nabla\phi$, et la conservation de la masse $\nabla\cdot\vect V=0$ donne l'équation de Laplace :
\[
\nabla^2\phi=0\quad\text{dans le fluide},\qquad
\frac{\partial\phi}{\partial n}=\vect V\cdot\vect n=0\quad\text{sur le profil},\qquad
\nabla\phi\to\vect U_\infty\quad\text{à l'infini}.
\]
Le problème est \textbf{linéaire} et \textbf{elliptique}. La pression s'obtient ensuite par Bernoulli, valable partout puisque l'écoulement est irrotationnel :
\[
C_p=\frac{p-p_\infty}{\tfrac12\rho U_\infty^2}=1-\left(\frac{V}{U_\infty}\right)^2 .
\]
Sur la paroi, la vitesse est tangente, donc $V=|V_t|$ et $C_p=1-V_t^2$ en variables adimensionnées.

\section{Les singularités élémentaires}
La linéarité permet de superposer des solutions élémentaires de Laplace. En 2D, deux suffisent, toutes deux centrées en un point et de vitesse en $1/r$ :
\begin{center}
\renewcommand{\arraystretch}{1.4}
\begin{tabular}{lccl}
\toprule
 & potentiel $\phi$ & vitesse & porte\\
\midrule
source $\sigma$ & $\dfrac{\sigma}{2\pi}\ln r$ & $\dfrac{\sigma}{2\pi r}\,\vect e_r$ & un débit $\sigma$\\
tourbillon $\Gamma$ & $\dfrac{\Gamma}{2\pi}\,\vartheta$ & $\dfrac{\Gamma}{2\pi r}\,\vect e_\vartheta$ & une circulation $\Gamma$ (sens trigo.)\\
\bottomrule
\end{tabular}
\end{center}
Le tourbillon n'est rien d'autre qu'une source dont la vitesse a été tournée de $90^\circ$. Cette remarque servira à obtenir les intégrales de panneau du tourbillon gratuitement à partir de celles de la source.

\begin{figure}[htbp]
\centering
\begin{tikzpicture}
\begin{scope}
\fill[cblue] (0,0) circle (2pt);
\foreach \a in {0,30,...,330}{\draw[cblue,-{Stealth[length=1.6mm]}] (\a:0.25) -- (\a:1.1);}
\node[font=\small] at (0,-1.5) {source $\sigma>0$ : débit sortant};
\end{scope}
\begin{scope}[xshift=6.5cm]
\fill[corange] (0,0) circle (2pt);
\foreach \r in {0.4,0.75,1.1}{\draw[corange,thick,postaction={decorate},decoration={markings,mark=at position 0.25 with {\arrow{Stealth[length=1.8mm]}},mark=at position 0.75 with {\arrow{Stealth[length=1.8mm]}}}] (\r,0) arc (0:360:\r);}
\node[font=\small] at (0,-1.5) {tourbillon $\Gamma>0$ : rotation trigonométrique};
\end{scope}
\end{tikzpicture}
\caption{Les deux singularités élémentaires 2D. Leur vitesse décroît en $1/r$ et satisfait automatiquement $\nabla^2\phi=0$ hors du centre et la condition à l'infini.}
\end{figure}

\section{Poser les singularités sur la frontière}
Plutôt que de chercher $\phi$ dans tout le plan, on cherche des \emph{distributions} de singularités le long du contour du profil. Une source répartie avec une densité linéique $\sigma(s)$ forme une \textbf{nappe de sources}, et un tourbillon réparti avec une densité $\gamma(s)$ une \textbf{nappe tourbillonnaire}. Le champ qu'elles induisent satisfait Laplace partout hors de la nappe et décroît à l'infini. Il ne reste qu'à choisir les densités pour que la condition de glissement soit vérifiée sur la paroi.

Ce choix est légitime : la troisième identité de Green montre que toute fonction harmonique à l'extérieur d'un contour s'écrit comme le potentiel d'une couche de sources plus une couche de doublets posées sur ce contour, et une couche de doublets équivaut à une nappe tourbillonnaire. Conséquence pratique majeure : on ne maille que la paroi, une courbe 1D pour un problème 2D, et la condition à l'infini est vérifiée d'office. C'est la famille des \textbf{méthodes d'éléments de frontière} (BEM).

\paragraph{Les sauts à travers une nappe.} En traversant une nappe :
\begin{itemize}
\item une nappe de sources fait sauter la vitesse \emph{normale} de $\sigma$ : elle vaut $+\sigma/2$ d'un côté et $-\sigma/2$ de l'autre, comptée selon la normale. La vitesse tangentielle est continue ;
\item une nappe tourbillonnaire fait sauter la vitesse \emph{tangentielle} de $\gamma$, soit $\pm\gamma/2$ de part et d'autre. La vitesse normale est continue.
\end{itemize}
Ces demi-sauts réapparaîtront au chapitre suivant comme « auto-influence » d'un panneau sur son propre point de collocation.

\section{Non-unicité et condition de Kutta}
Autour d'un profil, le domaine fluide n'est pas simplement connexe : un contour qui entoure le profil ne peut pas être réduit à un point sans traverser l'obstacle. On peut alors ajouter à toute solution un tourbillon libre de circulation $\Gamma$ quelconque : il vérifie Laplace, la condition à l'infini, et sa vitesse peut être rendue tangente à la paroi. Il existe donc une \emph{famille} de solutions, une par valeur de $\Gamma$.

Pour presque toutes ces valeurs, l'écoulement contourne le bord de fuite aigu, ce qui demande une vitesse infinie à l'angle (figure~\ref{fig:kutta}). Physiquement, la viscosité l'interdit : au démarrage, un tourbillon de démarrage est lâché et la circulation autour du profil s'ajuste jusqu'à ce que l'écoulement quitte proprement le bord de fuite.

\begin{definition}[ (condition de Kutta)]
Parmi toutes les solutions potentielles, on retient celle où l'écoulement quitte le bord de fuite sans le contourner, avec des vitesses finies et égales sur l'extrados et l'intrados. Cette condition fixe $\Gamma$, donc la portance.
\end{definition}

\begin{figure}[htbp]
\centering
\begin{tikzpicture}[x=5cm,y=5cm]
\foreach \xs/\titre/\k in {0/{$\Gamma=0$ : l'écoulement contourne le bord de fuite}/0, 1.3/{$\Gamma$ de Kutta : départ tangent au bord de fuite}/1}{
\begin{scope}[xshift=\xs*5cm]
\fill[cmuted!25,draw=cink,thick] plot[smooth,tension=0.6] coordinates {(1,0) (0.75,0.05) (0.45,0.085) (0.2,0.08) (0.05,0.045) (0,0) (0.05,-0.035) (0.2,-0.045) (0.45,-0.035) (0.75,-0.012) (1,0)};
\draw[-{Stealth[length=2mm]},cink,thick] (-0.25,-0.05) -- (-0.08,-0.05) node[midway,above,font=\scriptsize] {$U_\infty$, $\alpha$};
\ifnum\k=0
  \draw[cblue,thick,-{Stealth[length=1.8mm]}] (0.6,-0.11) .. controls (0.95,-0.07) and (1.06,-0.02) .. (1.0,0.02) .. controls (0.95,0.05) and (0.88,0.05) .. (0.82,0.045);
  \fill[corange] (0.82,0.045) circle (0.012);
  \node[corange,font=\scriptsize,align=center] at (0.75,0.17) {point d'arrêt arrière\\ sur l'extrados};
  \node[cblue,font=\scriptsize] at (1.08,-0.1) {$V\to\infty$};
\else
  \draw[cblue,thick,-{Stealth[length=1.8mm]}] (0.6,0.14) .. controls (0.85,0.1) and (0.95,0.04) .. (1.25,-0.02);
  \draw[cblue,thick,-{Stealth[length=1.8mm]}] (0.6,-0.1) .. controls (0.85,-0.07) and (0.95,-0.03) .. (1.25,-0.025);
  \fill[corange] (1,0) circle (0.012);
  \node[corange,font=\scriptsize] at (1.05,0.08) {bord de fuite};
\fi
\node[font=\small,text width=4.6cm,align=center] at (0.5,-0.27) {\titre};
\end{scope}}
\end{tikzpicture}
\caption{La condition de Kutta sélectionne, dans la famille des solutions potentielles, celle qui quitte proprement le bord de fuite.}\label{fig:kutta}
\end{figure}

La condition de Kutta est donc un \emph{morceau de physique visqueuse} injecté dans un modèle non visqueux. Elle explique aussi pourquoi le $C_L$ réel est un peu plus faible que le $C_L$ potentiel : la couche limite épaissie près du bord de fuite décambre le profil effectif, ce que notre couplage faible ne voit pas (ch.~\ref{ch:validation}).

\section{Portance, traînée, et deux résultats de référence}
\paragraph{Kutta-Joukowski et d'Alembert.} Pour un écoulement potentiel 2D autour d'un corps de circulation $\Gamma$ (comptée dans le sens horaire), la force par unité d'envergure est perpendiculaire à $\vect U_\infty$ :
\[
L'=\rho\,U_\infty\,\Gamma\quad\Longrightarrow\quad C_L=\frac{2\Gamma}{U_\infty c},\qquad D'=0 .
\]
La traînée nulle est le \textbf{paradoxe de d'Alembert}. Il a deux conséquences pour nous. D'une part, le $C_{D,p}$ obtenu en intégrant $C_p$ sur le contour discret n'est pas une traînée, mais une erreur de discrétisation qui doit tendre vers 0 : c'est un excellent indicateur de convergence. D'autre part, toute la traînée devra venir de la couche limite (ch.~\ref{ch:trainee}).

\paragraph{Théorie des profils minces.} Pour un profil mince et peu cambré, la portance est affine en incidence :
\[
C_L=2\pi\big(\alpha-\alpha_{L0}\big),\qquad
\alpha_{L0}=-\frac1\pi\int_0^\pi\frac{\dd y_c}{\dd x}\,\big(\cos\vartheta_0-1\big)\,\dd\vartheta_0,\qquad x=\tfrac12(1-\cos\vartheta_0).
\]
Pour le NACA 2412, $\alpha_{L0}=-2{,}08^\circ$. L'épaisseur, absente de cette théorie, augmente la pente d'environ $0{,}77\,t$ en relatif. Ces deux résultats serviront de contrôle à la méthode des panneaux (figure~\ref{fig:cl-alpha}).

\begin{resultat}
\begin{itemize}
\item Le problème potentiel est linéaire : on superpose des sources et des tourbillons posés sur la paroi, et on ne maille que le contour.
\item Il a une infinité de solutions, une par circulation. La condition de Kutta au bord de fuite choisit la bonne et fixe la portance.
\item La traînée potentielle est nulle. Une valeur non nulle de $C_{D,p}$ mesure l'erreur de discrétisation.
\end{itemize}
\end{resultat}

\exercices
\begin{exercice}
Un profil de corde $c=1$~m vole à $U_\infty=50$~m/s avec $C_L=0{,}5$. Calculer la circulation $\Gamma$ autour du profil, et la portance par mètre d'envergure pour $\rho=1{,}2$~kg/m\textsuperscript{3}.
\end{exercice}
\begin{exercice}
Par un bilan de débit sur un petit rectangle aplati qui enjambe une nappe de sources de densité $\sigma$, montrer que la vitesse normale saute de $\sigma$ à la traversée. Pourquoi, par symétrie, vaut-elle $+\sigma/2$ d'un côté et $-\sigma/2$ de l'autre pour une nappe plane isolée ?
\end{exercice}

\chapter{La méthode des panneaux de Hess-Smith}\label{ch:panneaux}
\iterbanner{V0}

\section{Le principe}
On remplace le contour par $N$ segments rectilignes, les \textbf{panneaux}. Sur chacun, on pose :
\begin{itemize}
\item une nappe de sources d'intensité \textbf{constante} $\sigma_j$, propre au panneau $j$ ;
\item une nappe tourbillonnaire d'intensité $\gamma$, constante, et \textbf{la même sur tous les panneaux}.
\end{itemize}
Cela fait $N+1$ inconnues : $\sigma_1,\dots,\sigma_N$ et $\gamma$. On écrit $N+1$ équations :
\begin{itemize}
\item la non-pénétration $\vect V\cdot\vect n_i=0$ au milieu de chaque panneau, appelé \textbf{point de collocation} : $N$ équations ;
\item la condition de Kutta : 1 équation.
\end{itemize}

\paragraph{Pourquoi des sources et un seul tourbillon ?} Les sources seules ont une circulation totale nulle, donc ne portent pas (Kutta-Joukowski). Il faut de la circulation, mais un tourbillon libre par panneau donnerait $2N$ inconnues pour $N+1$ équations : le problème serait sous-déterminé. Hess et Smith (1967) prennent le minimum : une seule intensité tourbillonnaire uniforme, que la condition de Kutta fixe. C'est la formulation la plus simple qui donne une solution unique avec portance.

\section{Géométrie des panneaux}
Le panneau $j$ va du point $(x_j,y_j)$ au point $(x_{j+1},y_{j+1})$ du contour. On en tire (figure~\ref{fig:panneau}) :
\[
L_j=\sqrt{(x_{j+1}-x_j)^2+(y_{j+1}-y_j)^2},\qquad
\vect t_j=\frac{(x_{j+1}-x_j,\ y_{j+1}-y_j)}{L_j},\qquad
\vect n_j=(t_{y,j},\,-t_{x,j}),
\]
et le point de collocation $(\bar x_j,\bar y_j)$ au milieu. Comme le contour est parcouru dans le sens trigonométrique, $\vect n_j$ (à droite de $\vect t_j$) pointe vers l'extérieur, c'est-à-dire vers le fluide.

\begin{figure}[htbp]
\centering
\begin{tikzpicture}[scale=1.1]
% panneau
\draw[cink,very thick] (0,0) -- (4,0);
\fill[cink] (0,0) circle (1.6pt) node[below left,font=\small] {début $(x_j,y_j)$};
\fill[cink] (4,0) circle (1.6pt) node[below right,font=\small] {fin $(x_{j+1},y_{j+1})$};
\fill[cgreen] (2,0) circle (1.8pt) node[below left=3pt and 1pt,font=\small,cgreen] {collocation};
\draw[arr,cblue] (2,0) -- (3,0) node[above,font=\small,cblue] {$\vect t_j$};
\draw[arr,cblue] (2,0) -- (2,-0.9) node[right,font=\small,cblue] {$\vect n_j$ (vers le fluide)};
\draw[arr,cmuted] (2,0) -- (2,0.7) node[left,font=\small,cmuted] {$\vect e_\eta$};
% repère local
\draw[arr,cmuted] (0,0) -- (0.9,0) ; \node[cmuted,font=\small] at (0.8,0.2) {$\xi$};
% point P
\coordinate (P) at (2.9,2.2);
\fill[corange] (P) circle (1.8pt) node[above right,font=\small,corange] {$P=(\xi,\eta)$};
\draw[corange,thin] (0,0) -- (P) node[pos=0.5,above left,font=\small] {$r_1$};
\draw[corange,thin] (4,0) -- (P) node[pos=0.5,right,font=\small] {$r_2$};
\draw[corange] (0.7,0) arc (0:37.2:0.7); \node[corange,font=\small] at (0.95,0.32) {$\theta_1$};
\draw[corange] (4.5,0) arc (0:116.6:0.5); \node[corange,font=\small] at (4.55,0.62) {$\theta_2$};
\draw[cmuted,dashed] (4,0) -- (4.7,0);
\node[font=\small,cmuted,align=left] at (-1.5,1.3) {intérieur du profil\\ (côté $\eta>0$)};
\node[font=\small,cmuted] at (-1.2,-1.0) {fluide ($\eta<0$)};
\end{tikzpicture}
\caption{Un panneau et son repère local. $\xi$ suit la tangente, $\eta$ la normale gauche, qui pointe vers l'intérieur du profil. Les angles $\theta_1$, $\theta_2$ sont ceux sous lesquels les extrémités voient le point $P$.}\label{fig:panneau}
\end{figure}

\section{La vitesse induite par un panneau}
Plaçons le panneau sur l'axe $\xi\in[0,L]$, $\eta=0$, et calculons la vitesse qu'il induit en un point $P=(\xi,\eta)$.

\paragraph{Source.} Un élément $\dd\xi'$ en $(\xi',0)$ est une source de débit $\sigma\,\dd\xi'$. Sa vitesse en $P$ est radiale :
\[
\dd\vect V=\frac{\sigma\,\dd\xi'}{2\pi}\,\frac{(\xi-\xi',\ \eta)}{(\xi-\xi')^2+\eta^2}.
\]
En intégrant :
\[
u_\sigma=\frac{\sigma}{2\pi}\int_0^L\frac{\xi-\xi'}{(\xi-\xi')^2+\eta^2}\,\dd\xi'
=\frac{\sigma}{4\pi}\Big[-\ln\big((\xi-\xi')^2+\eta^2\big)\Big]_0^L
=\frac{\sigma}{2\pi}\ln\frac{r_1}{r_2},
\]
\[
v_\sigma=\frac{\sigma}{2\pi}\int_0^L\frac{\eta}{(\xi-\xi')^2+\eta^2}\,\dd\xi'=\frac{\sigma}{2\pi}\big(\theta_2-\theta_1\big),
\qquad \theta_k=\operatorname{atan2}(\eta,\ \xi-\xi_k).
\]
La composante tangentielle dépend du rapport des distances aux extrémités, la composante normale de l'angle sous lequel on voit le panneau.

\paragraph{Tourbillon.} La vitesse d'un tourbillon est celle d'une source tournée de $+90^\circ$ : $(u,v)\to(-v,u)$. On obtient sans nouveau calcul
\[
u_\gamma=-\frac{\gamma}{2\pi}\big(\theta_2-\theta_1\big),\qquad v_\gamma=\frac{\gamma}{2\pi}\ln\frac{r_1}{r_2}.
\]

\paragraph{Auto-influence.} Quand $P$ est le point de collocation du panneau lui-même, $r_1=r_2$ donc $\ln(r_1/r_2)=0$. Pour l'angle, il faut savoir de quel côté on s'approche. Le fluide est du côté $\eta<0$. En faisant $\eta\to0^-$ avec $0<\xi<L$, on a $\theta_1\to0$ et $\theta_2\to-\pi$, donc $\theta_2-\theta_1=-\pi$. On retrouve les demi-sauts du chapitre précédent :
\begin{itemize}
\item la source induit sur elle-même $v_\sigma=-\sigma/2$ selon $\vect e_\eta$, soit $+\sigma/2$ selon la normale sortante : une source positive « souffle » vers le fluide ;
\item le tourbillon induit $u_\gamma=+\gamma/2$ selon la tangente.
\end{itemize}

\paragraph{Retour au repère global.} Avec $\vect e_\xi=\vect t_j=(t_x,t_y)$ et $\vect e_\eta=(-t_y,t_x)$, une vitesse locale $(u,v)$ vaut $(u\,t_x-v\,t_y,\ u\,t_y+v\,t_x)$ dans le repère global.

\begin{codebox}[ : \code{panel.influence\_matrices}]
Tout est vectorisé : on calcule d'un coup les $N\times N$ couples (collocation $i$, panneau $j$).
\begin{lstlisting}
dx = P.xc[:, None] - P.xs[None, :]          # collocation i - debut du panneau j
dy = P.yc[:, None] - P.ys[None, :]
xi  =  dx * tx + dy * ty                    # coordonnees locales du panneau j
eta = -dx * ty + dy * tx
r1, r2 = np.hypot(xi, eta), np.hypot(xi - L, eta)
lnr = np.log(r1 / r2)
dth = np.arctan2(eta, xi - L) - np.arctan2(eta, xi)    # theta2 - theta1
idx = np.arange(P.n)
lnr[idx, idx] = 0.0                         # auto-influence, cote fluide
dth[idx, idx] = -np.pi
us_l, vs_l = lnr / (2*np.pi), dth / (2*np.pi)      # source unitaire
uv_l, vv_l = -dth / (2*np.pi), lnr / (2*np.pi)     # tourbillon unitaire
\end{lstlisting}
Les quatre matrices renvoyées (\code{usx, usy, uvx, uvy}) contiennent les composantes globales de la vitesse induite en $i$ par une source unitaire et par un tourbillon unitaire posés sur le panneau $j$. Elles ne dépendent que de la géométrie, pas de l'incidence.
\end{codebox}

\section{Le système linéaire}
Notons $A^n_{ij}$ et $A^t_{ij}$ les composantes normale ($\cdot\,\vect n_i$) et tangentielle ($\cdot\,\vect t_i$) de la vitesse induite en $i$ par une source unitaire sur le panneau $j$, et $B^n_{ij}$, $B^t_{ij}$ les mêmes pour un tourbillon unitaire. La vitesse au point de collocation $i$ s'écrit
\[
\vect V_i=\vect U_\infty+\sum_j\sigma_j\,(\text{vitesse de la source }j)+\gamma\sum_j(\text{vitesse du tourbillon }j),
\qquad \vect U_\infty=(\cos\alpha,\ \sin\alpha).
\]

\paragraph{Non-pénétration.} Pour $i=1,\dots,N$ :
\[
\sum_{j=1}^N A^n_{ij}\,\sigma_j+\Big(\sum_{j=1}^N B^n_{ij}\Big)\gamma=-\vect U_\infty\cdot\vect n_i .
\]

\paragraph{Kutta discrète.} Les deux panneaux qui touchent le bord de fuite sont le premier (extrados, $\vect t_1$ dirigé vers l'amont) et le dernier (intrados, $\vect t_N$ dirigé vers l'aval). Leurs tangentes sont donc presque opposées. L'écoulement réel quitte le bord de fuite vers l'aval des deux côtés : $V_{t,1}<0$ et $V_{t,N}>0$. Demander des vitesses de même norme revient à
\begin{gather*}
V_{t,1}+V_{t,N}=0\\
\Longleftrightarrow\quad
\sum_j\big(A^t_{1j}+A^t_{Nj}\big)\sigma_j+\Big(\sum_jB^t_{1j}+\sum_jB^t_{Nj}\Big)\gamma=-\vect U_\infty\cdot(\vect t_1+\vect t_N).
\end{gather*}

\paragraph{Forme matricielle.}
\[
\begin{pmatrix}
A^n & B^n\vect 1\\[2pt]
A^t_{1\cdot}+A^t_{N\cdot} & \vect 1^\top\big(B^t_{1\cdot}+B^t_{N\cdot}\big)^\top
\end{pmatrix}
\begin{pmatrix}\boldsymbol\sigma\\ \gamma\end{pmatrix}
=-\begin{pmatrix}\big(\vect U_\infty\cdot\vect n_i\big)_{i}\\ \vect U_\infty\cdot(\vect t_1+\vect t_N)\end{pmatrix}.
\]
C'est un système \textbf{dense} de taille $N+1=201$. La matrice ne dépend que de la géométrie ; seul le second membre dépend de $\alpha$.

\begin{piege}
Écrire $V_{t,1}=V_{t,N}$ au lieu de $V_{t,1}+V_{t,N}=0$ est l'erreur classique : avec des tangentes opposées, cela imposerait des vitesses physiques \emph{opposées} au bord de fuite. Le symptôme est un $C_L$ faux et un $C_{L,C_p}$ très différent de $C_{L,KJ}$.
\end{piege}

\begin{codebox}[ : \code{panel.solve}]
\begin{lstlisting}
A = np.zeros((N + 1, N + 1))
A[:N, :N] = Asn                          # sources -> vitesse normale
A[:N, N]  = Avn.sum(axis=1)              # tourbillon uniforme -> vitesse normale
A[N, :N]  = Ast[0] + Ast[-1]             # Kutta : Vt_1 + Vt_N = 0
A[N, N]   = Avt[0].sum() + Avt[-1].sum()
lu = np.linalg.inv(A)                    # reutilisable pour toute incidence
b[:N] = -(Uinf[0] * P.nx + Uinf[1] * P.ny)
b[N]  = -(Uinf[0] * (P.tx[0] + P.tx[-1]) + Uinf[1] * (P.ty[0] + P.ty[-1]))
q = lu @ b;  sigma, gamma = q[:N], q[N]
\end{lstlisting}
Le dictionnaire optionnel \code{cache} conserve les panneaux, les matrices d'influence et l'inverse. Changer d'incidence ne coûte alors qu'un produit matrice-vecteur, ce que \code{alpha\_for\_cl} exploite (ch.~\ref{ch:modele}). L'inverse explicite est acceptable ici : $201\times201$, une fois par profil.
\end{codebox}

\section{Ce qu'on tire de la solution}
\paragraph{Vitesse tangentielle et pression.} Aux points de collocation,
\[
V_{t,i}=\vect U_\infty\cdot\vect t_i+\sum_jA^t_{ij}\sigma_j+\gamma\sum_jB^t_{ij},\qquad C_{p,i}=1-V_{t,i}^2 .
\]
$V_{t,i}$ est \emph{signée} selon $\vect t_i$ : elle est négative sur l'extrados (l'écoulement va du bord d'attaque vers le bord de fuite, à l'opposé du sens de parcours) et positive sur l'intrados. Le changement de signe repère le point d'arrêt (ch.~\ref{ch:cl}).

\paragraph{Efforts.} La force de pression sur le panneau $i$ vaut $-C_{p,i}\,\vect n_i\,L_i$ (en unités de $\tfrac12\rho U_\infty^2c$). On somme, puis on projette :
\[
\vect F=-\sum_iC_{p,i}\,\vect n_i\,L_i,\qquad
C_L=F_y\cos\alpha-F_x\sin\alpha,\qquad C_{D,p}=F_x\cos\alpha+F_y\sin\alpha,
\]
\[
C_{M,c/4}=-\sum_i\Big[(\bar x_i-\tfrac14)\,f_{y,i}-\bar y_i\,f_{x,i}\Big]\qquad\text{(positif à cabrer).}
\]

\paragraph{Deux contrôles gratuits.} La circulation totale de la nappe tourbillonnaire vaut $\gamma$ fois le périmètre, comptée dans le sens trigonométrique. La portance correspond à une circulation horaire $\Gamma=-\gamma\,\ell_{tot}$, d'où $C_{L,KJ}=-2\gamma\,\ell_{tot}$ par Kutta-Joukowski. On doit avoir $C_{L,KJ}\approx C_{L,C_p}$, et $C_{D,p}\to0$ (d'Alembert). Les deux écarts mesurent l'erreur de discrétisation.

\begin{figure}[htbp]
\centering
% figure générée par tools/make_figures.py : cp
\begin{tikzpicture}
\begin{groupplot}[group style={group size=3 by 1,horizontal sep=0.9cm,y descriptions at=edge left},
  width=0.34\linewidth,height=5.2cm,y dir=reverse,xmin=0,xmax=1,ymin=-2.6,ymax=1.1,
  xlabel={$x/c$},ylabel={$C_p$},title style={font=\small},tick label style={font=\scriptsize}]
\nextgroupplot[title={NACA 0012, $\alpha=0^\circ$}]
\addplot[cblue,very thick] coordinates {(0.00012336,0.98707) (0.00061668,0.88936) (0.0016028,0.72336) (0.0030808,0.53089) (0.0050492,0.34529) (0.0075061,0.18392) (0.010449,0.051559) (0.013875,-0.05372) (0.017781,-0.13625) (0.022162,-0.2006) (0.027016,-0.25076) (0.032336,-0.28994) (0.038117,-0.3206) (0.044355,-0.34463) (0.051042,-0.36343) (0.058172,-0.37807) (0.065738,-0.38937) (0.073733,-0.39793) (0.082148,-0.40422) (0.090976,-0.40862) (0.10021,-0.4114) (0.10983,-0.4128) (0.11984,-0.41299) (0.13023,-0.41212) (0.14098,-0.4103) (0.15209,-0.40765) (0.16354,-0.40424) (0.17532,-0.40015) (0.18742,-0.39545) (0.19983,-0.39018) (0.21253,-0.38441) (0.22552,-0.37819) (0.23878,-0.37155) (0.2523,-0.36454) (0.26606,-0.35719) (0.28006,-0.34955) (0.29427,-0.34165) (0.30868,-0.33352) (0.32328,-0.32519) (0.33806,-0.3167) (0.353,-0.30807) (0.36808,-0.29932) (0.38329,-0.29048) (0.39862,-0.28157) (0.41405,-0.27261) (0.42956,-0.26363) (0.44514,-0.25463) (0.46078,-0.24563) (0.47645,-0.23665) (0.49215,-0.2277) (0.50785,-0.21878) (0.52355,-0.2099) (0.53922,-0.20107) (0.55486,-0.19229) (0.57044,-0.18356) (0.58595,-0.17489) (0.60138,-0.16627) (0.61671,-0.15769) (0.63192,-0.14916) (0.647,-0.14067) (0.66194,-0.13221) (0.67672,-0.12377) (0.69132,-0.11535) (0.70573,-0.10693) (0.71994,-0.098492) (0.73394,-0.090035) (0.7477,-0.081542) (0.76122,-0.072998) (0.77448,-0.064388) (0.78747,-0.055696) (0.80017,-0.046906) (0.81258,-0.037999) (0.82468,-0.028959) (0.83646,-0.019766) (0.84791,-0.010401) (0.85902,-0.00084368) (0.86977,0.0089278) (0.88016,0.018936) (0.89017,0.029206) (0.89979,0.039764) (0.90902,0.050638) (0.91785,0.061862) (0.92627,0.073471) (0.93426,0.085506) (0.94183,0.098016) (0.94896,0.11106) (0.95565,0.12469) (0.96188,0.13901) (0.96766,0.15411) (0.97298,0.17012) (0.97784,0.18719) (0.98222,0.20556) (0.98612,0.2255) (0.98955,0.24743) (0.99249,0.27195) (0.99495,0.29995) (0.99692,0.33298) (0.9984,0.37378) (0.99938,0.42821) (0.99988,0.50906)};
\addplot[corange,very thick] coordinates {(0.00012336,0.98707) (0.00061668,0.88936) (0.0016028,0.72336) (0.0030808,0.53089) (0.0050492,0.34529) (0.0075061,0.18392) (0.010449,0.051559) (0.013875,-0.05372) (0.017781,-0.13625) (0.022162,-0.2006) (0.027016,-0.25076) (0.032336,-0.28994) (0.038117,-0.3206) (0.044355,-0.34463) (0.051042,-0.36343) (0.058172,-0.37807) (0.065738,-0.38937) (0.073733,-0.39793) (0.082148,-0.40422) (0.090976,-0.40862) (0.10021,-0.4114) (0.10983,-0.4128) (0.11984,-0.41299) (0.13023,-0.41212) (0.14098,-0.4103) (0.15209,-0.40765) (0.16354,-0.40424) (0.17532,-0.40015) (0.18742,-0.39545) (0.19983,-0.39018) (0.21253,-0.38441) (0.22552,-0.37819) (0.23878,-0.37155) (0.2523,-0.36454) (0.26606,-0.35719) (0.28006,-0.34955) (0.29427,-0.34165) (0.30868,-0.33352) (0.32328,-0.32519) (0.33806,-0.3167) (0.353,-0.30807) (0.36808,-0.29932) (0.38329,-0.29048) (0.39862,-0.28157) (0.41405,-0.27261) (0.42956,-0.26363) (0.44514,-0.25463) (0.46078,-0.24563) (0.47645,-0.23665) (0.49215,-0.2277) (0.50785,-0.21878) (0.52355,-0.2099) (0.53922,-0.20107) (0.55486,-0.19229) (0.57044,-0.18356) (0.58595,-0.17489) (0.60138,-0.16627) (0.61671,-0.15769) (0.63192,-0.14916) (0.647,-0.14067) (0.66194,-0.13221) (0.67672,-0.12377) (0.69132,-0.11535) (0.70573,-0.10693) (0.71994,-0.098492) (0.73394,-0.090035) (0.7477,-0.081542) (0.76122,-0.072998) (0.77448,-0.064388) (0.78747,-0.055696) (0.80017,-0.046906) (0.81258,-0.037999) (0.82468,-0.028959) (0.83646,-0.019766) (0.84791,-0.010401) (0.85902,-0.00084368) (0.86977,0.0089278) (0.88016,0.018936) (0.89017,0.029206) (0.89979,0.039764) (0.90902,0.050638) (0.91785,0.061862) (0.92627,0.073471) (0.93426,0.085506) (0.94183,0.098016) (0.94896,0.11106) (0.95565,0.12469) (0.96188,0.13901) (0.96766,0.15411) (0.97298,0.17012) (0.97784,0.18719) (0.98222,0.20556) (0.98612,0.2255) (0.98955,0.24743) (0.99249,0.27195) (0.99495,0.29995) (0.99692,0.33298) (0.9984,0.37378) (0.99938,0.42821) (0.99988,0.50906)};
\node[font=\scriptsize,anchor=south east,fill=white,inner sep=1pt] at (rel axis cs:0.98,0.02) {$C_L=0.000$};
\nextgroupplot[title={NACA 0012, $\alpha=5^\circ$}]
\addplot[cblue,very thick] coordinates {(0.00012336,-0.64686) (0.00061668,-1.1651) (0.0016028,-1.5826) (0.0030808,-1.8604) (0.0050492,-2.0079) (0.0075061,-2.058) (0.010449,-2.0448) (0.013875,-1.9949) (0.017781,-1.9261) (0.022162,-1.8492) (0.027016,-1.7707) (0.032336,-1.694) (0.038117,-1.6208) (0.044355,-1.5519) (0.051042,-1.4874) (0.058172,-1.4272) (0.065738,-1.3712) (0.073733,-1.3188) (0.082148,-1.2698) (0.090976,-1.2239) (0.10021,-1.1806) (0.10983,-1.1397) (0.11984,-1.101) (0.13023,-1.0642) (0.14098,-1.0291) (0.15209,-0.99554) (0.16354,-0.9633) (0.17532,-0.93229) (0.18742,-0.90238) (0.19983,-0.87348) (0.21253,-0.84548) (0.22552,-0.81832) (0.23878,-0.79194) (0.2523,-0.76627) (0.26606,-0.74128) (0.28006,-0.71692) (0.29427,-0.69316) (0.30868,-0.66997) (0.32328,-0.64733) (0.33806,-0.62521) (0.353,-0.60361) (0.36808,-0.5825) (0.38329,-0.56186) (0.39862,-0.54169) (0.41405,-0.52198) (0.42956,-0.50271) (0.44514,-0.48387) (0.46078,-0.46545) (0.47645,-0.44744) (0.49215,-0.42982) (0.50785,-0.41259) (0.52355,-0.39574) (0.53922,-0.37924) (0.55486,-0.36309) (0.57044,-0.34726) (0.58595,-0.33176) (0.60138,-0.31655) (0.61671,-0.30162) (0.63192,-0.28696) (0.647,-0.27254) (0.66194,-0.25836) (0.67672,-0.24438) (0.69132,-0.23058) (0.70573,-0.21696) (0.71994,-0.20348) (0.73394,-0.19013) (0.7477,-0.17688) (0.76122,-0.16372) (0.77448,-0.15062) (0.78747,-0.13756) (0.80017,-0.12451) (0.81258,-0.11146) (0.82468,-0.098382) (0.83646,-0.085252) (0.84791,-0.072049) (0.85902,-0.058747) (0.86977,-0.045323) (0.88016,-0.031751) (0.89017,-0.018004) (0.89979,-0.0040549) (0.90902,0.010128) (0.91785,0.024578) (0.92627,0.039331) (0.93426,0.05443) (0.94183,0.069922) (0.94896,0.085863) (0.95565,0.10232) (0.96188,0.11937) (0.96766,0.13712) (0.97298,0.15569) (0.97784,0.17524) (0.98222,0.19598) (0.98612,0.21819) (0.98955,0.24228) (0.99249,0.26883) (0.99495,0.29873) (0.99692,0.33345) (0.9984,0.37572) (0.99938,0.4313) (0.99988,0.51279)};
\addplot[corange,very thick] coordinates {(0.00012336,-0.11677) (0.00061668,0.34601) (0.0016028,0.68737) (0.0030808,0.8933) (0.0050492,0.98506) (0.0075061,0.99738) (0.010449,0.96181) (0.013875,0.90101) (0.017781,0.82926) (0.022162,0.75485) (0.027016,0.68226) (0.032336,0.61371) (0.038117,0.55015) (0.044355,0.49181) (0.051042,0.43857) (0.058172,0.39015) (0.065738,0.34616) (0.073733,0.30624) (0.082148,0.27002) (0.090976,0.23715) (0.10021,0.20732) (0.10983,0.18026) (0.11984,0.15573) (0.13023,0.1335) (0.14098,0.11337) (0.15209,0.095174) (0.16354,0.078756) (0.17532,0.063974) (0.18742,0.050701) (0.19983,0.038822) (0.21253,0.028232) (0.22552,0.018834) (0.23878,0.010539) (0.2523,0.0032645) (0.26606,-0.0030652) (0.28006,-0.0085219) (0.29427,-0.013172) (0.30868,-0.017079) (0.32328,-0.020299) (0.33806,-0.022889) (0.353,-0.024899) (0.36808,-0.026378) (0.38329,-0.027369) (0.39862,-0.027916) (0.41405,-0.028056) (0.42956,-0.027826) (0.44514,-0.02726) (0.46078,-0.026387) (0.47645,-0.025235) (0.49215,-0.02383) (0.50785,-0.022195) (0.52355,-0.020349) (0.53922,-0.018311) (0.55486,-0.016096) (0.57044,-0.013718) (0.58595,-0.011188) (0.60138,-0.008514) (0.61671,-0.0057039) (0.63192,-0.0027626) (0.647,0.00030676) (0.66194,0.0035028) (0.67672,0.0068259) (0.69132,0.010278) (0.70573,0.013862) (0.71994,0.017583) (0.73394,0.021448) (0.7477,0.025463) (0.76122,0.029637) (0.77448,0.033981) (0.78747,0.038505) (0.80017,0.04322) (0.81258,0.048141) (0.82468,0.053281) (0.83646,0.058655) (0.84791,0.06428) (0.85902,0.070173) (0.86977,0.076354) (0.88016,0.082842) (0.89017,0.089661) (0.89979,0.096835) (0.90902,0.10439) (0.91785,0.11236) (0.92627,0.12078) (0.93426,0.1297) (0.94183,0.13914) (0.94896,0.14919) (0.95565,0.15989) (0.96188,0.17134) (0.96766,0.18363) (0.97298,0.19689) (0.97784,0.2113) (0.98222,0.22706) (0.98612,0.24447) (0.98955,0.26394) (0.99249,0.28607) (0.99495,0.31179) (0.99692,0.34262) (0.9984,0.38134) (0.99938,0.4338) (0.99988,0.51279)};
\node[font=\scriptsize,anchor=south east,fill=white,inner sep=1pt] at (rel axis cs:0.98,0.02) {$C_L=0.603$};
\nextgroupplot[title={NACA 2412, $\alpha=2^\circ$}]
\addplot[cblue,very thick] coordinates {(-1.4838e-05,0.77064) (0.00020448,0.52203) (0.00092198,0.24185) (0.0021378,-0.021587) (0.0038514,-0.24115) (0.0060621,-0.41069) (0.0087685,-0.53586) (0.011969,-0.62605) (0.01566,-0.69031) (0.01984,-0.73587) (0.024505,-0.76813) (0.029651,-0.79094) (0.035273,-0.807) (0.041368,-0.81819) (0.047928,-0.82581) (0.054949,-0.83077) (0.062423,-0.8337) (0.070344,-0.83503) (0.078703,-0.83507) (0.087494,-0.83403) (0.096707,-0.83205) (0.10633,-0.82924) (0.11636,-0.82565) (0.12679,-0.82133) (0.13759,-0.81631) (0.14877,-0.81062) (0.16031,-0.80426) (0.17219,-0.79724) (0.18442,-0.78957) (0.19696,-0.78125) (0.20982,-0.77229) (0.22297,-0.76269) (0.2364,-0.75246) (0.25011,-0.74159) (0.26407,-0.73009) (0.27826,-0.71795) (0.29268,-0.70518) (0.30731,-0.69175) (0.32214,-0.67766) (0.33714,-0.66285) (0.3523,-0.64726) (0.36761,-0.63075) (0.38305,-0.61298) (0.39857,-0.59301) (0.41414,-0.57296) (0.42974,-0.55449) (0.44542,-0.53715) (0.46115,-0.52047) (0.47691,-0.50423) (0.49269,-0.48834) (0.50848,-0.47272) (0.52425,-0.45732) (0.54,-0.44213) (0.5557,-0.4271) (0.57135,-0.41223) (0.58692,-0.39749) (0.6024,-0.38287) (0.61777,-0.36835) (0.63302,-0.35393) (0.64814,-0.33957) (0.6631,-0.32528) (0.6779,-0.31104) (0.69252,-0.29683) (0.70695,-0.28263) (0.72116,-0.26844) (0.73515,-0.25424) (0.74891,-0.24002) (0.76242,-0.22575) (0.77566,-0.21142) (0.78863,-0.19702) (0.80131,-0.18253) (0.81369,-0.16793) (0.82576,-0.15322) (0.8375,-0.13837) (0.84891,-0.12336) (0.85997,-0.10818) (0.87068,-0.092813) (0.88102,-0.077237) (0.89098,-0.061432) (0.90055,-0.045378) (0.90973,-0.029051) (0.9185,-0.012424) (0.92686,0.0045331) (0.9348,0.021855) (0.94232,0.039582) (0.9494,0.057764) (0.95603,0.07646) (0.96222,0.095743) (0.96795,0.1157) (0.97323,0.13646) (0.97804,0.15817) (0.98238,0.18102) (0.98625,0.20531) (0.98965,0.2314) (0.99257,0.25988) (0.995,0.29161) (0.99695,0.32805) (0.99841,0.37189) (0.99939,0.42886) (0.99988,0.51151)};
\addplot[corange,very thick] coordinates {(0.00026156,0.9374) (0.0010289,0.9995) (0.0022837,0.96529) (0.0040239,0.86694) (0.006247,0.74009) (0.0089501,0.6104) (0.01213,0.49149) (0.015782,0.3885) (0.019902,0.30194) (0.024485,0.23031) (0.029527,0.17147) (0.035021,0.1233) (0.040961,0.083896) (0.047342,0.051689) (0.054155,0.025391) (0.061395,0.0039608) (0.069053,-0.013434) (0.077121,-0.027462) (0.085592,-0.038666) (0.094457,-0.047483) (0.10371,-0.054271) (0.11333,-0.059326) (0.12333,-0.062893) (0.13367,-0.065179) (0.14437,-0.066358) (0.1554,-0.066581) (0.16676,-0.065978) (0.17844,-0.064664) (0.19042,-0.062738) (0.20269,-0.060293) (0.21525,-0.05741) (0.22808,-0.054165) (0.24116,-0.050628) (0.25449,-0.046865) (0.26806,-0.042938) (0.28185,-0.038909) (0.29585,-0.034835) (0.31005,-0.03078) (0.32443,-0.026808) (0.33899,-0.022992) (0.3537,-0.019422) (0.36855,-0.016223) (0.38354,-0.013619) (0.39866,-0.012141) (0.41396,-0.011327) (0.42937,-0.009486) (0.44486,-0.0067567) (0.4604,-0.0035765) (0.47599,-0.00011111) (0.4916,0.0035457) (0.50723,0.0073323) (0.52285,0.011205) (0.53845,0.015132) (0.55402,0.01909) (0.56954,0.023059) (0.58499,0.027027) (0.60037,0.030985) (0.61565,0.034925) (0.63082,0.038845) (0.64587,0.042742) (0.66078,0.04662) (0.67553,0.05048) (0.69012,0.054327) (0.70452,0.058169) (0.71872,0.062013) (0.73272,0.065869) (0.74649,0.069748) (0.76001,0.073663) (0.77329,0.077626) (0.7863,0.081654) (0.79903,0.08576) (0.81147,0.089963) (0.82361,0.09428) (0.83542,0.09873) (0.84691,0.10333) (0.85806,0.10812) (0.86886,0.11309) (0.8793,0.1183) (0.88936,0.12375) (0.89904,0.12949) (0.90832,0.13554) (0.9172,0.14194) (0.92567,0.14873) (0.93372,0.15596) (0.94134,0.16368) (0.94852,0.17195) (0.95526,0.18085) (0.96155,0.19046) (0.96737,0.2009) (0.97274,0.2123) (0.97763,0.22482) (0.98205,0.23871) (0.986,0.25425) (0.98945,0.27189) (0.99242,0.29222) (0.9949,0.3162) (0.99689,0.34537) (0.99838,0.38256) (0.99938,0.43365) (0.99988,0.51151)};
\node[font=\scriptsize,anchor=south east,fill=white,inner sep=1pt] at (rel axis cs:0.98,0.02) {$C_L=0.498$};
\end{groupplot}
\end{tikzpicture}

\caption{Distributions de $C_p$ calculées par \code{panel.solve} (200 panneaux). Bleu : extrados ; orange : intrados. Axe vertical inversé, convention usuelle. Le pic d'aspiration du NACA 0012 passe de $C_p=-0{,}41$ à $\alpha=0$ à $-2{,}06$ à $\alpha=5^\circ$.}\label{fig:cp}
\end{figure}

\section{Vérifications}
\paragraph{Convergence.} La figure~\ref{fig:convergence} montre, pour le NACA 0012 à $5^\circ$, que les trois indicateurs décroissent à peu près en $1/N$ (pentes mesurées $-1{,}17$ pour $C_L$ et $-1{,}24$ pour $C_{D,p}$). C'est l'ordre attendu pour des panneaux plans à singularités constantes avec collocation au milieu. À 200 panneaux, $C_L$ est à $0{,}06$~\% de sa valeur limite.

\begin{figure}[htbp]
\centering
% figure générée par tools/make_figures.py : convergence
\begin{tikzpicture}
\begin{loglogaxis}[width=0.72\linewidth,height=6cm,xlabel={nombre de panneaux $N$},ylabel={erreur},
  legend style={draw=none,font=\small,at={(0.02,0.02)},anchor=south west},xmin=15,xmax=800,xtick={20,50,100,200,500},xticklabels={20,50,100,200,500}]
\addplot[cblue,very thick,mark=*,mark size=1.5pt] coordinates {(20,0.001539) (30,0.003499) (40,0.003447) (60,0.002396) (80,0.001768) (120,0.001137) (160,0.0008237) (200,0.000638) (320,0.0003657) (480,0.0002197) (640,0.0001488)};
\addlegendentry{$|C_L(N)-C_L^\infty|/C_L^\infty$}
\addplot[corange,very thick,mark=square*,mark size=1.5pt] coordinates {(20,0.01875) (30,0.01904) (40,0.01595) (60,0.01136) (80,0.008758) (120,0.006) (160,0.004562) (200,0.00368) (320,0.002328) (480,0.001562) (640,0.001176)};
\addlegendentry{$|C_{L,C_p}-C_{L,KJ}|/C_L^\infty$}
\addplot[cgreen,very thick,mark=triangle*,mark size=1.8pt] coordinates {(20,0.005424) (30,0.002714) (40,0.001664) (60,0.0008615) (80,0.0005582) (120,0.0003175) (160,0.0002189) (200,0.0001662) (320,9.574e-05) (480,6.09e-05) (640,4.459e-05)};
\addlegendentry{$|C_{D,p}|$ (d'Alembert)}
\addplot[cmuted,dashed,domain=40:640] {0.6/x};
\node[cmuted,font=\scriptsize] at (axis cs:330,0.0035) {pente $-1$};
\end{loglogaxis}
\end{tikzpicture}

\caption{Convergence en maillage sur le NACA 0012 à $\alpha=5^\circ$. La référence $C_L^\infty$ est calculée avec 2\,400 panneaux.}\label{fig:convergence}
\end{figure}

\paragraph{Pente et angle de portance nulle.} La figure~\ref{fig:cl-alpha} compare les panneaux à la théorie des profils minces. Les pentes mesurées valent $6{,}88$ à $6{,}89$~/rad, contre $2\pi(1+0{,}77\times0{,}12)=6{,}86$ attendu avec la correction d'épaisseur. Les angles de portance nulle valent $-2{,}13^\circ$ (2412) et $-4{,}26^\circ$ (4412), contre $-2{,}08^\circ$ et $-4{,}15^\circ$ en profil mince.

\begin{figure}[htbp]
\centering
% figure générée par tools/make_figures.py : cl_alpha
\begin{tikzpicture}
\begin{axis}[width=0.75\linewidth,height=6.4cm,xlabel={$\alpha$ (°)},ylabel={$C_L$ (non visqueux)},
  xmin=-4,xmax=10,legend columns=2,legend style={draw=none,font=\scriptsize,at={(0.5,-0.22)},anchor=north},legend cell align=left]
\addplot[cmuted,very thick,mark=*,mark size=1.3pt] coordinates {(-4,-0.4823) (-3,-0.36185) (-2,-0.2413) (-1,-0.12067) (0,5.5511e-17) (1,0.12067) (2,0.2413) (3,0.36185) (4,0.4823) (5,0.6026) (6,0.72271) (7,0.84261) (8,0.96224) (9,1.0816) (10,1.2006)};\addlegendentry{NACA 0012 (panneaux)}
\addplot[cmuted,dashed,domain=-4:10] {2*pi*(x-(-0.0000))*pi/180};\addlegendentry{profil mince, $\alpha_{L0}=-0.00^\circ$}
\addplot[cblue,very thick,mark=*,mark size=1.3pt] coordinates {(-4,-0.2265) (-3,-0.10581) (-2,0.014921) (-1,0.13565) (0,0.25635) (1,0.37698) (2,0.4975) (3,0.61788) (4,0.73808) (5,0.85806) (6,0.97778) (7,1.0972) (8,1.2163) (9,1.3351) (10,1.4534)};\addlegendentry{NACA 2412 (panneaux)}
\addplot[cblue,dashed,domain=-4:10] {2*pi*(x-(-2.0772))*pi/180};\addlegendentry{profil mince, $\alpha_{L0}=-2.08^\circ$}
\addplot[corange,very thick,mark=*,mark size=1.3pt] coordinates {(-4,0.028327) (-3,0.14925) (-2,0.27014) (-1,0.39097) (0,0.51169) (1,0.63227) (2,0.75268) (3,0.87287) (4,0.99282) (5,1.1125) (6,1.2318) (7,1.3508) (8,1.4694) (9,1.5875) (10,1.7052)};\addlegendentry{NACA 4412 (panneaux)}
\addplot[corange,dashed,domain=-4:10] {2*pi*(x-(-4.1545))*pi/180};\addlegendentry{profil mince, $\alpha_{L0}=-4.15^\circ$}
\end{axis}
\end{tikzpicture}

\caption{$C_L(\alpha)$ non visqueux de trois profils (traits pleins) et théorie des profils minces (tirets), avec l'angle $\alpha_{L0}$ calculé par l'intégrale du chapitre~\ref{ch:potentiel}. L'écart de pente vient de l'épaisseur.}\label{fig:cl-alpha}
\end{figure}

\begin{resultat}
\begin{itemize}
\item $N$ sources constantes, un tourbillon uniforme, $N$ conditions de glissement et une condition de Kutta : un système dense $(N+1)\times(N+1)$, inversé une fois par géométrie.
\item Sorties : $V_t$ signée, $C_p$, $C_L$, $C_M$, et deux contrôles ($C_{L,KJ}$, $C_{D,p}$).
\item Précision : $1/N$, soit moins de 0,1~\% sur $C_L$ à 200 panneaux.
\end{itemize}
\end{resultat}

\begin{projet}
Côté optimisation, deux propriétés de ce bloc comptent. D'abord, il est \emph{exactement} linéaire en $\alpha$ pour une géométrie donnée, ce qui rend $C_L$ quasi affine en $(m,\alpha)$ : un modèle quadratique représente $C_L$ avec $R^2=0{,}9998$. Ensuite, c'est une méthode d'éléments de frontière, comme la méthode des moments qu'utilisent beaucoup de solveurs de SER : on discrétise la surface et on résout un système dense. C'est aussi le type de système que visent les algorithmes quantiques de résolution linéaire cités dans la note d'intention.
\end{projet}

\exercices
\begin{exercice}
Vérifier le calcul de $v_\sigma$ : montrer que $\int_0^L\frac{\eta\,\dd\xi'}{(\xi-\xi')^2+\eta^2}=\theta_2-\theta_1$, avec $\theta_k=\operatorname{atan2}(\eta,\xi-\xi_k)$. Que vaut cette intégrale très loin du panneau ? Interpréter.
\end{exercice}
\begin{exercice}
(a) Quel est le coût, en ordre de grandeur, de l'inversion d'une matrice $201\times201$, puis d'une résolution pour une nouvelle incidence une fois l'inverse connue ? (b) En déduire pourquoi \code{alpha\_for\_cl}, qui appelle \code{solve} 5 à 8 fois, ne coûte presque rien de plus qu'un seul calcul.
\end{exercice}

\chapter{La couche limite intégrale : principe et mise en place}\label{ch:cl}
\iterbanner{V0}

\section{Pourquoi une méthode intégrale}
Les équations de Prandtl sont des EDP paraboliques en $(s,y)$. Les résoudre sur chaque face demanderait un maillage dans l'épaisseur, une loi de paroi ou un modèle de turbulence, et des centaines de stations. Les \textbf{méthodes intégrales} font un pari plus modeste : intégrer les équations dans l'épaisseur, ne suivre que quelques grandeurs globales du profil de vitesse, et fermer le système par des corrélations empiriques. On passe d'une EDP à une ou deux \textbf{EDO en $s$}, résolues en quelques millisecondes.

\section{Rappels : grandeurs intégrales et équation de von Kármán}
Pour un profil de vitesse $u(s,y)$ qui tend vers $U_e(s)$ en bord de couche :
\[
\delta^*=\int_0^\infty\Big(1-\frac{u}{U_e}\Big)\dd y,\qquad
\theta=\int_0^\infty\frac{u}{U_e}\Big(1-\frac{u}{U_e}\Big)\dd y,\qquad
H=\frac{\delta^*}{\theta},\qquad c_f=\frac{\tau_w}{\tfrac12\rho U_e^2}.
\]
$\delta^*$ est le déplacement des lignes de courant extérieures, $\theta$ mesure le déficit de quantité de mouvement, et $H$ est un indicateur de la \emph{forme} du profil : environ 2,6 pour un profil laminaire de Blasius, 1,3 à 1,4 pour un profil turbulent plein, et au-delà de 3,5 (laminaire) ou de 2 à 2,8 (turbulent) près du décollement.

L'intégration de l'équation de quantité de mouvement de Prandtl dans l'épaisseur donne l'\textbf{équation intégrale de von Kármán}, valable en laminaire comme en turbulent :
\[
\boxed{\ \frac{\dd\theta}{\dd s}+\big(2+H\big)\frac{\theta}{U_e}\frac{\dd U_e}{\dd s}=\frac{c_f}{2}\ }
\]
Elle se lit comme un bilan : le déficit de quantité de mouvement croît par frottement ($c_f/2$), et un gradient de pression adverse ($U_e'<0$) l'amplifie d'autant plus que le profil est creusé ($H$ grand). Une équation, trois inconnues ($\theta$, $H$, $c_f$) : toutes les méthodes intégrales sont des façons de fermer ce système.

\begin{center}
\begin{tabularx}{\linewidth}{lllX}
\toprule
Régime & Méthode & Inconnues & Fermeture \\
\midrule
laminaire & Thwaites (ch.~\ref{ch:thwaites}) & $\theta$ & $H$ et $c_f$ fonctions de $\lambda=\theta^2U_e'/\nu$ \\
transition & Michel (ch.~\ref{ch:transition}) & --- & seuil sur $\Re_\theta$ \\
turbulent & Head (ch.~\ref{ch:head}) & $\theta$, $H$ & équation d'entraînement, $H_1(H)$, $c_f$ de Ludwieg-Tillmann \\
\bottomrule
\end{tabularx}
\end{center}

\section{De la solution de panneaux aux deux faces}
La couche limite démarre au \textbf{point d'arrêt} et se développe séparément sur l'extrados et sur l'intrados. Il faut donc trouver ce point, puis réorganiser les données de chaque face dans le sens de l'écoulement.

\paragraph{Le point d'arrêt.} En parcourant le contour depuis le bord de fuite extrados, la vitesse tangentielle signée $V_t$ est d'abord négative (l'écoulement remonte le sens de parcours), s'annule au point d'arrêt, puis devient positive sur l'intrados. On cherche l'unique indice $k$ tel que $V_{t,k}<0\le V_{t,k+1}$, et on interpole linéairement le zéro :
\[
s_{\text{arrêt}}=\ell_k+w\,(\ell_{k+1}-\ell_k),\qquad w=\frac{V_{t,k}}{V_{t,k}-V_{t,k+1}},
\]
où $\ell_k$ est l'abscisse curviligne cumulée du point de collocation $k$ le long du contour.

\paragraph{Les deux faces.} Chaque face reçoit ensuite trois tableaux, orientés du point d'arrêt vers le bord de fuite :
\begin{itemize}
\item l'abscisse curviligne $s$ mesurée depuis le point d'arrêt (donc $s>0$) ;
\item l'abscisse $x/c$ des stations, pour repérer la transition et la coupure en $x_{stop}$ ;
\item la vitesse extérieure $U_e=|V_t|$ : $-V_t$ sur l'extrados, $+V_t$ sur l'intrados.
\end{itemize}

\begin{figure}[htbp]
\centering
% figure générée par tools/make_figures.py : vt_contour
\begin{tikzpicture}
\begin{groupplot}[group style={group size=2 by 1,horizontal sep=1.6cm},width=0.5\linewidth,height=5.4cm,
  title style={font=\small},tick label style={font=\scriptsize},legend style={draw=none,font=\scriptsize}]
\nextgroupplot[title={$V_t$ signée le long du contour},xlabel={abscisse $\ell$ depuis le BF extrados},ylabel={$V_t/U_\infty$},xmin=0,xmax=2.042]
\addplot[cink,very thick] coordinates {(0.00012493,-0.69892) (0.00062451,-0.75574) (0.0016232,-0.79253) (0.0031199,-0.81972) (0.0051131,-0.84166) (0.0076009,-0.8603) (0.010581,-0.8767) (0.01405,-0.89146) (0.018004,-0.90497) (0.02244,-0.91751) (0.027352,-0.92927) (0.032737,-0.94037) (0.038589,-0.95092) (0.044901,-0.96101) (0.051668,-0.97069) (0.058882,-0.98001) (0.066537,-0.98901) (0.074624,-0.99773) (0.083136,-1.0062) (0.092064,-1.0144) (0.1014,-1.0224) (0.11113,-1.0303) (0.12125,-1.0379) (0.13175,-1.0454) (0.14262,-1.0527) (0.15385,-1.0599) (0.16542,-1.0669) (0.17732,-1.0739) (0.18955,-1.0807) (0.20208,-1.0874) (0.21492,-1.0941) (0.22804,-1.1006) (0.24143,-1.1071) (0.25508,-1.1136) (0.26898,-1.1199) (0.28311,-1.1263) (0.29746,-1.1325) (0.312,-1.1388) (0.32674,-1.145) (0.34166,-1.1512) (0.35673,-1.1574) (0.37195,-1.1636) (0.3873,-1.1698) (0.40276,-1.176) (0.41832,-1.1822) (0.43396,-1.1884) (0.44968,-1.1946) (0.46544,-1.2009) (0.48125,-1.2072) (0.49707,-1.2136) (0.5129,-1.22) (0.52872,-1.2265) (0.54452,-1.2331) (0.56027,-1.2398) (0.57597,-1.2468) (0.59159,-1.2542) (0.60717,-1.2621) (0.6227,-1.27) (0.63814,-1.277) (0.65345,-1.2835) (0.66862,-1.2895) (0.68362,-1.2952) (0.69844,-1.3007) (0.71308,-1.3058) (0.72751,-1.3107) (0.74172,-1.3153) (0.7557,-1.3197) (0.76943,-1.3238) (0.78291,-1.3277) (0.79611,-1.3313) (0.80903,-1.3346) (0.82166,-1.3377) (0.83398,-1.3406) (0.84598,-1.3432) (0.85766,-1.3456) (0.86899,-1.3477) (0.87998,-1.3496) (0.89061,-1.3512) (0.90087,-1.3525) (0.91076,-1.3535) (0.92026,-1.3543) (0.92938,-1.3546) (0.9381,-1.3546) (0.94642,-1.3541) (0.95433,-1.3531) (0.96183,-1.3512) (0.96892,-1.3484) (0.97559,-1.3442) (0.98185,-1.3383) (0.9877,-1.3297) (0.99314,-1.3175) (0.99818,-1.3001) (1.0028,-1.2752) (1.0071,-1.2393) (1.011,-1.1877) (1.0146,-1.1141) (1.0179,-1.0107) (1.021,-0.87072) (1.0239,-0.69135) (1.0267,-0.47892) (1.0295,-0.25021) (1.0323,-0.022469) (1.0352,0.18632) (1.0383,0.36477) (1.0416,0.50981) (1.0452,0.62418) (1.0491,0.7131) (1.0533,0.78198) (1.058,0.8355) (1.0629,0.87732) (1.0683,0.91023) (1.0741,0.93632) (1.0803,0.95713) (1.0869,0.97381) (1.0938,0.98722) (1.1012,0.99802) (1.109,1.0067) (1.1172,1.0136) (1.1257,1.0191) (1.1346,1.0235) (1.1439,1.0268) (1.1536,1.0292) (1.1636,1.031) (1.174,1.0321) (1.1847,1.0326) (1.1957,1.0328) (1.2071,1.0325) (1.2188,1.0318) (1.2308,1.0309) (1.243,1.0297) (1.2556,1.0283) (1.2684,1.0267) (1.2815,1.025) (1.2948,1.0232) (1.3084,1.0212) (1.3222,1.0193) (1.3362,1.0173) (1.3504,1.0153) (1.3648,1.0133) (1.3793,1.0114) (1.3941,1.0097) (1.4089,1.0081) (1.4239,1.0068) (1.4391,1.0061) (1.4544,1.0056) (1.4698,1.0047) (1.4853,1.0034) (1.5009,1.0018) (1.5165,1.0001) (1.5321,0.99823) (1.5477,0.99633) (1.5634,0.99438) (1.579,0.99241) (1.5946,0.99041) (1.6102,0.9884) (1.6256,0.98639) (1.641,0.98439) (1.6564,0.98238) (1.6716,0.98039) (1.6866,0.9784) (1.7016,0.97641) (1.7164,0.97443) (1.731,0.97246) (1.7454,0.97048) (1.7596,0.9685) (1.7737,0.9665) (1.7875,0.9645) (1.801,0.96246) (1.8143,0.9604) (1.8274,0.9583) (1.8401,0.95616) (1.8526,0.95396) (1.8648,0.95169) (1.8766,0.94935) (1.8881,0.94692) (1.8993,0.9444) (1.9101,0.94176) (1.9206,0.93899) (1.9307,0.93608) (1.9404,0.93301) (1.9497,0.92976) (1.9586,0.92631) (1.9671,0.92264) (1.9752,0.91871) (1.9828,0.9145) (1.99,0.90997) (1.9968,0.90507) (2.0031,0.89974) (2.0089,0.89392) (2.0143,0.88753) (2.0192,0.88044) (2.0236,0.87252) (2.0276,0.86357) (2.0311,0.8533) (2.034,0.8413) (2.0365,0.82692) (2.0385,0.80909) (2.04,0.78577) (2.041,0.75256) (2.0415,0.69892)};
\addplot[cmuted,thin] coordinates {(0,0) (2.042,0)};
\addplot[only marks,mark=*,corange,mark size=2.4pt] coordinates {(1.0326,0)};
\node[corange,font=\scriptsize,anchor=north west] at (axis cs:1.063,-0.08) {point d'arrêt};
\node[font=\scriptsize,anchor=north] at (axis cs:0.5,-0.2) {extrados ($V_t<0$)};
\node[font=\scriptsize,anchor=north] at (axis cs:1.6,0.8) {intrados ($V_t>0$)};
\nextgroupplot[title={$U_e=|V_t|$ sur chaque face},xlabel={$x/c$},ylabel={$U_e/U_\infty$},xmin=0,xmax=1,legend pos=south east]
\addplot[cblue,very thick] coordinates {(0.0010289,0.022469) (0.00026156,0.25021) (-1.4838e-05,0.47892) (0.00020448,0.69135) (0.00092198,0.87072) (0.0021378,1.0107) (0.0038514,1.1141) (0.0060621,1.1877) (0.0087685,1.2393) (0.011969,1.2752) (0.01566,1.3001) (0.01984,1.3175) (0.024505,1.3297) (0.029651,1.3383) (0.035273,1.3442) (0.041368,1.3484) (0.047928,1.3512) (0.054949,1.3531) (0.062423,1.3541) (0.070344,1.3546) (0.078703,1.3546) (0.087494,1.3543) (0.096707,1.3535) (0.10633,1.3525) (0.11636,1.3512) (0.12679,1.3496) (0.13759,1.3477) (0.14877,1.3456) (0.16031,1.3432) (0.17219,1.3406) (0.18442,1.3377) (0.19696,1.3346) (0.20982,1.3313) (0.22297,1.3277) (0.2364,1.3238) (0.25011,1.3197) (0.26407,1.3153) (0.27826,1.3107) (0.29268,1.3058) (0.30731,1.3007) (0.32214,1.2952) (0.33714,1.2895) (0.3523,1.2835) (0.36761,1.277) (0.38305,1.27) (0.39857,1.2621) (0.41414,1.2542) (0.42974,1.2468) (0.44542,1.2398) (0.46115,1.2331) (0.47691,1.2265) (0.49269,1.22) (0.50848,1.2136) (0.52425,1.2072) (0.54,1.2009) (0.5557,1.1946) (0.57135,1.1884) (0.58692,1.1822) (0.6024,1.176) (0.61777,1.1698) (0.63302,1.1636) (0.64814,1.1574) (0.6631,1.1512) (0.6779,1.145) (0.69252,1.1388) (0.70695,1.1325) (0.72116,1.1263) (0.73515,1.1199) (0.74891,1.1136) (0.76242,1.1071) (0.77566,1.1006) (0.78863,1.0941) (0.80131,1.0874) (0.81369,1.0807) (0.82576,1.0739) (0.8375,1.0669) (0.84891,1.0599) (0.85997,1.0527) (0.87068,1.0454) (0.88102,1.0379) (0.89098,1.0303) (0.90055,1.0224) (0.90973,1.0144) (0.9185,1.0062) (0.92686,0.99773) (0.9348,0.98901) (0.94232,0.98001) (0.9494,0.97069) (0.95603,0.96101) (0.96222,0.95092) (0.96795,0.94037) (0.97323,0.92927) (0.97804,0.91751) (0.98238,0.90497) (0.98625,0.89146) (0.98965,0.8767) (0.99257,0.8603) (0.995,0.84166) (0.99695,0.81972) (0.99841,0.79253) (0.99939,0.75574) (0.99988,0.69892)};\addlegendentry{extrados}
\addplot[corange,very thick] coordinates {(0.0022837,0.18632) (0.0040239,0.36477) (0.006247,0.50981) (0.0089501,0.62418) (0.01213,0.7131) (0.015782,0.78198) (0.019902,0.8355) (0.024485,0.87732) (0.029527,0.91023) (0.035021,0.93632) (0.040961,0.95713) (0.047342,0.97381) (0.054155,0.98722) (0.061395,0.99802) (0.069053,1.0067) (0.077121,1.0136) (0.085592,1.0191) (0.094457,1.0235) (0.10371,1.0268) (0.11333,1.0292) (0.12333,1.031) (0.13367,1.0321) (0.14437,1.0326) (0.1554,1.0328) (0.16676,1.0325) (0.17844,1.0318) (0.19042,1.0309) (0.20269,1.0297) (0.21525,1.0283) (0.22808,1.0267) (0.24116,1.025) (0.25449,1.0232) (0.26806,1.0212) (0.28185,1.0193) (0.29585,1.0173) (0.31005,1.0153) (0.32443,1.0133) (0.33899,1.0114) (0.3537,1.0097) (0.36855,1.0081) (0.38354,1.0068) (0.39866,1.0061) (0.41396,1.0056) (0.42937,1.0047) (0.44486,1.0034) (0.4604,1.0018) (0.47599,1.0001) (0.4916,0.99823) (0.50723,0.99633) (0.52285,0.99438) (0.53845,0.99241) (0.55402,0.99041) (0.56954,0.9884) (0.58499,0.98639) (0.60037,0.98439) (0.61565,0.98238) (0.63082,0.98039) (0.64587,0.9784) (0.66078,0.97641) (0.67553,0.97443) (0.69012,0.97246) (0.70452,0.97048) (0.71872,0.9685) (0.73272,0.9665) (0.74649,0.9645) (0.76001,0.96246) (0.77329,0.9604) (0.7863,0.9583) (0.79903,0.95616) (0.81147,0.95396) (0.82361,0.95169) (0.83542,0.94935) (0.84691,0.94692) (0.85806,0.9444) (0.86886,0.94176) (0.8793,0.93899) (0.88936,0.93608) (0.89904,0.93301) (0.90832,0.92976) (0.9172,0.92631) (0.92567,0.92264) (0.93372,0.91871) (0.94134,0.9145) (0.94852,0.90997) (0.95526,0.90507) (0.96155,0.89974) (0.96737,0.89392) (0.97274,0.88753) (0.97763,0.88044) (0.98205,0.87252) (0.986,0.86357) (0.98945,0.8533) (0.99242,0.8413) (0.9949,0.82692) (0.99689,0.80909) (0.99838,0.78577) (0.99938,0.75256) (0.99988,0.69892)};\addlegendentry{intrados}
\end{groupplot}
\end{tikzpicture}

\caption{NACA 2412, $\alpha=2^\circ$. À gauche, la vitesse tangentielle signée le long du contour, parcouru depuis le bord de fuite extrados. Elle s'annule au point d'arrêt, situé ici en $x/c\approx0{,}001$ sur l'intrados. À droite, les vitesses extérieures que reçoit chaque face.}\label{fig:vt}
\end{figure}

\begin{codebox}[ : \code{boundary\_layer.split\_at\_stagnation}]
\begin{lstlisting}
k = np.nonzero((Vt[:-1] < 0) & (Vt[1:] >= 0))[0]    # changement de signe - -> +
if k.size != 1: raise RuntimeError("point d'arret ambigu")
s_nodes = np.concatenate([[0.0], np.cumsum(P.L)])
s_c = 0.5 * (s_nodes[:-1] + s_nodes[1:])          # abscisse des collocations
w = Vt[k] / (Vt[k] - Vt[k + 1])
s_stag = s_c[k] + w * (s_c[k + 1] - s_c[k])
up = np.arange(k, -1, -1)       # du point d'arret vers le BF extrados
lo = np.arange(k + 1, P.n)      # du point d'arret vers le BF intrados
upper = (s_stag - s_c[up], P.xc[up], -Vt[up], up)
lower = (s_c[lo] - s_stag, P.xc[lo],  Vt[lo], lo)
\end{lstlisting}
Si le signe change plusieurs fois (forme pathologique, très forte incidence), l'exception est rattrapée par \code{evaluate}, qui renvoie un résultat marqué invalide au lieu de planter.
\end{codebox}

\section{L'organisation de la marche sur une face}
La fonction \code{march\_surface(s, x, Ue, Re, ...)} traite une face en cinq temps :
\begin{enumerate}
\item couper la face en $x_{stop}=0{,}985$ (ch.~\ref{ch:trainee} explique pourquoi) ;
\item calculer la solution laminaire de Thwaites sur \emph{toute} la face, ce qui ne coûte qu'une somme cumulée ;
\item chercher la première station où un critère de transition est franchi, et interpoler sa position ;
\item intégrer les équations de Head de la transition jusqu'au bout ;
\item renvoyer un objet \code{SurfaceBL} avec $s$, $x$, $U_e$, $\theta$, $H$, $c_f$, un masque turbulent, $x_{tr}$, le mode de transition, $x_{sep}$ et des drapeaux.
\end{enumerate}
Les valeurs au bout de la face, $\theta_{TE}$, $H_{TE}$ et $U_{e,TE}$, alimentent Squire-Young.

\begin{piege}
$U_e$ est mis à $\max(U_e,10^{-6})$ avant tout calcul. La première station est très proche du point d'arrêt, où $U_e\to0$, et plusieurs formules divisent par $U_e$ ou par $U_e^6$. Le seuil évite une division par zéro sans changer le résultat.
\end{piege}

\exercices
\begin{exercice}
Calculer $\delta^*/\delta$, $\theta/\delta$ et $H$ pour le profil linéaire $u/U_e=y/\delta$ ($0\le y\le\delta$), puis pour le profil en puissance $u/U_e=(y/\delta)^{1/7}$. Lequel ressemble à un profil laminaire, lequel à un profil turbulent ?
\end{exercice}
\begin{exercice}
Pourquoi le point d'arrêt d'un profil cambré à incidence positive est-il sur l'intrados, légèrement en aval du bord d'attaque géométrique ? Dans quel sens se déplace-t-il quand $\alpha$ augmente ?
\end{exercice}

\chapter{Couche limite laminaire : la méthode de Thwaites}\label{ch:thwaites}
\iterbanner{V0}

\section{De von Kármán à une équation fermée}
On multiplie l'équation de von Kármán par $2U_e\theta/\nu$ :
\[
\frac{U_e}{\nu}\frac{\dd\theta^2}{\dd s}+2\big(2+H\big)\frac{\theta^2}{\nu}\frac{\dd U_e}{\dd s}=\frac{U_e\,\theta\,c_f}{\nu}.
\]
Deux groupements sans dimension apparaissent. Le premier est le \textbf{paramètre de gradient de pression}
\[
\lambda=\frac{\theta^2}{\nu}\frac{\dd U_e}{\dd s}.
\]
Le second est le frottement réduit $\ell=\tfrac12c_f\,\Re_\theta=\tfrac{\tau_w\theta}{\mu U_e}$. L'équation devient
\[
\frac{U_e}{\nu}\frac{\dd\theta^2}{\dd s}=2\big[\ell-(2+H)\lambda\big]\equiv F .
\]

\paragraph{L'hypothèse de Thwaites.} En laminaire, la forme du profil de vitesse est gouvernée principalement par le gradient de pression local. On suppose donc que $H$ et $\ell$ sont des fonctions du seul $\lambda$ : c'est une famille de profils à un paramètre. Thwaites a tracé $F(\lambda)$ pour toutes les solutions exactes connues (Falkner-Skan, Howarth, etc.) et constaté que les points s'alignent presque sur une droite (figure~\ref{fig:thwaites-f}) :
\[
F(\lambda)\approx a-b\lambda,\qquad a=0{,}45,\quad b=6 .
\]

\paragraph{La quadrature.} Avec cette droite, l'équation devient linéaire en $\theta^2$ :
\[
U_e\frac{\dd\theta^2}{\dd s}+6\,\theta^2\frac{\dd U_e}{\dd s}=0{,}45\,\nu
\quad\Longleftrightarrow\quad
\frac{1}{U_e^5}\frac{\dd}{\dd s}\big(\theta^2U_e^6\big)=0{,}45\,\nu .
\]
On a multiplié par $U_e^5$ pour faire apparaître une dérivée exacte. C'est l'origine de la puissance 6, et donc du 5 dans l'intégrale. En partant du point d'arrêt, où $U_e=0$ :
\[
\boxed{\ \theta^2(s)=\frac{0{,}45\,\nu}{U_e^6(s)}\int_0^sU_e^5(s')\,\dd s'\ }
\]
On n'a plus d'EDO à résoudre : $\theta$ s'obtient par une simple intégrale cumulée de $U_e^5$.

\section{Les corrélations de fermeture}
Une fois $\theta$ connu, on calcule $\lambda$, puis $H$ et $c_f$ par les corrélations de Thwaites, sous la forme de Cebeci et Bradshaw :
\[
0\le\lambda\le0{,}1:\quad \ell=0{,}22+1{,}57\lambda-1{,}8\lambda^2,\qquad H=2{,}61-3{,}75\lambda+5{,}24\lambda^2,
\]
\[
-0{,}1\le\lambda<0:\quad \ell=0{,}22+1{,}402\lambda+\frac{0{,}018\lambda}{\lambda+0{,}107},\qquad H=2{,}088+\frac{0{,}0731}{\lambda+0{,}14},
\]
\[
c_f=\frac{2\,\ell(\lambda)}{\Re_\theta},\qquad \Re_\theta=\frac{U_e\theta}{\nu}.
\]
Le \textbf{décollement laminaire} correspond à $\ell=0$, soit $\lambda\approx-0{,}09$ (constante \code{LAMBDA\_SEP}).

\begin{figure}[htbp]
\centering
% figure générée par tools/make_figures.py : thwaites_fermeture
\begin{tikzpicture}
\begin{groupplot}[group style={group size=2 by 1,horizontal sep=1.6cm},width=0.5\linewidth,height=5.4cm,
  xlabel={$\lambda=\theta^2U_e'/\nu$},xmin=-0.1,xmax=0.1,title style={font=\small},tick label style={font=\scriptsize},
  legend style={draw=none,font=\scriptsize}]
\nextgroupplot[title={Corrélations de fermeture},legend pos=north east,ymin=0,ymax=4]
\addplot[cblue,very thick] coordinates {(-0.09,3.55) (-0.088403,3.5048) (-0.086807,3.4622) (-0.08521,3.4222) (-0.083613,3.3844) (-0.082017,3.3487) (-0.08042,3.3149) (-0.078824,3.2829) (-0.077227,3.2525) (-0.07563,3.2236) (-0.074034,3.1961) (-0.072437,3.17) (-0.07084,3.145) (-0.069244,3.1211) (-0.067647,3.0983) (-0.06605,3.0765) (-0.064454,3.0556) (-0.062857,3.0356) (-0.061261,3.0164) (-0.059664,2.9979) (-0.058067,2.9802) (-0.056471,2.9631) (-0.054874,2.9467) (-0.053277,2.9309) (-0.051681,2.9157) (-0.050084,2.901) (-0.048487,2.8868) (-0.046891,2.8731) (-0.045294,2.8599) (-0.043697,2.8471) (-0.042101,2.8347) (-0.040504,2.8227) (-0.038908,2.8111) (-0.037311,2.7999) (-0.035714,2.789) (-0.034118,2.7784) (-0.032521,2.7681) (-0.030924,2.7582) (-0.029328,2.7485) (-0.027731,2.7391) (-0.026134,2.73) (-0.024538,2.7211) (-0.022941,2.7125) (-0.021345,2.7041) (-0.019748,2.6959) (-0.018151,2.6879) (-0.016555,2.6802) (-0.014958,2.6726) (-0.013361,2.6652) (-0.011765,2.658) (-0.010168,2.651) (-0.0085714,2.6442) (-0.0069748,2.6375) (-0.0053782,2.631) (-0.0037815,2.6246) (-0.0021849,2.6184) (-0.00058824,2.6123) (0.0010084,2.6062) (0.002605,2.6003) (0.0042017,2.5943) (0.0057983,2.5884) (0.007395,2.5826) (0.0089916,2.5767) (0.010588,2.5709) (0.012185,2.5651) (0.013782,2.5593) (0.015378,2.5536) (0.016975,2.5479) (0.018571,2.5422) (0.020168,2.5365) (0.021765,2.5309) (0.023361,2.5253) (0.024958,2.5197) (0.026555,2.5141) (0.028151,2.5086) (0.029748,2.5031) (0.031345,2.4976) (0.032941,2.4922) (0.034538,2.4867) (0.036134,2.4813) (0.037731,2.476) (0.039328,2.4706) (0.040924,2.4653) (0.042521,2.46) (0.044118,2.4548) (0.045714,2.4495) (0.047311,2.4443) (0.048908,2.4391) (0.050504,2.434) (0.052101,2.4288) (0.053697,2.4237) (0.055294,2.4187) (0.056891,2.4136) (0.058487,2.4086) (0.060084,2.4036) (0.061681,2.3986) (0.063277,2.3937) (0.064874,2.3888) (0.066471,2.3839) (0.068067,2.379) (0.069664,2.3742) (0.071261,2.3694) (0.072857,2.3646) (0.074454,2.3598) (0.07605,2.3551) (0.077647,2.3504) (0.079244,2.3457) (0.08084,2.3411) (0.082437,2.3365) (0.084034,2.3319) (0.08563,2.3273) (0.087227,2.3228) (0.088824,2.3183) (0.09042,2.3138) (0.092017,2.3093) (0.093613,2.3049) (0.09521,2.3005) (0.096807,2.2961) (0.098403,2.2917) (0.1,2.2874)};\addlegendentry{$H(\lambda)$}
\addplot[corange,very thick] coordinates {(-0.09,-0.014741) (-0.088403,0.10491) (-0.086807,0.20919) (-0.08521,0.30146) (-0.083613,0.38419) (-0.082017,0.45921) (-0.08042,0.5279) (-0.078824,0.59135) (-0.077227,0.65039) (-0.07563,0.7057) (-0.074034,0.75782) (-0.072437,0.80719) (-0.07084,0.85418) (-0.069244,0.89909) (-0.067647,0.94217) (-0.06605,0.98364) (-0.064454,1.0237) (-0.062857,1.0624) (-0.061261,1.1) (-0.059664,1.1366) (-0.058067,1.1723) (-0.056471,1.2071) (-0.054874,1.2412) (-0.053277,1.2745) (-0.051681,1.3073) (-0.050084,1.3394) (-0.048487,1.371) (-0.046891,1.4022) (-0.045294,1.4329) (-0.043697,1.4631) (-0.042101,1.493) (-0.040504,1.5225) (-0.038908,1.5517) (-0.037311,1.5805) (-0.035714,1.6091) (-0.034118,1.6374) (-0.032521,1.6655) (-0.030924,1.6933) (-0.029328,1.7209) (-0.027731,1.7482) (-0.026134,1.7754) (-0.024538,1.8024) (-0.022941,1.8292) (-0.021345,1.8559) (-0.019748,1.8824) (-0.018151,1.9087) (-0.016555,1.935) (-0.014958,1.961) (-0.013361,1.987) (-0.011765,2.0128) (-0.010168,2.0385) (-0.0085714,2.0642) (-0.0069748,2.0897) (-0.0053782,2.1151) (-0.0037815,2.1404) (-0.0021849,2.1656) (-0.00058824,2.1908) (0.0010084,2.2158) (0.002605,2.2408) (0.0042017,2.2656) (0.0057983,2.2904) (0.007395,2.3151) (0.0089916,2.3397) (0.010588,2.3642) (0.012185,2.3886) (0.013782,2.413) (0.015378,2.4372) (0.016975,2.4613) (0.018571,2.4854) (0.020168,2.5093) (0.021765,2.5332) (0.023361,2.5569) (0.024958,2.5806) (0.026555,2.6042) (0.028151,2.6277) (0.029748,2.6511) (0.031345,2.6744) (0.032941,2.6976) (0.034538,2.7208) (0.036134,2.7438) (0.037731,2.7668) (0.039328,2.7896) (0.040924,2.8124) (0.042521,2.835) (0.044118,2.8576) (0.045714,2.8801) (0.047311,2.9025) (0.048908,2.9248) (0.050504,2.947) (0.052101,2.9691) (0.053697,2.9911) (0.055294,3.0131) (0.056891,3.0349) (0.058487,3.0567) (0.060084,3.0783) (0.061681,3.0999) (0.063277,3.1214) (0.064874,3.1428) (0.066471,3.1641) (0.068067,3.1853) (0.069664,3.2064) (0.071261,3.2274) (0.072857,3.2483) (0.074454,3.2691) (0.07605,3.2899) (0.077647,3.3105) (0.079244,3.3311) (0.08084,3.3516) (0.082437,3.3719) (0.084034,3.3922) (0.08563,3.4124) (0.087227,3.4325) (0.088824,3.4525) (0.09042,3.4724) (0.092017,3.4923) (0.093613,3.512) (0.09521,3.5316) (0.096807,3.5512) (0.098403,3.5706) (0.1,3.59)};\addlegendentry{$10\,\ell(\lambda)$}
\addplot[cmuted,dashed] coordinates {(-0.09,0) (-0.09,4)};
\node[cmuted,font=\scriptsize,rotate=90,anchor=south] at (axis cs:-0.09,2.6) {décollement};
\nextgroupplot[title={$F(\lambda)=2[\ell-(2+H)\lambda]$ et sa droite},legend pos=north east]
\addplot[cblue,very thick] coordinates {(-0.09,0.99605) (-0.088403,0.99426) (-0.086807,0.99015) (-0.08521,0.98434) (-0.083613,0.97726) (-0.082017,0.96921) (-0.08042,0.96043) (-0.078824,0.9511) (-0.077227,0.94135) (-0.07563,0.93127) (-0.074034,0.92094) (-0.072437,0.91043) (-0.07084,0.89978) (-0.069244,0.88903) (-0.067647,0.87821) (-0.06605,0.86734) (-0.064454,0.85644) (-0.062857,0.84553) (-0.061261,0.83462) (-0.059664,0.82372) (-0.058067,0.81283) (-0.056471,0.80197) (-0.054874,0.79113) (-0.053277,0.78032) (-0.051681,0.76955) (-0.050084,0.75881) (-0.048487,0.74811) (-0.046891,0.73744) (-0.045294,0.72682) (-0.043697,0.71623) (-0.042101,0.70568) (-0.040504,0.69518) (-0.038908,0.68471) (-0.037311,0.67428) (-0.035714,0.66389) (-0.034118,0.65354) (-0.032521,0.64322) (-0.030924,0.63294) (-0.029328,0.6227) (-0.027731,0.61249) (-0.026134,0.60232) (-0.024538,0.59217) (-0.022941,0.58207) (-0.021345,0.57199) (-0.019748,0.56195) (-0.018151,0.55193) (-0.016555,0.54195) (-0.014958,0.53199) (-0.013361,0.52207) (-0.011765,0.51217) (-0.010168,0.50229) (-0.0085714,0.49245) (-0.0069748,0.48262) (-0.0053782,0.47283) (-0.0037815,0.46305) (-0.0021849,0.4533) (-0.00058824,0.44358) (0.0010084,0.43387) (0.002605,0.42419) (0.0042017,0.41452) (0.0057983,0.40488) (0.007395,0.39525) (0.0089916,0.38564) (0.010588,0.37605) (0.012185,0.36648) (0.013782,0.35692) (0.015378,0.34739) (0.016975,0.33787) (0.018571,0.32836) (0.020168,0.31888) (0.021765,0.30941) (0.023361,0.29996) (0.024958,0.29052) (0.026555,0.2811) (0.028151,0.2717) (0.029748,0.26231) (0.031345,0.25293) (0.032941,0.24358) (0.034538,0.23423) (0.036134,0.2249) (0.037731,0.21558) (0.039328,0.20628) (0.040924,0.19699) (0.042521,0.18772) (0.044118,0.17846) (0.045714,0.16921) (0.047311,0.15997) (0.048908,0.15074) (0.050504,0.14153) (0.052101,0.13233) (0.053697,0.12314) (0.055294,0.11396) (0.056891,0.1048) (0.058487,0.095641) (0.060084,0.086495) (0.061681,0.07736) (0.063277,0.068234) (0.064874,0.059119) (0.066471,0.050012) (0.068067,0.040916) (0.069664,0.031828) (0.071261,0.022748) (0.072857,0.013677) (0.074454,0.0046148) (0.07605,-0.0044399) (0.077647,-0.013487) (0.079244,-0.022526) (0.08084,-0.031559) (0.082437,-0.040584) (0.084034,-0.049603) (0.08563,-0.058615) (0.087227,-0.067621) (0.088824,-0.076622) (0.09042,-0.085616) (0.092017,-0.094606) (0.093613,-0.10359) (0.09521,-0.11257) (0.096807,-0.12154) (0.098403,-0.13051) (0.1,-0.13948)};\addlegendentry{$F$ exact (corrélations)}
\addplot[corange,dashed,very thick,domain=-0.09:0.1] {0.45-6*x};\addlegendentry{$0{,}45-6\lambda$}
\end{groupplot}
\end{tikzpicture}

\caption{À gauche, les corrélations $H(\lambda)$ et $\ell(\lambda)$ utilisées par le code. À droite, la fonction $F(\lambda)$ qu'elles impliquent, comparée à la droite $0{,}45-6\lambda$ de Thwaites. L'écart reste faible sur toute la plage utile, ce qui justifie la quadrature.}\label{fig:thwaites-f}
\end{figure}

\paragraph{Intuition physique de $\lambda$.} $\theta^2/\nu$ est le temps de diffusion visqueuse à travers la couche, et $1/U_e'$ le temps caractéristique de variation de l'écoulement extérieur. $\lambda$ compare les deux :
\begin{itemize}
\item $\lambda>0$ (accélération) : profil plein, $H$ petit, frottement élevé ;
\item $\lambda<0$ (décélération) : profil creusé près de la paroi, $H$ grand ;
\item $\lambda\approx-0{,}09$ : le frottement s'annule, la couche décolle.
\end{itemize}

\section{Deux cas exacts}
\paragraph{Plaque plane ($U_e=1$).} L'intégrale vaut $s$, d'où $\theta^2=0{,}45\,\nu s$, soit
\[
\theta=0{,}671\,\frac{s}{\sqrt{\Re_s}},
\]
contre $0{,}664$ pour la solution exacte de Blasius : 1~\% d'écart. Avec $\lambda=0$, on obtient $H=2{,}61$ (Blasius : 2,59).

\paragraph{Point d'arrêt ($U_e=ks$, Hiemenz).} L'intégrale vaut $k^5s^6/6$, d'où $\theta^2=0{,}075\,\nu/k$ : l'épaisseur est \emph{constante} près du point d'arrêt, et $\lambda=0{,}075$. C'est ce qui permet de démarrer la quadrature en $s=0$ sans singularité.

\begin{figure}[htbp]
\centering
% figure générée par tools/make_figures.py : plaque_plane
\begin{tikzpicture}
\begin{groupplot}[group style={group size=2 by 1,horizontal sep=1.6cm},width=0.5\linewidth,height=5.2cm,
  xlabel={$x/c$},xmin=0,xmax=1,title style={font=\small},tick label style={font=\scriptsize},
  legend style={draw=none,font=\scriptsize},legend pos=north west]
\nextgroupplot[title={Laminaire (Thwaites), $Re=10^7$},ylabel={$10^3\,\theta/c$}]
\addplot[cblue,very thick] coordinates {(0.0001,0.0015) (0.020148,0.030074) (0.040196,0.042504) (0.060244,0.052046) (0.080292,0.060091) (0.10034,0.067179) (0.12039,0.073588) (0.14044,0.079482) (0.16048,0.084968) (0.18053,0.090121) (0.20058,0.094994) (0.22063,0.09963) (0.24068,0.10406) (0.26073,0.10831) (0.28077,0.11239) (0.30082,0.11634) (0.32087,0.12015) (0.34092,0.12385) (0.36097,0.12744) (0.38101,0.13093) (0.40106,0.13433) (0.42111,0.13765) (0.44116,0.14089) (0.46121,0.14406) (0.48125,0.14715) (0.5013,0.15019) (0.52135,0.15316) (0.5414,0.15608) (0.56145,0.15894) (0.5815,0.16176) (0.60154,0.16452) (0.62159,0.16724) (0.64164,0.16992) (0.66169,0.17255) (0.68174,0.17515) (0.70178,0.1777) (0.72183,0.18022) (0.74188,0.18271) (0.76193,0.18516) (0.78198,0.18758) (0.80202,0.18997) (0.82207,0.19233) (0.84212,0.19466) (0.86217,0.19697) (0.88222,0.19924) (0.90227,0.20149) (0.92231,0.20372) (0.94236,0.20592) (0.96241,0.2081) (0.98246,0.21026)};\addlegendentry{Thwaites (code)}
\addplot[cgreen,dashed,very thick] coordinates {(0.0001,0.0020998) (0.020148,0.029805) (0.040196,0.042098) (0.060244,0.051538) (0.080292,0.059498) (0.10034,0.066513) (0.12039,0.072855) (0.14044,0.078688) (0.16048,0.084117) (0.18053,0.089217) (0.20058,0.09404) (0.22063,0.098628) (0.24068,0.10301) (0.26073,0.10722) (0.28077,0.11126) (0.30082,0.11517) (0.32087,0.11894) (0.34092,0.1226) (0.36097,0.12615) (0.38101,0.12961) (0.40106,0.13298) (0.42111,0.13626) (0.44116,0.13947) (0.46121,0.1426) (0.48125,0.14567) (0.5013,0.14867) (0.52135,0.15161) (0.5414,0.1545) (0.56145,0.15733) (0.5815,0.16012) (0.60154,0.16286) (0.62159,0.16555) (0.64164,0.1682) (0.66169,0.1708) (0.68174,0.17337) (0.70178,0.1759) (0.72183,0.1784) (0.74188,0.18086) (0.76193,0.18328) (0.78198,0.18568) (0.80202,0.18805) (0.82207,0.19038) (0.84212,0.19269) (0.86217,0.19497) (0.88222,0.19722) (0.90227,0.19945) (0.92231,0.20165) (0.94236,0.20383) (0.96241,0.20599) (0.98246,0.20813)};\addlegendentry{Blasius}
\nextgroupplot[title={Turbulent dès $x=0$ (Head)}]
\addplot[corange,very thick] coordinates {(0.0001,0.0015) (0.020148,0.054417) (0.040196,0.097004) (0.060244,0.13614) (0.080292,0.17316) (0.10034,0.20865) (0.12039,0.24297) (0.14044,0.27633) (0.16048,0.3089) (0.18053,0.34077) (0.20058,0.37205) (0.22063,0.4028) (0.24068,0.43307) (0.26073,0.4629) (0.28077,0.49235) (0.30082,0.52143) (0.32087,0.55017) (0.34092,0.57861) (0.36097,0.60676) (0.38101,0.63464) (0.40106,0.66226) (0.42111,0.68964) (0.44116,0.71679) (0.46121,0.74373) (0.48125,0.77047) (0.5013,0.79701) (0.52135,0.82336) (0.5414,0.84954) (0.56145,0.87554) (0.5815,0.90139) (0.60154,0.92707) (0.62159,0.95261) (0.64164,0.978) (0.66169,1.0033) (0.68174,1.0284) (0.70178,1.0534) (0.72183,1.0782) (0.74188,1.1029) (0.76193,1.1276) (0.78198,1.1521) (0.80202,1.1765) (0.82207,1.2008) (0.84212,1.2249) (0.86217,1.249) (0.88222,1.273) (0.90227,1.2969) (0.92231,1.3207) (0.94236,1.3444) (0.96241,1.368) (0.98246,1.3915)};\addlegendentry{Head (code)}
\addplot[cgreen,dashed,very thick] coordinates {(0.0001,0.00090428) (0.020148,0.063051) (0.040196,0.10956) (0.060244,0.15144) (0.080292,0.19056) (0.10034,0.22776) (0.12039,0.26349) (0.14044,0.29805) (0.16048,0.33163) (0.18053,0.36437) (0.20058,0.3964) (0.22063,0.42779) (0.24068,0.45862) (0.26073,0.48893) (0.28077,0.51879) (0.30082,0.54821) (0.32087,0.57725) (0.34092,0.60593) (0.36097,0.63427) (0.38101,0.6623) (0.40106,0.69004) (0.42111,0.7175) (0.44116,0.74469) (0.46121,0.77165) (0.48125,0.79837) (0.5013,0.82486) (0.52135,0.85115) (0.5414,0.87724) (0.56145,0.90313) (0.5815,0.92884) (0.60154,0.95437) (0.62159,0.97973) (0.64164,1.0049) (0.66169,1.03) (0.68174,1.0549) (0.70178,1.0796) (0.72183,1.1042) (0.74188,1.1287) (0.76193,1.153) (0.78198,1.1772) (0.80202,1.2013) (0.82207,1.2253) (0.84212,1.2491) (0.86217,1.2728) (0.88222,1.2965) (0.90227,1.32) (0.92231,1.3434) (0.94236,1.3667) (0.96241,1.3899) (0.98246,1.413)};\addlegendentry{loi en $1/7$}
\end{groupplot}
\end{tikzpicture}

\caption{Validation sur plaque plane à $\Re=10^7$ (tests V2 et V3 de \code{validation.py}). À gauche, Thwaites contre Blasius : 1,0~\% d'écart, la différence entre les constantes 0,671 et 0,664. À droite, Head contre la loi turbulente en $1/7$ (ch.~\ref{ch:head}) : moins de 5~\% d'écart au-delà de $x=0{,}3$.}\label{fig:plaque}
\end{figure}

\begin{codebox}[ : bloc Thwaites de \code{march\_surface}]
\begin{lstlisting}
s0 = np.concatenate([[0.0], s])          # on ajoute le point d'arret, s = 0
U0 = np.concatenate([[0.0], Ue])         # ou Ue = 0
integ = np.concatenate([[0.0], np.cumsum(0.5 * (U0[1:]**5 + U0[:-1]**5) * np.diff(s0))])[1:]
theta_lam = np.sqrt(0.45 / Re * integ / Ue**6)       # nu = 1/Re
lam = theta_lam**2 * dUe * Re                         # dUe = np.gradient(Ue, s)
H_lam, ell = thwaites_closure(lam)
Re_th = Ue * theta_lam * Re
\end{lstlisting}
La quadrature est une méthode des trapèzes sur $U_e^5$. Le point $s=0$, $U_e=0$ est ajouté explicitement : sans lui, l'intégrale démarrerait à la première collocation et $\theta$ serait faux sur les premières stations. \code{thwaites\_closure} borne $\lambda$ à $[-0{,}1;\,0{,}1]$ avant d'évaluer les corrélations.
\end{codebox}

\begin{piege}
$\theta$ dépend de tout l'historique amont (l'intégrale), alors que $\lambda$, donc $H$, est \emph{local}. Thwaites réagit donc mal aux variations brutales de $U_e$ : bulbe de décollement, choc, ou la chute de vitesse au bord de fuite fermé. De plus, au-delà de $\lambda<-0{,}09$, la méthode n'a plus de sens. On ne la prolonge pas : on bascule en turbulent (ch.~\ref{ch:transition}).
\end{piege}

\exercices
\begin{exercice}
Retrouver le résultat de Hiemenz : pour $U_e=ks$, calculer $\theta^2$ par la formule de Thwaites, puis $\lambda$. Que vaut $H$ au point d'arrêt d'après les corrélations ?
\end{exercice}
\begin{exercice}
Pourquoi la puissance 5 ? Reprendre le calcul avec une droite quelconque $F=a-b\lambda$ et montrer qu'on obtient $\theta^2U_e^{b}=a\,\nu\int_0^sU_e^{b-1}\,\dd s'$.
\end{exercice}

\chapter{La transition laminaire--turbulent}\label{ch:transition}
\iterbanner{V0}

La position de la transition est le paramètre qui pèse le plus sur la traînée d'un profil à $\Re\sim10^6$ : une couche limite turbulente frotte environ cinq fois plus qu'une couche laminaire. C'est aussi, on le verra, la principale source d'irrégularité de $C_D(m,p,t)$.

\section{Trois événements, le premier l'emporte}
Le code déclenche la transition à la première station où l'un des trois critères suivants est franchi.

\paragraph{1. Le critère de Michel (transition naturelle).} Corrélation empirique à un paramètre (Michel, ONERA, 1951) : la transition a lieu quand le Reynolds de quantité de mouvement dépasse une valeur qui dépend du Reynolds local,
\[
\Re_\theta>\Re_{\theta,tr}(\Re_x)=1{,}174\left(1+\frac{22\,400}{\Re_x}\right)\Re_x^{0{,}46},\qquad \Re_x=\frac{U_e\,s}{\nu}.
\]
Il représente la transition « naturelle » par amplification des ondes de Tollmien-Schlichting dans un environnement calme. Il ne contient ni le gradient de pression explicitement (seulement via l'historique dans $\theta$) ni le taux de turbulence extérieure.

\paragraph{2. Le décollement laminaire.} Si $\lambda<-0{,}09$ avant que Michel ne soit atteint, la couche laminaire décolle. Physiquement, elle forme le plus souvent un \textbf{bulbe de décollement} : la couche décollée transitionne très vite et recolle en turbulent. Le modèle suppose un bulbe court, réduit à une transition instantanée au point de décollement.

\paragraph{3. Le déclenchement forcé.} Une abscisse \code{x\_trip} imposée par l'utilisateur, qui représente une bande de rugosité. C'est la même sémantique que la commande \texttt{XTR} d'XFOIL : la transition a lieu au premier des deux, critère libre ou \code{x\_trip}. L'option \code{free\_transition=False} désactive Michel, pour pouvoir imposer une transition \emph{plus tardive} que celle de Michel, par exemple celle trouvée par XFOIL.

\begin{figure}[htbp]
\centering
% figure générée par tools/make_figures.py : transition_michel
\begin{tikzpicture}
\begin{groupplot}[group style={group size=2 by 1,horizontal sep=1.6cm},width=0.5\linewidth,height=5.4cm,
  xlabel={$x/c$},xmin=0,xmax=1,ymin=0,ymax=2000,title style={font=\small},tick label style={font=\scriptsize},
  legend style={draw=none,font=\scriptsize},legend pos=north west]
\nextgroupplot[title={extrados : transition en $x/c=0.279$ (michel)}]
\addplot[cblue,very thick] coordinates {(0.0038514,84.921) (0.0060621,102.74) (0.0087685,121.8) (0.011969,141.78) (0.01566,162.43) (0.01984,183.55) (0.024505,204.98) (0.029651,226.61) (0.035273,248.36) (0.041368,270.15) (0.047928,291.96) (0.054949,313.74) (0.062423,335.49) (0.070344,357.2) (0.078703,378.85) (0.087494,400.46) (0.096707,422.03) (0.10633,443.56) (0.11636,465.06) (0.12679,486.55) (0.13759,508.04) (0.14877,529.53) (0.16031,551.04) (0.17219,572.58) (0.18442,594.17) (0.19696,615.81) (0.20982,637.52) (0.22297,659.32) (0.2364,681.22) (0.25011,703.22) (0.26407,725.36) (0.27826,747.65) (0.29268,770.11) (0.30731,792.76) (0.32214,815.64) (0.33714,838.79) (0.3523,862.29) (0.36761,886.23) (0.38305,910.85) (0.39857,936.73) (0.41414,962.69) (0.42974,987.74) (0.44542,1012.1) (0.46115,1036.1) (0.47691,1059.9) (0.49269,1083.5) (0.50848,1106.9) (0.52425,1130.2) (0.54,1153.4) (0.5557,1176.6) (0.57135,1199.7) (0.58692,1222.7) (0.6024,1245.7) (0.61777,1268.7) (0.63302,1291.7) (0.64814,1314.8) (0.6631,1337.9) (0.6779,1361) (0.69252,1384.2) (0.70695,1407.6) (0.72116,1431.1) (0.73515,1454.8) (0.74891,1478.6) (0.76242,1502.8) (0.77566,1527.2) (0.78863,1551.9) (0.80131,1577) (0.81369,1602.5) (0.82576,1628.5) (0.8375,1655) (0.84891,1682.1) (0.85997,1709.9) (0.87068,1738.4) (0.88102,1767.8) (0.89098,1798.1) (0.90055,1829.5) (0.90973,1862) (0.9185,1895.9) (0.92686,1931.2) (0.9348,1968.2) (0.94232,2007.1) (0.9494,2048.2) (0.95603,2091.8) (0.96222,2138.4) (0.96795,2188.4) (0.97323,2242.6) (0.97804,2301.8) (0.98238,2367.2)};\addlegendentry{$Re_\theta$ (Thwaites)}
\addplot[corange,very thick] coordinates {(0.0038514,254.52) (0.0060621,268.16) (0.0087685,282.68) (0.011969,297.72) (0.01566,313.1) (0.01984,328.77) (0.024505,344.68) (0.029651,360.8) (0.035273,377.11) (0.041368,393.57) (0.047928,410.16) (0.054949,426.86) (0.062423,443.63) (0.070344,460.45) (0.078703,477.29) (0.087494,494.13) (0.096707,510.95) (0.10633,527.72) (0.11636,544.43) (0.12679,561.06) (0.13759,577.59) (0.14877,593.99) (0.16031,610.27) (0.17219,626.39) (0.18442,642.34) (0.19696,658.12) (0.20982,673.71) (0.22297,689.09) (0.2364,704.25) (0.25011,719.19) (0.26407,733.89) (0.27826,748.33) (0.29268,762.51) (0.30731,776.42) (0.32214,790.04) (0.33714,803.36) (0.3523,816.36) (0.36761,829.01) (0.38305,841.27) (0.39857,852.97) (0.41414,864.32) (0.42974,875.56) (0.44542,886.67) (0.46115,897.59) (0.47691,908.3) (0.49269,918.78) (0.50848,929.02) (0.52425,939.03) (0.54,948.78) (0.5557,958.29) (0.57135,967.53) (0.58692,976.51) (0.6024,985.23) (0.61777,993.68) (0.63302,1001.9) (0.64814,1009.8) (0.6631,1017.4) (0.6779,1024.7) (0.69252,1031.8) (0.70695,1038.5) (0.72116,1045) (0.73515,1051.2) (0.74891,1057.1) (0.76242,1062.6) (0.77566,1067.9) (0.78863,1072.8) (0.80131,1077.4) (0.81369,1081.7) (0.82576,1085.6) (0.8375,1089.2) (0.84891,1092.5) (0.85997,1095.4) (0.87068,1098) (0.88102,1100.2) (0.89098,1102) (0.90055,1103.4) (0.90973,1104.5) (0.9185,1105.1) (0.92686,1105.4) (0.9348,1105.2) (0.94232,1104.5) (0.9494,1103.5) (0.95603,1101.9) (0.96222,1099.8) (0.96795,1097.2) (0.97323,1094) (0.97804,1090.1) (0.98238,1085.5)};\addlegendentry{seuil de Michel}
\addplot[cmuted,dashed] coordinates {(0.2795,0) (0.2795,2000)};
\nextgroupplot[title={intrados : transition en $x/c=0.478$ (michel)}]
\addplot[cblue,very thick] coordinates {(0.0040239,28.372) (0.006247,39.504) (0.0089501,52.466) (0.01213,66.77) (0.015782,82.089) (0.019902,98.191) (0.024485,114.91) (0.029527,132.12) (0.035021,149.71) (0.040961,167.62) (0.047342,185.77) (0.054155,204.13) (0.061395,222.64) (0.069053,241.28) (0.077121,260.03) (0.085592,278.87) (0.094457,297.77) (0.10371,316.72) (0.11333,335.71) (0.12333,354.73) (0.13367,373.77) (0.14437,392.82) (0.1554,411.87) (0.16676,430.91) (0.17844,449.93) (0.19042,468.92) (0.20269,487.86) (0.21525,506.75) (0.22808,525.58) (0.24116,544.32) (0.25449,562.95) (0.26806,581.46) (0.28185,599.83) (0.29585,618.03) (0.31005,636.03) (0.32443,653.8) (0.33899,671.28) (0.3537,688.43) (0.36855,705.14) (0.38354,721.27) (0.39866,736.44) (0.41396,751.01) (0.42937,766.2) (0.44486,781.87) (0.4604,797.7) (0.47599,813.57) (0.4916,829.42) (0.50723,845.18) (0.52285,860.84) (0.53845,876.36) (0.55402,891.72) (0.56954,906.92) (0.58499,921.92) (0.60037,936.74) (0.61565,951.35) (0.63082,965.75) (0.64587,979.95) (0.66078,993.95) (0.67553,1007.7) (0.69012,1021.3) (0.70452,1034.7) (0.71872,1047.9) (0.73272,1060.9) (0.74649,1073.8) (0.76001,1086.5) (0.77329,1099.1) (0.7863,1111.6) (0.79903,1124) (0.81147,1136.3) (0.82361,1148.6) (0.83542,1160.9) (0.84691,1173.1) (0.85806,1185.5) (0.86886,1198) (0.8793,1210.6) (0.88936,1223.4) (0.89904,1236.4) (0.90832,1249.8) (0.9172,1263.6) (0.92567,1277.8) (0.93372,1292.7) (0.94134,1308.3) (0.94852,1324.7) (0.95526,1342.3) (0.96155,1361) (0.96737,1381.4) (0.97274,1403.7) (0.97763,1428.5) (0.98205,1456.4)};\addlegendentry{$Re_\theta$ (Thwaites)}
\addplot[corange,very thick] coordinates {(0.0040239,300.71) (0.006247,247.32) (0.0089501,235.02) (0.01213,237.11) (0.015782,245.52) (0.019902,257.09) (0.024485,270.39) (0.029527,284.69) (0.035021,299.6) (0.040961,314.88) (0.047342,330.38) (0.054155,346.02) (0.061395,361.72) (0.069053,377.43) (0.077121,393.12) (0.085592,408.77) (0.094457,424.33) (0.10371,439.81) (0.11333,455.18) (0.12333,470.44) (0.13367,485.56) (0.14437,500.55) (0.1554,515.39) (0.16676,530.08) (0.17844,544.61) (0.19042,558.98) (0.20269,573.19) (0.21525,587.23) (0.22808,601.11) (0.24116,614.82) (0.25449,628.36) (0.26806,641.73) (0.28185,654.94) (0.29585,667.99) (0.31005,680.89) (0.32443,693.63) (0.33899,706.23) (0.3537,718.71) (0.36855,731.07) (0.38354,743.36) (0.39866,755.68) (0.41396,768.01) (0.42937,780.01) (0.44486,791.69) (0.4604,803.1) (0.47599,814.26) (0.4916,825.19) (0.50723,835.9) (0.52285,846.39) (0.53845,856.66) (0.55402,866.71) (0.56954,876.55) (0.58499,886.18) (0.60037,895.59) (0.61565,904.79) (0.63082,913.77) (0.64587,922.54) (0.66078,931.08) (0.67553,939.41) (0.69012,947.51) (0.70452,955.38) (0.71872,963.02) (0.73272,970.43) (0.74649,977.61) (0.76001,984.53) (0.77329,991.21) (0.7863,997.64) (0.79903,1003.8) (0.81147,1009.7) (0.82361,1015.4) (0.83542,1020.7) (0.84691,1025.8) (0.85806,1030.6) (0.86886,1035) (0.8793,1039.2) (0.88936,1043) (0.89904,1046.6) (0.90832,1049.7) (0.9172,1052.6) (0.92567,1055) (0.93372,1057.1) (0.94134,1058.7) (0.94852,1060) (0.95526,1060.8) (0.96155,1061.1) (0.96737,1060.8) (0.97274,1060) (0.97763,1058.6) (0.98205,1056.4)};\addlegendentry{seuil de Michel}
\addplot[cmuted,dashed] coordinates {(0.4782,0) (0.4782,2000)};
\end{groupplot}
\end{tikzpicture}

\caption{NACA 2412, $\alpha=2^\circ$, $\Re=3\cdot10^6$ : $\Re_\theta$ laminaire calculé par Thwaites (bleu) et seuil de Michel (orange) le long de chaque face. La transition a lieu au croisement, interpolé entre les deux stations qui l'encadrent.}\label{fig:michel}
\end{figure}

\section{Interpoler la position de transition}
Chaque critère s'écrit comme une fonction $g_i$ aux stations, avec la convention « franchi si $g\ge0$ » :
\[
g^{\text{Michel}}=\Re_\theta-\Re_{\theta,tr}(\Re_x),\qquad g^{\text{décollement}}=-0{,}09-\lambda,\qquad g^{\text{forcé}}=x-x_{trip}.
\]
À la première station $i$ où l'un d'eux devient positif, on interpole linéairement le zéro entre $i-1$ et $i$ :
\[
w=\frac{-g_{i-1}}{g_i-g_{i-1}}\in[0,1],\qquad s_{tr}=s_{i-1}+w\,(s_i-s_{i-1}).
\]
La condition initiale de Head (ch.~\ref{ch:head}) est prise en $s_{tr}$, avec $\theta$ interpolé au même endroit.

\paragraph{Pourquoi c'est important.} Sans interpolation, $x_{tr}$ ne peut prendre que les abscisses des stations. Quand $(m,p,t)$ varie continûment, $x_{tr}$ saute d'une station à la suivante et $C_D$ devient une fonction en escalier. Mesuré sur le NACA 2412, $C_D$ variait de $\pm1{,}5$~\% quand on changeait le nombre de panneaux. Avec l'interpolation, il varie de $\pm0{,}1$~\%. Pour un ajustement quadratique, des marches sont du bruit pur.

\begin{codebox}[ : recherche de la transition dans \code{march\_surface}]
\begin{lstlisting}
crit = {}
if x_trip is not None:
    crit["forced"] = x - x_trip
crit["laminar_sep"] = np.where(s > 0.002, LAMBDA_SEP - lam, -1.0)
if free_transition:
    crit["michel"] = Re_th - michel_transition(Re_x)
i_tr, mode, w = n, "none", 0.0
for i in range(n):
    hit = [k for k, g in crit.items() if g[i] >= 0]
    if hit:
        i_tr, mode = i, hit[0]
        if i > 0:
            g0, g1 = crit[mode][i - 1], crit[mode][i]
            w = float(np.clip(-g0 / (g1 - g0), 0.0, 1.0)) if g1 != g0 else 1.0
        break
\end{lstlisting}
Le décollement laminaire est ignoré sur les deux premiers millièmes de corde ($s\le0{,}002$) : juste après le point d'arrêt, $U_e'$ est mal estimé par différences finies et $\lambda$ peut être faussement très négatif. L'ordre des clés du dictionnaire fixe la priorité en cas d'égalité sur une même station : forcé, puis décollement laminaire, puis Michel.
\end{codebox}

\section{Comportement sur un profil}
\begin{figure}[htbp]
\centering
% figure générée par tools/make_figures.py : xtr_alpha
\begin{tikzpicture}
\begin{axis}[width=0.7\linewidth,height=5.6cm,xlabel={$\alpha$ (°)},ylabel={$x_{tr}/c$},xmin=-4,xmax=10,ymin=0,ymax=1,
  legend style={draw=none,font=\scriptsize,at={(0.02,0.02)},anchor=south west}]
\addplot[cblue,very thick,mark=*,mark size=1.4pt] coordinates {(-4,0.66149) (-3,0.59365) (-2,0.52539) (-1,0.46) (0,0.40359) (1,0.34723) (2,0.27945) (3,0.20801) (4,0.14399) (5,0.098468) (6,0.070947) (7,0.054723) (8,0.022071) (9,0.013933) (10,0.010145)};\addlegendentry{extrados}
\addplot[corange,very thick,mark=square*,mark size=1.4pt] coordinates {(-4,0.072542) (-3,0.096003) (-2,0.13131) (-1,0.18121) (0,0.24686) (1,0.33239) (2,0.47818) (3,0.61065) (4,0.7494) (5,0.8636) (6,0.9027) (7,0.91429) (8,0.92399) (9,0.93213) (10,0.93901)};\addlegendentry{intrados}
\addplot[only marks,mark=o,mark size=3.5pt,cink,thick] coordinates {(6,0.9027) (7,0.91429) (8,0.92399) (9,0.93213) (10,0.93901) (8,0.022071) (9,0.013933) (10,0.010145)};\addlegendentry{décollement laminaire}
\end{axis}
\end{tikzpicture}

\caption{Position de transition du NACA 2412 en fonction de l'incidence ($\Re=3\cdot10^6$). Quand $\alpha$ augmente, la transition remonte vers le bord d'attaque sur l'extrados (gradient adverse plus fort après le pic d'aspiration) et recule sur l'intrados. Cercles : transition par décollement laminaire plutôt que par Michel.}\label{fig:xtr-alpha}
\end{figure}

La figure~\ref{fig:xtr-alpha} montre le comportement attendu, et une bascule de mode. Sur l'intrados, à partir de $\alpha=6^\circ$, la couche laminaire décolle avant d'atteindre le seuil de Michel. Sur l'extrados, la même bascule se produit vers $\alpha=8^\circ$, sous le pic d'aspiration. Chaque bascule crée une rupture de pente dans $x_{tr}(\alpha)$, donc dans $C_D(\alpha)$.

\begin{figure}[htbp]
\centering
% figure générée par tools/make_figures.py : cd_xtrip
\begin{tikzpicture}
\begin{axis}[width=0.7\linewidth,height=5.4cm,xlabel={transition imposée $x_{tr}/c$ (deux faces)},ylabel={$C_D$ (counts)},
  xmin=0,xmax=0.8,legend style={draw=none,font=\scriptsize},legend pos=north east]
\addplot[cblue,very thick,mark=*,mark size=1.2pt] coordinates {(0.02,97.593) (0.05,95.511) (0.08,93.023) (0.11,90.368) (0.14,87.62) (0.17,84.83) (0.2,82.031) (0.23,79.231) (0.26,76.435) (0.29,73.664) (0.32,70.911) (0.35,68.203) (0.38,65.532) (0.41,63.877) (0.44,63.043) (0.47,62.211) (0.5,61.374) (0.53,60.543) (0.56,59.706) (0.59,58.873) (0.62,58.04) (0.65,57.206) (0.68,56.373) (0.71,55.538) (0.74,54.704) (0.77,53.87) (0.8,53.035)};\addlegendentry{transition imposée}
\addplot[corange,dashed,very thick] coordinates {(0,69.17) (0.8,69.17)};\addlegendentry{transition libre (Michel)}
\end{axis}
\end{tikzpicture}

\caption{$C_D$ du NACA 2412 ($\alpha=2^\circ$) quand on impose la même position de transition sur les deux faces. De $x_{tr}=0{,}02$ à $0{,}8$, la traînée passe de 98 à 53~counts. En transition libre, Michel donne 69~counts.}\label{fig:cd-xtrip}
\end{figure}

La figure~\ref{fig:cd-xtrip} chiffre l'enjeu : selon l'endroit où la transition a lieu, la traînée du même profil varie presque du simple au double. Toute erreur sur le critère de transition se retrouve directement dans $C_D$.

\begin{piege}
Michel est connu pour être \textbf{trop précoce à bas Reynolds}. Sur le NACA 0012 à $\alpha=0$ et $\Re=10^6$, il fait transitionner en $x/c=0{,}46$, et donne $C_D=75$~counts contre environ 54 couramment cités pour XFOIL, à reconfirmer au test V5. Si l'on impose la transition au décollement laminaire ($x\approx0{,}62$), on obtient 62~counts. L'erreur de transition domine donc l'erreur de couche limite. La méthode $e^N$ d'XFOIL suit l'amplification des ondes instables et voit l'effet stabilisant d'un gradient favorable ; Michel ne le voit pas.
\end{piege}

\begin{projet}
La bascule entre transition de Michel et transition par décollement laminaire est la principale non-régularité de $C_D(m,p,t)$ sur la grille. C'est elle qui limite le $R^2$ du modèle quadratique à environ 0,91 (ch.~\ref{ch:espace}). Deux leviers existent si la phase~4 en a besoin. On peut forcer la transition près du bord d'attaque (cas « tout turbulent », plus régulier et défendable pour un drone de série aux surfaces imparfaites). On peut aussi caler le critère sur XFOIL.
\end{projet}

\exercices
\begin{exercice}
Sur plaque plane, Thwaites donne $\Re_\theta=0{,}671\sqrt{\Re_x}$. Écrire l'équation qui donne le $\Re_x$ de transition selon Michel, et la résoudre numériquement (quelques itérations suffisent). En déduire $x_{tr}$ pour $\Re=3\cdot10^6$.
\end{exercice}
\begin{exercice}
On fait varier continûment l'épaisseur $t$ d'un profil. Expliquer pourquoi, sans interpolation de $x_{tr}$, la courbe $C_D(t)$ présenterait de petites marches. Quel serait l'effet de ces marches sur un modèle quadratique ajusté par moindres carrés, et sur un optimiseur par gradient ?
\end{exercice}

\chapter{Couche limite turbulente : la méthode d'entraînement de Head}\label{ch:head}
\iterbanner{V0}

\section{Pourquoi une deuxième équation}
En laminaire, la forme du profil de vitesse suit le gradient de pression local, et un seul paramètre ($\lambda$) suffit. En turbulent, ce n'est plus vrai : la forme moyenne dépend de l'histoire de la couche, de la façon dont la turbulence s'est développée en amont. Il faut donc suivre $H$ comme une inconnue à part entière, et lui donner sa propre équation d'évolution.

\section{L'idée de Head : l'entraînement}
Le débit volumique qui traverse la couche limite, par unité d'envergure, vaut
\[
Q=\int_0^\delta u\,\dd y=U_e\,(\delta-\delta^*)=U_e\,\theta\,H_1,\qquad H_1=\frac{\delta-\delta^*}{\theta}.
\]
Ce débit ne peut augmenter que si la couche \textbf{entraîne} du fluide extérieur irrotationnel à travers sa frontière, l'interface turbulent--non turbulent (figure~\ref{fig:entrainement}). Head (1958) suppose que la vitesse d'entraînement, rapportée à $U_e$, ne dépend que de la forme du profil :
\[
\frac{1}{U_e}\frac{\dd Q}{\dd s}=\frac{1}{U_e}\frac{\dd(U_e\theta H_1)}{\dd s}=F(H_1),
\]
et que $H_1$ est lui-même une fonction de $H$. Avec l'équation de von Kármán, on obtient un système fermé de deux EDO.

\begin{figure}[htbp]
\centering
\begin{tikzpicture}[x=1.1cm,y=1.1cm]
\fill[pattern=north east lines,pattern color=cmuted] (0,-0.25) rectangle (10,0);
\draw[cink,thick] (0,0) -- (10,0);
\draw[cblue,very thick,decorate,decoration={snake,amplitude=0.6mm,segment length=4mm}] (0,0.8) .. controls (4,1.3) and (7,1.7) .. (10,2.1);
\node[cblue,font=\small] at (8.2,2.25) {frontière turbulente $\delta(s)$};
\foreach \x in {2.5,5.5,8}{\draw[corange,thick,-{Stealth[length=2mm]}] (\x-0.3,{0.8+0.13*\x+0.7}) -- (\x,{0.8+0.13*\x+0.12});}
\node[corange,font=\small] at (3.2,2.4) {entraînement de fluide extérieur};
\draw[cink,-{Stealth[length=2.5mm]},thick] (-1.2,1.0) -- (-0.2,1.0) node[midway,above,font=\small] {$U_e$};
\node[font=\small,align=center] at (5,0.55) {débit $Q=U_e(\delta-\delta^*)$ croissant};
\end{tikzpicture}
\caption{L'idée de Head : le débit dans la couche limite turbulente n'augmente que par entraînement de fluide extérieur à travers la frontière turbulente.}\label{fig:entrainement}
\end{figure}

\section{Le système résolu}
\[
\frac{\dd\theta}{\dd s}=\frac{c_f}{2}-\big(H+2\big)\frac{\theta}{U_e}\frac{\dd U_e}{\dd s},\qquad
\frac{\dd E}{\dd s}=U_e\,F(H_1),\qquad E=U_e\,\theta\,H_1,
\]
avec les fermetures empiriques suivantes (Head, puis Cebeci et Bradshaw) :
\[
F(H_1)=0{,}0306\,(H_1-3)^{-0{,}6169},
\]
\[
H_1(H)=\begin{cases}3{,}3+0{,}8234\,(H-1{,}1)^{-1{,}287} & H\le1{,}6,\\ 3{,}3+1{,}5501\,(H-0{,}6778)^{-3{,}064} & H>1{,}6,\end{cases}
\]
\[
c_f=0{,}246\cdot10^{-0{,}678\,H}\,\Re_\theta^{-0{,}268}\qquad\text{(Ludwieg-Tillmann).}
\]
Les inconnues intégrées sont $\theta$ et $E$. À chaque pas, on en déduit $H_1=E/(U_e\theta)$, puis $H$ en inversant $H_1(H)$, puis $c_f$.

\begin{figure}[htbp]
\centering
% figure générée par tools/make_figures.py : head_fermeture
\begin{tikzpicture}
\begin{groupplot}[group style={group size=2 by 1,horizontal sep=1.6cm},width=0.5\linewidth,height=5.2cm,
  title style={font=\small},tick label style={font=\scriptsize}]
\nextgroupplot[title={$H_1(H)$},xlabel={$H$},ylabel={$H_1=(\delta-\delta^*)/\theta$},ymin=3,ymax=12]
\addplot[cblue,very thick] coordinates {(1.15,42.208) (1.1639,31.694) (1.1777,25.35) (1.1916,21.151) (1.2055,18.19) (1.2193,16.001) (1.2332,14.326) (1.2471,13.006) (1.2609,11.943) (1.2748,11.071) (1.2887,10.344) (1.3025,9.7296) (1.3164,9.2043) (1.3303,8.7507) (1.3441,8.3556) (1.358,8.0086) (1.3718,7.7018) (1.3857,7.4288) (1.3996,7.1845) (1.4134,6.9648) (1.4273,6.7662) (1.4412,6.586) (1.455,6.4218) (1.4689,6.2716) (1.4828,6.1338) (1.4966,6.0069) (1.5105,5.8898) (1.5244,5.7814) (1.5382,5.6809) (1.5521,5.5873) (1.566,5.5001) (1.5798,5.4186) (1.5937,5.3423) (1.6076,5.2376) (1.6214,5.1517) (1.6353,5.0708) (1.6492,4.9944) (1.663,4.9224) (1.6769,4.8544) (1.6908,4.7901) (1.7046,4.7294) (1.7185,4.6718) (1.7324,4.6173) (1.7462,4.5656) (1.7601,4.5166) (1.7739,4.47) (1.7878,4.4258) (1.8017,4.3838) (1.8155,4.3438) (1.8294,4.3058) (1.8433,4.2696) (1.8571,4.2351) (1.871,4.2022) (1.8849,4.1708) (1.8987,4.1409) (1.9126,4.1123) (1.9265,4.085) (1.9403,4.0588) (1.9542,4.0339) (1.9681,4.01) (1.9819,3.9871) (1.9958,3.9652) (2.0097,3.9442) (2.0235,3.9241) (2.0374,3.9048) (2.0513,3.8863) (2.0651,3.8685) (2.079,3.8514) (2.0929,3.835) (2.1067,3.8193) (2.1206,3.8042) (2.1345,3.7896) (2.1483,3.7756) (2.1622,3.7621) (2.1761,3.7491) (2.1899,3.7366) (2.2038,3.7246) (2.2176,3.713) (2.2315,3.7018) (2.2454,3.691) (2.2592,3.6806) (2.2731,3.6705) (2.287,3.6609) (2.3008,3.6515) (2.3147,3.6424) (2.3286,3.6337) (2.3424,3.6253) (2.3563,3.6171) (2.3702,3.6092) (2.384,3.6016) (2.3979,3.5942) (2.4118,3.587) (2.4256,3.5801) (2.4395,3.5734) (2.4534,3.5669) (2.4672,3.5606) (2.4811,3.5545) (2.495,3.5486) (2.5088,3.5429) (2.5227,3.5374) (2.5366,3.532) (2.5504,3.5268) (2.5643,3.5217) (2.5782,3.5168) (2.592,3.512) (2.6059,3.5074) (2.6197,3.5029) (2.6336,3.4985) (2.6475,3.4942) (2.6613,3.4901) (2.6752,3.4861) (2.6891,3.4822) (2.7029,3.4784) (2.7168,3.4747) (2.7307,3.4711) (2.7445,3.4676) (2.7584,3.4642) (2.7723,3.4609) (2.7861,3.4577) (2.8,3.4546)};
\addplot[cmuted,dashed] coordinates {(2.4,3) (2.4,12)};
\node[cmuted,font=\scriptsize,anchor=south west] at (axis cs:2.42,8) {$H_{sep}=2{,}4$};
\nextgroupplot[title={entraînement $F(H_1)$},xlabel={$H_1$},ylabel={$F=\frac{1}{U_e}\frac{d(U_e\theta H_1)}{ds}$}]
\addplot[corange,very thick] coordinates {(3.35,0.058477) (3.4227,0.052051) (3.4954,0.047197) (3.5681,0.043374) (3.6408,0.040269) (3.7134,0.037686) (3.7861,0.035497) (3.8588,0.033612) (3.9315,0.031969) (4.0042,0.030521) (4.0769,0.029233) (4.1496,0.028079) (4.2223,0.027036) (4.295,0.02609) (4.3676,0.025225) (4.4403,0.024432) (4.513,0.023701) (4.5857,0.023025) (4.6584,0.022397) (4.7311,0.021812) (4.8038,0.021266) (4.8765,0.020754) (4.9492,0.020273) (5.0218,0.01982) (5.0945,0.019393) (5.1672,0.018989) (5.2399,0.018606) (5.3126,0.018243) (5.3853,0.017898) (5.458,0.01757) (5.5307,0.017257) (5.6034,0.016958) (5.6761,0.016672) (5.7487,0.016399) (5.8214,0.016137) (5.8941,0.015886) (5.9668,0.015645) (6.0395,0.015413) (6.1122,0.01519) (6.1849,0.014975) (6.2576,0.014768) (6.3303,0.014568) (6.4029,0.014375) (6.4756,0.014189) (6.5483,0.014009) (6.621,0.013835) (6.6937,0.013666) (6.7664,0.013503) (6.8391,0.013345) (6.9118,0.013191) (6.9845,0.013042) (7.0571,0.012898) (7.1298,0.012757) (7.2025,0.012621) (7.2752,0.012488) (7.3479,0.012359) (7.4206,0.012233) (7.4933,0.01211) (7.566,0.011991) (7.6387,0.011875) (7.7113,0.011761) (7.784,0.011651) (7.8567,0.011543) (7.9294,0.011438) (8.0021,0.011335) (8.0748,0.011234) (8.1475,0.011136) (8.2202,0.01104) (8.2929,0.010947) (8.3655,0.010855) (8.4382,0.010765) (8.5109,0.010677) (8.5836,0.010591) (8.6563,0.010507) (8.729,0.010425) (8.8017,0.010344) (8.8744,0.010265) (8.9471,0.010187) (9.0197,0.010111) (9.0924,0.010037) (9.1651,0.0099634) (9.2378,0.0098916) (9.3105,0.0098211) (9.3832,0.009752) (9.4559,0.0096841) (9.5286,0.0096175) (9.6013,0.009552) (9.6739,0.0094877) (9.7466,0.0094245) (9.8193,0.0093624) (9.892,0.0093013) (9.9647,0.0092413) (10.037,0.0091823) (10.11,0.0091243) (10.183,0.0090672) (10.255,0.0090111) (10.328,0.0089558) (10.401,0.0089015) (10.474,0.008848) (10.546,0.0087953) (10.619,0.0087434) (10.692,0.0086924) (10.764,0.0086421) (10.837,0.0085925) (10.91,0.0085437) (10.982,0.0084957) (11.055,0.0084483) (11.128,0.0084016) (11.2,0.0083556) (11.273,0.0083102) (11.346,0.0082655) (11.418,0.0082214) (11.491,0.0081779) (11.564,0.008135) (11.637,0.0080927) (11.709,0.008051) (11.782,0.0080098) (11.855,0.0079692) (11.927,0.0079291) (12,0.0078895)};
\end{groupplot}
\end{tikzpicture}

\caption{Les fermetures de Head. $H_1$ décroît quand $H$ croît : un profil creusé ($H$ grand) a peu de débit pour son déficit de quantité de mouvement. L'entraînement $F$ décroît avec $H_1$.}\label{fig:head-fermeture}
\end{figure}

\paragraph{Mécanisme du décollement turbulent.} En gradient adverse ($U_e'<0$) :
\begin{enumerate}
\item le terme $-(H+2)\theta U_e'/U_e$ devient positif et $\theta$ croît vite ;
\item $E$ croît plus lentement que $U_e\theta$, donc $H_1$ diminue et $H$ augmente : le profil se creuse ;
\item $c_f$ diminue avec $H$ (facteur $10^{-0{,}678H}$) ;
\item quand $H$ atteint 2 à 2,8, le frottement s'annule : c'est le décollement.
\end{enumerate}
Le seuil est flou dans la littérature. Le code prend $H_{sep}=2{,}4$ (constante \code{H\_SEP\_TURB}).

\section{Conditions initiales à la transition}
À la transition, on suppose $\theta$ continu (la quantité de mouvement ne saute pas) et on impose $H=1{,}4$, valeur typique d'une couche turbulente jeune. Le facteur de forme saute donc d'environ 2,6 (laminaire) à 1,4 (turbulent) : c'est le saut visible sur la figure~\ref{fig:bl-2412}. On démarre en $s_{tr}$ interpolé, avec $E_0=U_e(s_{tr})\,\theta(s_{tr})\,H_1(1{,}4)$, où $H_1(1{,}4)=7{,}18$.

\begin{codebox}[ : la fonction \code{rhs} et son intégration]
\begin{lstlisting}
def rhs(si, y):
    th, E = y                          # E = Ue * theta * H1
    ue = Ue_f(si)                      # interpolation lineaire de Ue(s)
    H1 = max(E / (ue * th), H1_sep)    # au-dela du decollement : H gele a 2.4
    Hh = float(H_of_H1(H1))
    cft = float(cf_ludwieg_tillmann(Hh, ue * th * Re))
    dth = 0.5 * cft - (Hh + 2) * th / ue * dUe_f(si)
    dE = ue * 0.0306 * (H1 - 3.0) ** -0.6169
    return [dth, dE]

sol = solve_ivp(rhs, (s_tr, st[-1]), y0, t_eval=st, method="RK45", rtol=1e-6, atol=1e-12)
\end{lstlisting}
Deux protections rendent l'intégration robuste. D'abord, $H_1$ est borné par le bas à $H_1(H_{sep})$ : au-delà du décollement, $H$ est gelé à 2,4 au lieu de diverger, et le drapeau \code{turbulent\_separation} est levé avec l'abscisse $x_{sep}$. Ensuite, \code{H\_of\_H1} borne $H_1$ au-dessus de 3,3, car l'inversion de $H_1(H)$ est singulière en $H_1\to3{,}3$.
\end{codebox}

\section{Le résultat sur un profil}
\begin{figure}[htbp]
\centering
% figure générée par tools/make_figures.py : cl_2412
\begin{tikzpicture}
\begin{groupplot}[group style={group size=3 by 1,horizontal sep=1.2cm},width=0.35\linewidth,height=5cm,
  xlabel={$x/c$},xmin=0,xmax=1,title style={font=\small},tick label style={font=\scriptsize},
  legend style={draw=none,font=\scriptsize}]
\nextgroupplot[title={$10^3\,\theta/c$},legend pos=north west]
\addplot[cblue,very thick] coordinates {(0.0038514,0.025409) (0.0060621,0.028835) (0.0087685,0.03276) (0.011969,0.037061) (0.01566,0.041645) (0.01984,0.046438) (0.024505,0.051386) (0.029651,0.056445) (0.035273,0.061585) (0.041368,0.066783) (0.047928,0.072022) (0.054949,0.077292) (0.062423,0.082584) (0.070344,0.087895) (0.078703,0.093223) (0.087494,0.098568) (0.096707,0.10393) (0.10633,0.10932) (0.11636,0.11473) (0.12679,0.12018) (0.13759,0.12565) (0.14877,0.13118) (0.16031,0.13675) (0.17219,0.14237) (0.18442,0.14805) (0.19696,0.1538) (0.20982,0.15963) (0.22297,0.16553) (0.2364,0.17153) (0.25011,0.17762) (0.26407,0.18382) (0.27826,0.19014) (0.29268,0.22309) (0.30731,0.25758) (0.32214,0.29173) (0.33714,0.32588) (0.3523,0.3603) (0.36761,0.39523) (0.38305,0.43102) (0.39857,0.46781) (0.41414,0.50511) (0.42974,0.54227) (0.44542,0.57929) (0.46115,0.61635) (0.47691,0.6536) (0.49269,0.69108) (0.50848,0.72884) (0.52425,0.7669) (0.54,0.8053) (0.5557,0.84403) (0.57135,0.88312) (0.58692,0.92258) (0.6024,0.96242) (0.61777,1.0027) (0.63302,1.0433) (0.64814,1.0844) (0.6631,1.126) (0.6779,1.1681) (0.69252,1.2107) (0.70695,1.2538) (0.72116,1.2976) (0.73515,1.3421) (0.74891,1.3873) (0.76242,1.4333) (0.77566,1.4802) (0.78863,1.5281) (0.80131,1.5771) (0.81369,1.6272) (0.82576,1.6787) (0.8375,1.7316) (0.84891,1.7861) (0.85997,1.8425) (0.87068,1.9009) (0.88102,1.9615) (0.89098,2.0247) (0.90055,2.0906) (0.90973,2.1599) (0.9185,2.2327) (0.92686,2.3098) (0.9348,2.3917) (0.94232,2.4792) (0.9494,2.5733) (0.95603,2.6751) (0.96222,2.786) (0.96795,2.908) (0.97323,3.0436) (0.97804,3.1965) (0.98238,3.3621)};\addlegendentry{extrados}
\addplot[corange,very thick] coordinates {(0.0040239,0.025927) (0.006247,0.025829) (0.0089501,0.028018) (0.01213,0.031211) (0.015782,0.034992) (0.019902,0.039175) (0.024485,0.04366) (0.029527,0.048382) (0.035021,0.053298) (0.040961,0.058375) (0.047342,0.063589) (0.054155,0.068922) (0.061395,0.074361) (0.069053,0.079893) (0.077121,0.085512) (0.085592,0.091209) (0.094457,0.09698) (0.10371,0.10282) (0.11333,0.10872) (0.12333,0.11469) (0.13367,0.12072) (0.14437,0.1268) (0.1554,0.13294) (0.16676,0.13912) (0.17844,0.14535) (0.19042,0.15162) (0.20269,0.15793) (0.21525,0.16427) (0.22808,0.17063) (0.24116,0.17701) (0.25449,0.1834) (0.26806,0.18979) (0.28185,0.19616) (0.29585,0.20251) (0.31005,0.20882) (0.32443,0.21507) (0.33899,0.22123) (0.3537,0.22728) (0.36855,0.23316) (0.38354,0.2388) (0.39866,0.244) (0.41396,0.24893) (0.42937,0.2542) (0.44486,0.25975) (0.4604,0.26543) (0.47599,0.27118) (0.4916,0.30364) (0.50723,0.33927) (0.52285,0.37397) (0.53845,0.40795) (0.55402,0.44136) (0.56954,0.47429) (0.58499,0.50677) (0.60037,0.53886) (0.61565,0.57057) (0.63082,0.60192) (0.64587,0.63291) (0.66078,0.66356) (0.67553,0.69387) (0.69012,0.72383) (0.70452,0.75346) (0.71872,0.78276) (0.73272,0.81175) (0.74649,0.84042) (0.76001,0.86881) (0.77329,0.89693) (0.7863,0.92478) (0.79903,0.9524) (0.81147,0.97984) (0.82361,1.0071) (0.83542,1.0343) (0.84691,1.0614) (0.85806,1.0885) (0.86886,1.1157) (0.8793,1.143) (0.88936,1.1706) (0.89904,1.1985) (0.90832,1.2269) (0.9172,1.2559) (0.92567,1.2857) (0.93372,1.3166) (0.94134,1.3488) (0.94852,1.3826) (0.95526,1.4184) (0.96155,1.4568) (0.96737,1.4984) (0.97274,1.5441) (0.97763,1.5953) (0.98205,1.6499)};\addlegendentry{intrados}
\nextgroupplot[title={facteur de forme $H$},ymin=1,ymax=3]
\addplot[cblue,very thick] coordinates {(0.0038514,2.4339) (0.0060621,2.4605) (0.0087685,2.4851) (0.011969,2.507) (0.01566,2.526) (0.01984,2.5422) (0.024505,2.5557) (0.029651,2.5671) (0.035273,2.5765) (0.041368,2.5845) (0.047928,2.5913) (0.054949,2.5972) (0.062423,2.6024) (0.070344,2.6073) (0.078703,2.6121) (0.087494,2.6167) (0.096707,2.6214) (0.10633,2.6263) (0.11636,2.6316) (0.12679,2.6373) (0.13759,2.6436) (0.14877,2.6505) (0.16031,2.6582) (0.17219,2.6667) (0.18442,2.6763) (0.19696,2.687) (0.20982,2.6991) (0.22297,2.7128) (0.2364,2.7284) (0.25011,2.7462) (0.26407,2.7668) (0.27826,2.7907) (0.29268,1.4161) (0.30731,1.4263) (0.32214,1.4321) (0.33714,1.4354) (0.3523,1.4375) (0.36761,1.4389) (0.38305,1.4403) (0.39857,1.4417) (0.41414,1.443) (0.42974,1.4438) (0.44542,1.4441) (0.46115,1.4442) (0.47691,1.4442) (0.49269,1.4443) (0.50848,1.4444) (0.52425,1.4446) (0.54,1.4448) (0.5557,1.4452) (0.57135,1.4457) (0.58692,1.4463) (0.6024,1.4471) (0.61777,1.4479) (0.63302,1.4489) (0.64814,1.4501) (0.6631,1.4514) (0.6779,1.4528) (0.69252,1.4544) (0.70695,1.4562) (0.72116,1.4582) (0.73515,1.4604) (0.74891,1.4629) (0.76242,1.4655) (0.77566,1.4684) (0.78863,1.4716) (0.80131,1.4752) (0.81369,1.4791) (0.82576,1.4833) (0.8375,1.488) (0.84891,1.4931) (0.85997,1.4988) (0.87068,1.5051) (0.88102,1.512) (0.89098,1.5197) (0.90055,1.5282) (0.90973,1.5378) (0.9185,1.5485) (0.92686,1.5607) (0.9348,1.5744) (0.94232,1.5903) (0.9494,1.6032) (0.95603,1.6194) (0.96222,1.6381) (0.96795,1.66) (0.97323,1.6862) (0.97804,1.7182) (0.98238,1.7548)};
\addplot[corange,very thick] coordinates {(0.0040239,2.2874) (0.006247,2.3541) (0.0089501,2.3887) (0.01213,2.412) (0.015782,2.431) (0.019902,2.4477) (0.024485,2.4628) (0.029527,2.4767) (0.035021,2.4894) (0.040961,2.5012) (0.047342,2.5121) (0.054155,2.5223) (0.061395,2.5319) (0.069053,2.541) (0.077121,2.5497) (0.085592,2.5581) (0.094457,2.5662) (0.10371,2.5741) (0.11333,2.5818) (0.12333,2.5894) (0.13367,2.5969) (0.14437,2.6042) (0.1554,2.6117) (0.16676,2.6189) (0.17844,2.6262) (0.19042,2.6336) (0.20269,2.6409) (0.21525,2.6481) (0.22808,2.6552) (0.24116,2.6621) (0.25449,2.6686) (0.26806,2.6746) (0.28185,2.6799) (0.29585,2.6844) (0.31005,2.6876) (0.32443,2.6893) (0.33899,2.6887) (0.3537,2.6851) (0.36855,2.6762) (0.38354,2.657) (0.39866,2.6364) (0.41396,2.6416) (0.42937,2.6693) (0.44486,2.6931) (0.4604,2.7101) (0.47599,2.7245) (0.4916,1.4106) (0.50723,1.4185) (0.52285,1.4234) (0.53845,1.4263) (0.55402,1.428) (0.56954,1.4289) (0.58499,1.4292) (0.60037,1.4293) (0.61565,1.4291) (0.63082,1.4287) (0.64587,1.4283) (0.66078,1.4279) (0.67553,1.4274) (0.69012,1.427) (0.70452,1.4266) (0.71872,1.4262) (0.73272,1.4259) (0.74649,1.4257) (0.76001,1.4255) (0.77329,1.4254) (0.7863,1.4255) (0.79903,1.4256) (0.81147,1.4258) (0.82361,1.4262) (0.83542,1.4268) (0.84691,1.4275) (0.85806,1.4283) (0.86886,1.4294) (0.8793,1.4307) (0.88936,1.4322) (0.89904,1.434) (0.90832,1.4361) (0.9172,1.4386) (0.92567,1.4415) (0.93372,1.4448) (0.94134,1.4487) (0.94852,1.4533) (0.95526,1.4587) (0.96155,1.4649) (0.96737,1.4724) (0.97274,1.4814) (0.97763,1.4924) (0.98205,1.5045)};
\nextgroupplot[title={$10^3\,c_f$},ymin=0,ymax=8]
\addplot[cblue,very thick] coordinates {(0.0038514,6.941) (0.0060621,5.5144) (0.0087685,4.4794) (0.011969,3.7164) (0.01566,3.1445) (0.01984,2.7081) (0.024505,2.369) (0.029651,2.1007) (0.035273,1.8848) (0.041368,1.708) (0.047928,1.5609) (0.054949,1.4367) (0.062423,1.3304) (0.070344,1.2382) (0.078703,1.1571) (0.087494,1.0852) (0.096707,1.0206) (0.10633,0.96212) (0.11636,0.90864) (0.12679,0.85934) (0.13759,0.81355) (0.14877,0.77073) (0.16031,0.73041) (0.17219,0.69221) (0.18442,0.6558) (0.19696,0.62089) (0.20982,0.58722) (0.22297,0.55455) (0.2364,0.52266) (0.25011,0.49131) (0.26407,0.46025) (0.27826,0.42919) (0.29268,4.3902) (0.30731,4.1622) (0.32214,3.994) (0.33714,3.8616) (0.3523,3.7517) (0.36761,3.6564) (0.38305,3.5699) (0.39857,3.4905) (0.41414,3.4186) (0.42974,3.3554) (0.44542,3.2999) (0.46115,3.2497) (0.47691,3.2034) (0.49269,3.1602) (0.50848,3.1193) (0.52425,3.0805) (0.54,3.0434) (0.5557,3.0077) (0.57135,2.9734) (0.58692,2.9401) (0.6024,2.9077) (0.61777,2.8761) (0.63302,2.8452) (0.64814,2.8149) (0.6631,2.7851) (0.6779,2.7556) (0.69252,2.7263) (0.70695,2.6972) (0.72116,2.6682) (0.73515,2.6391) (0.74891,2.6099) (0.76242,2.5805) (0.77566,2.5507) (0.78863,2.5204) (0.80131,2.4895) (0.81369,2.4579) (0.82576,2.4255) (0.8375,2.392) (0.84891,2.3574) (0.85997,2.3215) (0.87068,2.2841) (0.88102,2.2449) (0.89098,2.2037) (0.90055,2.1603) (0.90973,2.1142) (0.9185,2.0652) (0.92686,2.0126) (0.9348,1.9561) (0.94232,1.8947) (0.9494,1.8431) (0.95603,1.7833) (0.96222,1.7181) (0.96795,1.6462) (0.97323,1.5661) (0.97804,1.4755) (0.98238,1.3797)};
\addplot[corange,very thick] coordinates {(0.0040239,25.307) (0.006247,16.678) (0.0089501,11.981) (0.01213,9.112) (0.015782,7.2118) (0.019902,5.8823) (0.024485,4.9133) (0.029527,4.1837) (0.035021,3.6194) (0.040961,3.173) (0.047342,2.813) (0.054155,2.5176) (0.061395,2.2717) (0.069053,2.0642) (0.077121,1.8871) (0.085592,1.7343) (0.094457,1.6013) (0.10371,1.4845) (0.11333,1.3811) (0.12333,1.2891) (0.13367,1.2066) (0.14437,1.1324) (0.1554,1.0652) (0.16676,1.0042) (0.17844,0.94858) (0.19042,0.89781) (0.20269,0.85136) (0.21525,0.80883) (0.22808,0.76987) (0.24116,0.73423) (0.25449,0.7017) (0.26806,0.67213) (0.28185,0.64542) (0.29585,0.62153) (0.31005,0.60051) (0.32443,0.58251) (0.33899,0.56787) (0.3537,0.55731) (0.36855,0.55263) (0.38354,0.55915) (0.39866,0.56867) (0.41396,0.55229) (0.42937,0.51493) (0.44486,0.4839) (0.4604,0.46051) (0.47599,0.44056) (0.4916,4.3816) (0.50723,4.2031) (0.52285,4.0659) (0.53845,3.9562) (0.55402,3.8655) (0.56954,3.7884) (0.58499,3.7217) (0.60037,3.6628) (0.61565,3.6102) (0.63082,3.5626) (0.64587,3.5192) (0.66078,3.4792) (0.67553,3.4421) (0.69012,3.4075) (0.70452,3.3751) (0.71872,3.3445) (0.73272,3.3155) (0.74649,3.2878) (0.76001,3.2613) (0.77329,3.2358) (0.7863,3.2112) (0.79903,3.1872) (0.81147,3.1638) (0.82361,3.1407) (0.83542,3.1178) (0.84691,3.095) (0.85806,3.0722) (0.86886,3.0492) (0.8793,3.0258) (0.88936,3.0018) (0.89904,2.9771) (0.90832,2.9515) (0.9172,2.9247) (0.92567,2.8965) (0.93372,2.8664) (0.94134,2.834) (0.94852,2.799) (0.95526,2.7607) (0.96155,2.7185) (0.96737,2.6714) (0.97274,2.6181) (0.97763,2.5566) (0.98205,2.4924)};
\end{groupplot}
\end{tikzpicture}

\caption{Couche limite complète du NACA 2412 à $\alpha=2^\circ$, $\Re=3\cdot10^6$, calculée par \code{march\_surface}. On voit la transition (saut de $H$ et de $c_f$) en $x/c=0{,}28$ sur l'extrados et $0{,}48$ sur l'intrados, puis la montée de $H$ en fin d'extrados sous l'effet du gradient adverse.}\label{fig:bl-2412}
\end{figure}

La validation sur plaque plane turbulente (figure~\ref{fig:plaque}, à droite) donne $\theta$ à moins de 5~\% de la loi en $1/7$, $\theta=0{,}036\,x\,\Re_x^{-1/5}$, et $H\to1{,}34$.

\begin{piege}
\textbf{Le faux décollement de bord de fuite.} Près d'un bord de fuite fermé, l'écoulement potentiel décélère fortement (ch.~\ref{ch:trainee}), et ce gradient adverse est exagéré par rapport à l'écoulement réel. Sur la base de 512 profils à $\alpha=2^\circ$, 145 profils épais ou cambrés atteignent $H>2{,}4$ entre $x/c=0{,}957$ et $0{,}983$. Les déclarer invalides retirerait 28~\% de l'espace de conception à cause d'un artefact local. Le code les tolère (seuil \code{x\_sep\_crit}, ch.~\ref{ch:modele}).
\end{piege}

\begin{piege}
Head est une méthode « d'équilibre » : $F$ ne dépend que de la forme locale. Elle réagit donc instantanément, donc mal, aux changements brusques (choc, recollement de bulbe). Les méthodes à retard (\emph{lag-entrainment} de Green, ou le \emph{shear-lag} d'XFOIL) ajoutent une équation pour la contrainte de cisaillement et corrigent ce défaut.
\end{piege}

\exercices
\begin{exercice}
Calculer $H_1$ pour $H=1{,}3$, $1{,}4$ et $2{,}4$ avec les corrélations de Head. Que vaut $F$ pour $H=1{,}4$ ? Pour une couche limite d'épaisseur de quantité de mouvement $\theta=10^{-3}c$, quelle est alors l'épaisseur $\delta-\delta^*$ ?
\end{exercice}
\begin{exercice}
Montrer, à partir du système de Head et du sens de variation des fermetures, qu'en gradient adverse $H$ augmente. On pourra raisonner à $\theta$ fixé sur une petite distance et regarder comment évolue $H_1=E/(U_e\theta)$ quand $U_e$ diminue.
\end{exercice}

\chapter{La traînée : formule de Squire-Young}\label{ch:trainee}
\iterbanner{V0}

\section{La traînée est dans le sillage}
L'écoulement potentiel ne donne aucune traînée (d'Alembert). Toute la traînée de profil vient donc de la viscosité : le frottement à la paroi, et la \emph{traînée de forme}, qui est l'effet de l'épaississement de la couche limite sur la distribution de pression. Notre modèle ne calcule pas cette seconde part directement, puisque la pression vient de l'écoulement potentiel sans rétroaction. Il faut un autre chemin.

\paragraph{Le bilan de quantité de mouvement.} Prenons un grand volume de contrôle autour du profil, fermé loin en aval là où la pression est revenue à $p_\infty$ (figure~\ref{fig:sillage}). Le bilan de quantité de mouvement selon $x$ dit que la traînée totale est égale au défaut de flux de quantité de mouvement dans le sillage :
\[
D'=\rho\int_{\text{sillage}}u\,(U_\infty-u)\,\dd y=\rho\,U_\infty^2\,\theta_\infty
\quad\Longrightarrow\quad
C_D=\frac{2\,\theta_\infty}{c}.
\]
Ce résultat contient automatiquement le frottement \emph{et} la traînée de forme, sans avoir à les séparer. C'est le principe de la mesure de traînée par sondage de sillage en soufflerie. Il reste à relier $\theta_\infty$ à ce que l'on connaît, l'état de la couche limite au bord de fuite.

\begin{figure}[htbp]
\centering
\begin{tikzpicture}[x=1.25cm,y=1.25cm,scale=1.05]
\fill[cmuted!25,draw=cink,thick] plot[smooth,tension=0.6] coordinates {(4,0) (3,0.2) (1.8,0.36) (0.8,0.36) (0.2,0.2) (0,0) (0.2,-0.15) (0.8,-0.2) (1.8,-0.16) (3,-0.06) (4,0)};
\fill[cblue!12] (4,0.12) .. controls (6,0.2) and (8,0.35) .. (11,0.5) -- (11,-0.5) .. controls (8,-0.35) and (6,-0.2) .. (4,-0.1) -- cycle;
\draw[cblue,dashed] (4,0.12) .. controls (6,0.2) and (8,0.35) .. (11,0.5);
\draw[cblue,dashed] (4,-0.1) .. controls (6,-0.2) and (8,-0.35) .. (11,-0.5);
\draw[cmuted,dashed] (-1,-1.6) rectangle (11,1.6);
\node[cmuted,font=\small,anchor=south west] at (-1,1.6) {volume de contrôle};
% profil de vitesse au BF et à l'infini
\draw[cink,thin] (4,-0.6) -- (4,0.6);
\draw[corange,thick] plot[smooth] coordinates {(4.0,-0.6) (4.6,-0.25) (4.15,-0.1) (4.0,0) (4.2,0.12) (4.6,0.3) (4.65,0.6)};
\node[corange,font=\small,align=center] at (4.5,-1.15) {bord de fuite :\\ $\theta_{TE}$, $H_{TE}$, $U_{e,TE}<U_\infty$};
\draw[cink,thin] (10.2,-0.9) -- (10.2,0.9);
\draw[corange,thick] plot[smooth] coordinates {(10.75,-0.9) (10.7,-0.5) (10.5,-0.2) (10.4,0) (10.5,0.2) (10.7,0.5) (10.75,0.9)};
\node[corange,font=\small,align=center] at (8.6,-1.15) {loin en aval :\\ $\theta_\infty$, $H\to1$, $U_e=U_\infty$};
\draw[-{Stealth[length=2.5mm]},cink,thick] (-0.9,0.9) -- (0,0.9) node[midway,above,font=\small] {$U_\infty$};
\end{tikzpicture}
\caption{La traînée se lit dans le sillage lointain : $C_D=2\theta_\infty/c$. Squire et Young relient $\theta_\infty$ à l'état de la couche limite au bord de fuite.}\label{fig:sillage}
\end{figure}

\section{De $\theta_{TE}$ à $\theta_\infty$ : l'hypothèse de Squire et Young}
Dans le sillage, il n'y a plus de paroi, donc $c_f=0$. L'équation de von Kármán, appliquée à chaque moitié du sillage, devient
\[
\frac{\dd\theta}{\dd s}=-(H+2)\frac{\theta}{U_e}\frac{\dd U_e}{\dd s}
\quad\Longleftrightarrow\quad
\frac{\dd\ln\theta}{\dd\ln U_e}=-(H+2).
\]
Entre le bord de fuite et l'infini aval, $U_e$ remonte de $U_{e,TE}$ à $U_\infty$ et le sillage se « remplit » : $H$ descend de $H_{TE}$ à 1 (profil uniforme). Squire et Young (1938) supposent que $H$ varie \emph{linéairement en $\ln U_e$} entre ces deux états :
\[
H=1+\big(H_{TE}-1\big)\,\frac{\ln(U_\infty/U_e)}{\ln(U_\infty/U_{e,TE})}.
\]
Posons $u=\ln(U_e/U_\infty)$, qui va de $u_{TE}<0$ à 0, et $\kappa=(H_{TE}-1)/u_{TE}$, de sorte que $H=1+\kappa u$. L'intégration donne
\[
\ln\frac{\theta_\infty}{\theta_{TE}}=-\int_{u_{TE}}^0\big(3+\kappa u\big)\,\dd u
=3u_{TE}+\frac{\kappa}{2}u_{TE}^2
=u_{TE}\Big(3+\frac{H_{TE}-1}{2}\Big)=\frac{H_{TE}+5}{2}\,\ln\frac{U_{e,TE}}{U_\infty}.
\]
D'où la \textbf{formule de Squire-Young}, sommée sur les deux faces :
\[
\boxed{\ C_D=\sum_{\text{faces}}2\,\frac{\theta_{TE}}{c}\left(\frac{U_{e,TE}}{U_\infty}\right)^{\frac{H_{TE}+5}{2}}\ }
\]
Comme $U_{e,TE}<U_\infty$, le facteur est inférieur à 1 : le sillage s'accélère en aval du bord de fuite, ce qui réduit $\theta$. L'exposant croît avec $H_{TE}$ : une couche épaisse et creusée au bord de fuite « perd » davantage en se remplissant.

\begin{codebox}[ : \code{squire\_young} et la décomposition de $C_D$]
\begin{lstlisting}
def squire_young(theta_te, H_te, Ue_te):
    return 2.0 * theta_te * Ue_te ** ((H_te + 5.0) / 2.0)

CD = (squire_young(bl_up.theta_te, bl_up.H_te, bl_up.Ue_te)
      + squire_young(bl_lo.theta_te, bl_lo.H_te, bl_lo.Ue_te))
CD_f = (_friction_drag(bl_up, tx_up) + _friction_drag(bl_lo, tx_lo))   # sum cf Ue^2 |tx| ds
\end{lstlisting}
\code{CD\_f} intègre le frottement le long de la paroi, projeté sur la direction de la corde : $C_{D,f}=\sum c_f\,U_e^2\,|t_x|\,\Delta s$. La différence $C_D-C_{D,f}$ est la traînée de forme. Pour le NACA 2412 à $2^\circ$ : $C_D=69$~counts, dont 55 de frottement et 14 de forme.
\end{codebox}

\section{L'artefact du bord de fuite fermé et \code{x\_stop}}
\paragraph{Le problème.} Un bord de fuite fermé d'angle non nul est un \textbf{point d'arrêt} de l'écoulement potentiel : la vitesse y tombe à zéro. Sur le dernier pour-cent de corde, $U_e$ chute brutalement (figure~\ref{fig:bf}). Pour le NACA 0012 à $\alpha=0$, elle passe de 0,88 en $x=0{,}985$ à 0,70 sur la dernière station, et $\theta$ y triple sous l'effet de ce gradient adverse violent. La quantité $\theta\,U_e^{(H+5)/2}$, censée être à peu près conservée en aval du point où l'on applique Squire-Young, ne l'est plus, et $C_D$ est surestimé d'environ 22~\%.

Cette décélération n'existe pas dans l'écoulement réel : l'épaisseur de déplacement des deux couches limites et le sillage « ouvrent » le bord de fuite, et l'écoulement extérieur ne voit jamais de point d'arrêt. C'est exactement ce que corrige un couplage fort comme celui d'XFOIL.

\paragraph{La correction.} On arrête la marche de couche limite en $x_{stop}=0{,}985$ et on applique Squire-Young là, juste avant la zone artificielle. La figure~\ref{fig:bf} montre que la sensibilité au choix de $x_{stop}$ est faible tant qu'on reste en dehors de cette zone : au plus $\pm1{,}5$~\% sur $C_D$ entre 0,97 et 0,99, contre un saut de +22~\% (0012) ou +26~\% (2412) si l'on va jusqu'au bout.

\begin{figure}[htbp]
\centering
% figure générée par tools/make_figures.py : bord_de_fuite
\begin{tikzpicture}
\begin{groupplot}[group style={group size=3 by 1,horizontal sep=1.2cm},width=0.35\linewidth,height=5cm,
  title style={font=\small},tick label style={font=\scriptsize},legend style={draw=none,font=\scriptsize}]
\nextgroupplot[title={$U_e/U_\infty$ (NACA 0012, $\alpha=0$)},xlabel={$x/c$},xmin=0.85,xmax=1]
\addplot[cblue,very thick,mark=*,mark size=1pt] coordinates {(0.85902,1.0004) (0.86977,0.99553) (0.88016,0.99049) (0.89017,0.98529) (0.89979,0.97992) (0.90902,0.97435) (0.91785,0.96858) (0.92627,0.96256) (0.93426,0.95629) (0.94183,0.94973) (0.94896,0.94284) (0.95565,0.93558) (0.96188,0.92789) (0.96766,0.91972) (0.97298,0.91098) (0.97784,0.90156) (0.98222,0.89131) (0.98612,0.88006) (0.98955,0.86751) (0.99249,0.85326) (0.99495,0.83669) (0.99692,0.81672) (0.9984,0.79134) (0.99938,0.75617) (0.99988,0.70067)};
\addplot[cmuted,dashed] coordinates {(0.985,0.6) (0.985,1.0)};
\node[cmuted,font=\scriptsize,anchor=south east,rotate=90] at (axis cs:0.985,0.62) {$x_{stop}$};
\nextgroupplot[title={$10^3\,\theta/c$},xlabel={$x/c$},xmin=0.85,xmax=1]
\addplot[corange,very thick,mark=*,mark size=1pt] coordinates {(0.85902,1.4903) (0.86977,1.5303) (0.88016,1.5713) (0.89017,1.6135) (0.89979,1.6571) (0.90902,1.7022) (0.91785,1.7491) (0.92627,1.7981) (0.93426,1.8497) (0.94183,1.9042) (0.94896,1.9622) (0.95565,2.0246) (0.96188,2.0921) (0.96766,2.1658) (0.97298,2.2475) (0.97784,2.3392) (0.98222,2.444) (0.98612,2.5661) (0.98955,2.7125) (0.99249,2.8951) (0.99495,3.1352) (0.99692,3.4815) (0.9984,4.0841) (0.99938,5.6476) (0.99988,7.6063)};
\addplot[cmuted,dashed] coordinates {(0.985,1) (0.985,4)};
\nextgroupplot[title={$C_D$ (counts) selon $x_{stop}$},xlabel={$x_{stop}$},xmin=0.95,xmax=1,legend pos=north west]
\addplot[cblue,very thick,mark=*,mark size=1.3pt] coordinates {(0.95,64.755) (0.96,65.116) (0.97,65.784) (0.975,66.091) (0.98,66.382) (0.985,66.681) (0.99,67.269) (0.995,67.87) (1,81.584)};\addlegendentry{0012, $\alpha=0$}
\addplot[corange,very thick,mark=square*,mark size=1.3pt] coordinates {(0.95,66.989) (0.96,67.393) (0.97,68.147) (0.975,68.501) (0.98,68.842) (0.985,69.17) (0.99,69.774) (0.995,69.962) (1,86.997)};\addlegendentry{2412, $\alpha=2^\circ$}
\end{groupplot}
\end{tikzpicture}

\caption{L'artefact de bord de fuite. À gauche et au centre, $U_e$ et $\theta$ sur l'extrados du NACA 0012 à $\alpha=0$, sans coupure : chute de vitesse et explosion de $\theta$ sur le dernier pour-cent. À droite, $C_D$ selon l'abscisse d'arrêt : plateau jusqu'à 0,995, puis saut.}\label{fig:bf}
\end{figure}

\begin{codebox}[ : la coupure en tête de \code{march\_surface}]
\begin{lstlisting}
keep = x <= x_stop
keep[: np.argmax(x > 0.5)] = True     # ne coupe jamais cote bord d'attaque
s, x, Ue = s[keep], x[keep], Ue[keep]
\end{lstlisting}
La deuxième ligne protège le bord d'attaque. Sur l'intrados d'un profil cambré, les premières stations peuvent avoir un $x$ légèrement négatif ou décroissant près du point d'arrêt, et ne doivent jamais être retirées.
\end{codebox}

\begin{piege}
Squire-Young suppose un sillage qui se remplit régulièrement. Avec un décollement étendu, cette hypothèse tombe, et la formule n'a plus de sens. C'est une des raisons pour lesquelles un décollement turbulent en amont de $x/c=0{,}9$ rend le résultat invalide (ch.~\ref{ch:modele}).
\end{piege}

\begin{resultat}
\begin{itemize}
\item $C_D=2\theta_\infty$ par bilan de quantité de mouvement : frottement et forme sont inclus.
\item Squire-Young transporte $\theta$ du bord de fuite à l'infini avec l'exposant $(H_{TE}+5)/2$.
\item On l'applique en $x_{stop}=0{,}985$, pour éviter le point d'arrêt artificiel du bord de fuite fermé. Sensibilité : au plus $\pm1{,}5$~\%.
\end{itemize}
\end{resultat}

\exercices
\begin{exercice}
Pour le NACA 2412 à $\alpha=2^\circ$ et $\Re=3\cdot10^6$, le code obtient en $x_{stop}$ : extrados $\theta=3{,}362\cdot10^{-3}$, $H=1{,}755$, $U_e=0{,}905$ ; intrados $\theta=1{,}650\cdot10^{-3}$, $H=1{,}504$, $U_e=0{,}873$. Calculer $C_D$ en counts, et la contribution de chaque face.
\end{exercice}
\begin{exercice}
Refaire la démonstration de l'exposant en supposant cette fois que $H$ reste constant, égal à $H_{TE}$, dans tout le sillage. Quel exposant obtient-on ? Pourquoi cette hypothèse est-elle physiquement moins bonne ?
\end{exercice}

\chapter{L'assemblage : \code{evaluate} et \code{AeroResult}}\label{ch:modele}
\iterbanner{V0}

\section{La fonction \code{evaluate}}
C'est la seule fonction de la phase~1 que les autres phases appellent. Elle enchaîne tout ce qui précède :
\begin{lstlisting}
def evaluate(m, p, t, alpha=0.0, Re=3e6, n_half=100, x_trip=(None, None),
             x_stop=0.985, free_transition=True, x_sep_crit=0.9, return_details=False):
    X, Y = naca4(m, p, t, n_half=n_half)                     # ch. 2
    inv = solve(X, Y, alpha)                                 # ch. 4
    up, lo = split_at_stagnation(inv)                        # ch. 5 (exception rattrapee)
    bl_up = march_surface(*up[:3], Re, x_trip=x_trip[0], x_stop=x_stop, ...)   # ch. 6-8
    bl_lo = march_surface(*lo[:3], Re, x_trip=x_trip[1], x_stop=x_stop, ...)
    CD = squire_young(...up...) + squire_young(...lo...)     # ch. 9
    ...
\end{lstlisting}
\begin{center}
\renewcommand{\arraystretch}{1.15}
\begin{tabularx}{\linewidth}{lX}
\toprule
Argument & Rôle \\
\midrule
\code{m, p, t} & variables de conception, en fractions de corde \\
\code{alpha} & incidence en degrés \\
\code{Re} & Reynolds de corde ($3\cdot10^6$ par défaut, à figer en équipe) \\
\code{n\_half} & demi-nombre de panneaux (100, soit 200 panneaux) \\
\code{x\_trip} & transition forcée (extrados, intrados) ; \code{None} = libre \\
\code{free\_transition} & \code{False} désactive Michel, pour imposer une transition tardive \\
\code{x\_stop} & abscisse d'arrêt de la couche limite et d'application de Squire-Young \\
\code{x\_sep\_crit} & un décollement turbulent en amont de cette abscisse invalide le calcul \\
\code{return\_details} & renvoie aussi la solution de panneaux et les deux \code{SurfaceBL} \\
\bottomrule
\end{tabularx}
\end{center}

\section{Le résultat : \code{AeroResult}}
\begin{center}
\renewcommand{\arraystretch}{1.15}
\begin{tabularx}{\linewidth}{lX}
\toprule
Champ & Signification \\
\midrule
\code{CL}, \code{CL\_inv} & portance non visqueuse (identiques : le couplage faible ne corrige pas $C_L$) \\
\code{CD} & traînée de profil, Squire-Young sommé sur les deux faces \\
\code{CD\_f} & part de frottement ; \code{CD - CD\_f} est la traînée de forme \\
\code{CM} & moment de tangage au quart de corde, positif à cabrer \\
\code{x\_tr\_up}, \code{x\_tr\_lo} & positions de transition (interpolées) \\
\code{tr\_mode\_up}, \code{tr\_mode\_lo} & \code{'michel'}, \code{'laminar\_sep'}, \code{'forced'} ou \code{'none'} \\
\code{x\_sep\_up}, \code{x\_sep\_lo} & abscisse de décollement turbulent, \code{nan} si attaché \\
\code{r\_le} & rayon de bord d'attaque $1{,}1019\,t^2$, pour la phase SER \\
\code{valid} & \code{False} si le résultat sort du domaine du modèle \\
\code{flags} & liste de diagnostics, préfixés \code{up:} ou \code{lo:} \\
\bottomrule
\end{tabularx}
\end{center}
\code{to\_dict()} aplatit le tout pour l'écriture en CSV (\code{generer\_base.py}).

\section{La logique de validité}
Un optimiseur génétique explore les coins de l'espace. Le modèle ne doit jamais planter, mais il doit dire quand son résultat n'est plus crédible. La règle est :
\[
\code{valid}=\big(x_{sep}\ge x_{sep,crit}\big)\ \text{et}\ C_D\ \text{fini}\ \text{et aucun drapeau d'échec},
\]
où $x_{sep}$ est le plus amont des deux décollements turbulents ($+\infty$ si aucun). Les drapeaux d'échec sont \code{head\_ivp\_failed} (l'intégrateur n'a pas abouti) et \code{laminar\_separation\_no\_transition}.
\begin{itemize}
\item Un décollement turbulent \textbf{en aval} de $x_{sep,crit}=0{,}9$ est toléré, avec le drapeau \code{te\_separation} : c'est le plus souvent l'artefact de bord de fuite du chapitre~\ref{ch:head}, et $C_D$ reste plausible.
\item Un décollement \textbf{en amont} signifie que le modèle sort de son domaine : couplage faible, Squire-Young et Kutta au bord de fuite n'ont plus de sens. Le profil est déclaré invalide, donc infaisable pour l'optimiseur.
\item Un point d'arrêt ambigu (plusieurs changements de signe de $V_t$) produit un résultat invalide avec $C_D=$\,\code{nan}.
\end{itemize}

\begin{figure}[htbp]
\centering
% figure générée par tools/make_figures.py : polaire
\begin{tikzpicture}
\begin{groupplot}[group style={group size=2 by 1,horizontal sep=1.6cm},width=0.5\linewidth,height=5.6cm,
  title style={font=\small},tick label style={font=\scriptsize},legend style={draw=none,font=\scriptsize}]
\nextgroupplot[title={Polaire $C_L(C_D)$},xlabel={$C_D$ (counts)},ylabel={$C_L$},legend pos=south east]
\addplot[cblue,very thick,mark=*,mark size=1.3pt] coordinates {(69.584,-0.2265) (68.211,-0.10581) (67.478,0.014921) (67.288,0.13565) (67.371,0.25635) (68.1,0.37698) (69.17,0.4975) (72.579,0.61788) (77.134,0.73808) (82.671,0.85806) (90.04,0.97778) (98.424,1.0972) (114.76,1.2163) (129.4,1.3351) (145.75,1.4534)};\addlegendentry{valide}

\nextgroupplot[title={Décomposition de $C_D$},xlabel={$\alpha$ (°)},ylabel={counts},legend pos=north west,ymin=0]
\addplot[cblue,very thick] coordinates {(-4,69.584) (-3,68.211) (-2,67.478) (-1,67.288) (0,67.371) (1,68.1) (2,69.17) (3,72.579) (4,77.134) (5,82.671) (6,90.04) (7,98.424) (8,114.76) (9,129.4) (10,145.75)};\addlegendentry{$C_D$ total (Squire-Young)}
\addplot[cgreen,very thick] coordinates {(-4,54.092) (-3,55.071) (-2,56.389) (-1,57.475) (0,57.105) (1,57.167) (2,54.939) (3,55.464) (4,54.027) (5,52.844) (6,53.087) (7,54.737) (8,56.028) (9,56.418) (10,56.281)};\addlegendentry{frottement $C_{D,f}$}
\addplot[corange,very thick] coordinates {(-4,15.492) (-3,13.141) (-2,11.089) (-1,9.8128) (0,10.266) (1,10.933) (2,14.231) (3,17.115) (4,23.107) (5,29.827) (6,36.953) (7,43.687) (8,58.727) (9,72.98) (10,89.471)};\addlegendentry{forme $C_D-C_{D,f}$}
\end{groupplot}
\end{tikzpicture}

\caption{Polaire du NACA 2412 à $\Re=3\cdot10^6$, de $-4^\circ$ à $10^\circ$. Tous les points restent valides : à partir de $7^\circ$, l'extrados décolle, mais en aval de $x/c=0{,}9$ (drapeau \code{te\_separation}). À droite, la traînée de forme croît avec l'incidence alors que le frottement reste à peu près constant. La finesse maximale, environ 111, est atteinte vers $7^\circ$.}\label{fig:polaire}
\end{figure}

\paragraph{Lecture de la polaire.} Elle a la forme attendue : un minimum de traînée plat entre $-2^\circ$ et $0^\circ$ (67~counts), puis une croissance plus rapide aux fortes incidences, portée par la traînée de forme (épaississement de la couche limite de l'extrados sous gradient adverse) et par la remontée de la transition. En revanche, la portance reste linéaire : le modèle ne décroche pas (ch.~\ref{ch:validation}).

\section{Incidence imposée ou portance imposée : \code{alpha\_for\_cl}}
Deux façons de comparer des profils :
\begin{itemize}
\item à \textbf{incidence imposée} ($\alpha=2^\circ$) : un appel à \code{evaluate} ;
\item à \textbf{portance imposée} ($C_L^\star=0{,}4$, vol de croisière $L=W$) : il faut d'abord trouver l'incidence qui donne $C_L^\star$.
\end{itemize}
Comme $C_L(\alpha)$ non visqueux est quasi linéaire et monotone, une recherche de racine de Brent sur $[-10^\circ;15^\circ]$ converge en quelques itérations :
\begin{lstlisting}
def alpha_for_cl(m, p, t, CL_target, bracket=(-10.0, 15.0), n_half=100):
    X, Y = naca4(m, p, t, n_half=n_half)
    cache = {}                       # matrices d'influence et inverse partagees
    f = lambda a: solve(X, Y, a, cache=cache).CL - CL_target
    return brentq(f, *bracket, xtol=1e-6)
\end{lstlisting}
Grâce au cache, seule la première évaluation construit et inverse la matrice. Les suivantes ne font qu'un produit matrice-vecteur. Pour le NACA 2412, $C_L^\star=0{,}4$ donne $\alpha=1{,}19^\circ$. \code{evaluate\_at\_cl} enchaîne \code{alpha\_for\_cl} puis \code{evaluate}.

\section{Sensibilité aux paramètres numériques}
Le test V4 de \code{validation.py} vérifie que les paramètres numériques ne pilotent pas le résultat (NACA 2412, $\alpha=2^\circ$, $\Re=3\cdot10^6$) :
\begin{center}
\begin{tabular}{lcccc}
\toprule
\code{n\_half} & 60 & 100 & 150 & 200 \\
\midrule
$C_L$ & 0,4950 & 0,4975 & 0,4987 & 0,4993 \\
$C_D$ & 0,00690 & 0,00692 & 0,00692 & 0,00692 \\
$x_{tr,up}$ & 0,279 & 0,279 & 0,280 & 0,280 \\
\bottomrule
\end{tabular}

\medskip
\begin{tabular}{lc}
\toprule
\code{x\_stop} & $C_D$ \\
\midrule
0,970 & 0,00681 \\
0,980 & 0,00688 \\
0,985 & 0,00692 \\
0,990 & 0,00698 \\
\bottomrule
\end{tabular}
\end{center}
$C_D$ est stable au count près vis-à-vis du maillage, grâce à l'interpolation de la transition. Le choix de $x_{stop}$ déplace $C_D$ de façon presque uniforme sur tous les profils, ce qui ne change pas leur classement.

\begin{projet}
Pour l'optimisation, ce chapitre fixe le \emph{contrat} du simulateur : des entrées $(m,p,t)$ et un point de vol, des sorties toujours définies, et un booléen \code{valid} qui devient une contrainte. Le choix entre incidence imposée et portance imposée change la forme du problème. Il est discuté au chapitre suivant.
\end{projet}

\exercices
\begin{exercice}
Pour le NACA 2412, la méthode des panneaux donne une pente $\dd C_L/\dd\alpha=6{,}88$~/rad et $\alpha_{L0}=-2{,}13^\circ$. Estimer à la main l'incidence qui donne $C_L=0{,}4$, et comparer à la valeur $1{,}19^\circ$ trouvée par \code{alpha\_for\_cl}.
\end{exercice}
\begin{exercice}
À $\alpha=7^\circ$, le NACA 2412 renvoie \code{valid=True} avec deux drapeaux : \code{up:turbulent\_\allowbreak separation} et \code{te\_separation}. Expliquer ce que signifie chacun, et pourquoi le résultat reste valide. Que faudrait-il observer pour qu'il ne le soit plus ?
\end{exercice}

\chapter{L'espace de conception et les interfaces}\label{ch:espace}
\iterbanner{V0}

\section{Pourquoi une grille discrète}
La formulation QUBO travaille sur des bits. On fixe donc dès la phase~1 une grille de niveaux pour $(m,p,t)$, et c'est \emph{cette} grille que partagent les trois méthodes : NSGA-II (variables entières), la boucle QUBO (bits) et l'énumération exhaustive (front exact). Si elles ne travaillaient pas sur la même grille, la comparaison n'aurait pas de sens.

\begin{center}
\begin{tabular}{lccc}
\toprule
Variable & niveaux & valeurs & bits \\
\midrule
cambrure $m$ & 8 & $0,\ 1,\ \dots,\ 7$~\% & 3 \\
position $p$ & 8 & $0{,}20,\ 0{,}25,\ \dots,\ 0{,}55$ & 3 \\
épaisseur $t$ & 8 & $6,\ 8,\ \dots,\ 20$~\% & 3 \\
\bottomrule
\end{tabular}
\end{center}
Cela fait $8^3=512$ profils et 9 bits, donc \textbf{9 qubits}.

\paragraph{L'encodage binaire.} Chaque variable prend $2^k$ niveaux d'indice $i=\sum_{j=0}^{k-1}2^j\,b_j$, bit de poids faible en premier. Les 9 bits sont concaténés dans l'ordre $(m,p,t)$ :
\begin{lstlisting}
def encode(self, idx):                 # (im, ip, it) -> 9 bits
    bits = []
    for i, nb in zip(idx, self.bits_per_var):
        bits += [(i >> j) & 1 for j in range(nb)]
    return np.array(bits, dtype=int)
\end{lstlisting}
Un polynôme quadratique en $(m,p,t)$ reste quadratique en bits : chaque variable est affine en bits, et $b^2=b$. C'est ce qui permet de passer du modèle de substitution au QUBO (cours compagnon, chapitre QUBO).

\section{Le simulateur compté : \code{AeroSimulator}}
La métrique de comparaison entre NSGA-II et la boucle QUBO est le \textbf{nombre de simulations distinctes} nécessaires pour retrouver le front exact. Il faut donc le compter au même endroit pour les deux méthodes :
\begin{lstlisting}
class AeroSimulator:
    def __call__(self, idx):
        idx = tuple(int(i) for i in idx)
        if idx not in self.cache:                 # seule une NOUVELLE simulation compte
            m, p, t = self.space.values(idx)
            r = evaluate_at_cl(...) if self.case.mode == "cl" else evaluate(...)
            self.cache[idx] = r
            self.history.append(idx)
        return self.cache[idx]

    @property
    def n_sims(self):
        return len(self.cache)
\end{lstlisting}
NSGA-II réévalue souvent les mêmes individus d'une génération à l'autre ; sans cache, il serait pénalisé à tort. La méthode \code{constraints(r)} renvoie les contraintes au format de pymoo ($g\le0$) : épaisseur minimale, portance minimale en mode incidence, et validité.

\paragraph{Le point de vol : \code{FlightCase}.} Il regroupe $\Re=3\cdot10^6$, le mode (\code{'alpha'} ou \code{'cl'}), $\alpha=2^\circ$ ou $C_L^\star=0{,}4$, $C_{L,min}=0{,}3$ et $t_{min}=8$~\%. Ces valeurs sont à figer en équipe avant toute comparaison.

\section{La base exhaustive}
\code{generer\_base.py} évalue les 512 profils dans chaque mode, en quelques secondes, et écrit un CSV par mode. C'est à la fois la vérité de référence pour le front exact, une table de consultation pour développer les phases suivantes, et un outil de diagnostic.

\begin{figure}[htbp]
\centering
% figure générée par tools/make_figures.py : base_cartes
\begin{tikzpicture}
\begin{groupplot}[group style={group size=2 by 1,horizontal sep=3.2cm},width=0.36\linewidth,height=5.6cm,
  xlabel={cambrure $m$ (\%)},ylabel={épaisseur $t$ (\%)},xtick={0,...,7},ytick={6,8,...,20},
  enlargelimits=false,axis on top,title style={font=\small},tick label style={font=\scriptsize},
  colormap={bl}{color=(white) color=(cblue!45) color=(cblue!85!black)}]
\nextgroupplot[title={$\alpha=2^\circ$ imposé},colorbar,colorbar style={ylabel={$C_D$ (counts)},width=2.5mm,tick label style={font=\scriptsize}}]
\addplot[matrix plot*,mesh/cols=8,point meta=explicit] coordinates {(0,6) [61.11] (1,6) [61.4] (2,6) [58.78] (3,6) [53] (4,6) [52.2] (5,6) [66.3] (6,6) [68.25] (7,6) [71.72] (0,8) [64.25] (1,8) [64.6] (2,8) [63.92] (3,8) [62.85] (4,8) [67.1] (5,8) [68.91] (6,8) [71.27] (7,8) [74.91] (0,10) [66.61] (1,10) [66.91] (2,10) [66.75] (3,10) [66.72] (4,10) [69.89] (5,10) [71.46] (6,10) [74.32] (7,10) [78.15] (0,12) [68.68] (1,12) [69] (2,12) [69.17] (3,12) [69.72] (4,12) [72.31] (5,12) [73.97] (6,12) [77.35] (7,12) [81.31] (0,14) [70.65] (1,14) [71.02] (2,14) [71.46] (3,14) [72.79] (4,14) [74.63] (5,14) [76.47] (6,14) [80.18] (7,14) [85.14] (0,16) [72.56] (1,16) [73.01] (2,16) [73.6] (3,16) [75.08] (4,16) [76.57] (5,16) [79.46] (6,16) [84.27] (7,16) [89.53] (0,18) [74.37] (1,18) [74.8] (2,18) [75.38] (3,18) [77.05] (4,18) [78.86] (5,18) [83.49] (6,18) [88.57] (7,18) [94.14] (0,20) [75.8] (1,20) [76.25] (2,20) [77.9] (3,20) [80.23] (4,20) [82.94] (5,20) [87.74] (6,20) [93.08] (7,20) [98.99]};

\nextgroupplot[title={$C_L=0{,}4$ imposé},colorbar,colorbar style={ylabel={$C_D$ (counts)},width=2.5mm,tick label style={font=\scriptsize}}]
\addplot[matrix plot*,mesh/cols=8,point meta=explicit] coordinates {(0,6) [63.37] (1,6) [61.6] (2,6) [59.09] (3,6) [61.54] (4,6) [60.08] (5,6) [59.19] (6,6) [59.3] (7,6) [62.39] (0,8) [66.09] (1,8) [65.51] (2,8) [62.85] (3,8) [63.84] (4,8) [62.74] (5,8) [61.62] (6,8) [61.86] (7,8) [63.42] (0,10) [69.3] (1,10) [67.78] (2,10) [66.19] (3,10) [66.04] (4,10) [65.25] (5,10) [63.92] (6,10) [64.5] (7,10) [65.52] (0,12) [71.67] (1,12) [69.66] (2,12) [68.36] (3,12) [68.18] (4,12) [67.63] (5,12) [66.67] (6,12) [67.05] (7,12) [67.15] (0,14) [73.57] (1,14) [71.5] (2,14) [70.39] (3,14) [70.26] (4,14) [69.89] (5,14) [69.15] (6,14) [69.49] (7,14) [69.98] (0,16) [75.24] (1,16) [73.3] (2,16) [72.37] (3,16) [72.29] (4,16) [72.05] (5,16) [71.46] (6,16) [71.83] (7,16) [72.86] (0,18) [76.67] (1,18) [74.96] (2,18) [74.21] (3,18) [74.38] (4,18) [74.03] (5,18) [73.59] (6,18) [73.79] (7,18) [75.68] (0,20) [77.85] (1,20) [76.33] (2,20) [75.64] (3,20) [76.88] (4,20) [75.71] (5,20) [75.21] (6,20) [76.79] (7,20) [78.65]};

\end{groupplot}
\end{tikzpicture}

\caption{$C_D$ (counts) sur la tranche $p=0{,}40$ de la grille, dans les deux modes ($\Re=3\cdot10^6$). À incidence imposée, la traînée croît à la fois avec l'épaisseur et la cambrure. À portance imposée, la cambrure ne coûte presque plus rien : elle remplace l'incidence. Les croix marquent les profils invalides.}\label{fig:base}
\end{figure}

\begin{center}
\begin{tabularx}{\linewidth}{Xcc}
\toprule
 & $\alpha=2^\circ$ imposé & $C_L=0{,}4$ imposé \\
\midrule
profils valides & 512 / 512 & 507 / 512 \\
profils faisables (valides, $t\ge8$~\%, et $C_L\ge0{,}3$ en mode $\alpha$) & 392 & 446 \\
plage de $C_D$ (valides) & 51 à 124 counts & 51 à 88 counts \\
corrélation de $C_D$ avec $t$ / avec $m$ & 0,60 / 0,62 & 0,74 / 0,04 \\
meilleur profil faisable $(m,p,t)$ & (4~\%, 0,35, 8~\%) & (4~\%, 0,55, 8~\%) \\
\quad son $C_D$ & 55 counts & 54 counts \\
incidence & $2^\circ$ pour tous & $-5{,}5^\circ$ à $+3{,}5^\circ$ \\
drapeau \code{te\_separation} & 145 profils & 64 profils \\
\bottomrule
\end{tabularx}
\end{center}

\paragraph{Lecture.} Le choix du mode change la forme du problème :
\begin{itemize}
\item à incidence imposée, on compare des profils qui ne portent pas le même poids. Un profil très cambré « gagne » de la portance qu'il paie en traînée, d'où la corrélation de 0,62 avec $m$. Il faut une contrainte $C_L\ge C_{L,min}$ ;
\item à portance imposée, on compare à mission égale, comme un avionneur. La corrélation avec $m$ disparaît, l'épaisseur domine, et la contrainte de portance disparaît du QUBO. En contrepartie, l'incidence devient une \emph{sortie} de l'aéro et une \emph{entrée} de la SER : l'attitude vue par le radar dépend du profil. C'est un vrai couplage multiphysique.
\end{itemize}
La recommandation de la note d'interfaces est le mode portance imposée, avec le mode incidence comme repli.

\section{Vers la phase 2 : la géométrie pour la SER}
La SER a besoin d'une surface. \code{geometry\_export.extrude} transforme le profil en un tronçon d'aile extrudé selon l'envergure, facetté en triangles, avec les \emph{mêmes points} que le calcul de panneaux (figure~\ref{fig:extrusion}). Chaque panneau 2D devient un rectangle de deux triangles, et les deux flancs sont triangulés en bandes entre points d'extrados et d'intrados de même indice.

\begin{figure}[htbp]
\centering
\begin{tikzpicture}[x={(1cm,0cm)},y={(0cm,1cm)},z={(0.55cm,0.35cm)},scale=1.25]
\def\prof{(4,0,0) .. controls (3,0.25,0) and (1.5,0.45,0) .. (0.6,0.4,0) .. controls (0.1,0.32,0) and (0,0.1,0) .. (0,0,0) .. controls (0,-0.08,0) and (0.3,-0.2,0) .. (1,-0.22,0) .. controls (2,-0.2,0) and (3,-0.08,0) .. (4,0,0)}
\fill[cblue!8] (4,0,0) -- (4,0,4) -- (0,0,4) -- (0,0,0) -- cycle;
\foreach \z in {0,4}{
  \draw[cink,thick] plot[smooth,tension=0.6] coordinates {(4,0,\z) (3,0.2,\z) (2,0.34,\z) (1,0.38,\z) (0.4,0.3,\z) (0.08,0.15,\z) (0,0,\z) (0.1,-0.1,\z) (0.5,-0.18,\z) (1.5,-0.2,\z) (2.5,-0.14,\z) (3.3,-0.06,\z) (4,0,\z)};}
\foreach \x/\y in {4/0,3/0.2,2/0.34,1/0.38,0.4/0.3,0.08/0.15,0/0}{\draw[cmuted,thin] (\x,\y,0) -- (\x,\y,4);}
\foreach \x/\y/\xx/\yy in {3/0.2/4/0,2/0.34/3/0.2,1/0.38/2/0.34}{\draw[cmuted,thin] (\x,\y,0) -- (\xx,\yy,4);}
\draw[arr,cmuted] (0,0,0) -- (1.2,0,0) node[below,font=\small] {$x$ (corde)};
\draw[arr,cmuted] (4,0,0) -- (4,0,5) node[right,font=\small] {$z$ (envergure)};
\node[font=\small,align=center] at (2.6,-0.9,0) {chaque panneau 2D $\to$ un rectangle\\ de deux triangles};
\draw[corange,very thick,-{Stealth[length=2.5mm]}] (-2.2,-0.4,2) -- (-0.2,-0.02,2) node[pos=0,left,font=\small,corange,align=right] {radar de face,\\ tourné de $\alpha$};
\end{tikzpicture}
\caption{L'extrusion 2,5D exportée en STL. L'incidence n'est pas appliquée à la géométrie : elle est écrite dans les métadonnées JSON, et c'est la direction radar qui doit être tournée.}\label{fig:extrusion}
\end{figure}

Conventions du fichier : repère corps ($x$ le long de la corde vers l'aval, $y$ vers l'extrados, $z$ selon l'envergure, origine au bord d'attaque), normales sortantes, triangles dégénérés retirés au bord d'attaque et au bord de fuite fermé. Deux contrôles garantissent une surface fermée et bien orientée : $\sum_k\vect n_kA_k=0$, et le volume calculé par le théorème de la divergence égal à l'aire du profil fois l'envergure. Pour le NACA 2412 avec 200 panneaux, le STL compte 796 triangles.

\begin{piege}
Le facettage est un \emph{paramètre numérique} pour la SER, pas une variable de conception. À haute fréquence, chaque facette plane produit des reflets spéculaires qui n'existent pas sur la forme lisse. Il faut figer \code{n\_half} pour tout le hackathon et vérifier la convergence de la SER avec lui, comme on l'a fait pour $C_D$.
\end{piege}

\section{Vers la phase 4 : la sortie est-elle « QUBO-compatible » ?}
La boucle quantique remplace le simulateur par un modèle quadratique ajusté par moindres carrés, en variables normalisées $z_i=2x_i/7-1\in[-1,1]$ :
\[
\widehat{C_D}(z)=c_0+\sum_ic_iz_i+\sum_{i\le l}c_{il}\,z_iz_l\qquad(10\ \text{coefficients}).
\]
\code{analyse\_quadratique.py} mesure, sur la base exhaustive, la qualité de ce modèle.

\begin{figure}[htbp]
\centering
% figure générée par tools/make_figures.py : quadratique
\begin{tikzpicture}
\begin{axis}[width=0.62\linewidth,height=6cm,xlabel={$C_D$ du modèle aéro (counts)},ylabel={$C_D$ quadratique (counts)},
  legend style={draw=none,font=\scriptsize,at={(0.02,0.98)},anchor=north west},legend cell align=left]
\addplot[only marks,mark=*,mark size=1.1pt,cblue,fill opacity=0.6,draw opacity=0.6] coordinates {(61.11,58.04) (64.25,60.84) (66.61,63.59) (68.68,66.3) (70.65,68.97) (72.56,71.59) (74.37,74.17) (75.8,76.71) (61.11,57.16) (64.25,59.78) (66.61,62.36) (68.68,64.89) (70.65,67.39) (72.56,69.84) (74.37,72.25) (75.8,74.62) (61.11,57.22) (64.25,59.66) (66.61,62.07) (68.68,64.43) (70.65,66.75) (72.56,69.03) (74.37,71.27) (75.8,73.46) (61.11,58.22) (64.25,60.49) (66.61,62.73) (68.68,64.92) (70.65,67.06) (72.56,69.17) (74.37,71.23) (75.8,73.25) (61.11,60.17) (64.25,62.27) (66.61,64.33) (68.68,66.35) (70.65,68.32) (72.56,70.25) (74.37,72.14) (75.8,73.99) (61.11,63.06) (64.25,64.99) (66.61,66.88) (68.68,68.72) (70.65,70.52) (72.56,72.28) (74.37,74) (75.8,75.67) (61.11,66.9) (64.25,68.65) (66.61,70.37) (68.68,72.04) (70.65,73.67) (72.56,75.25) (74.37,76.8) (75.8,78.3) (61.11,71.68) (64.25,73.26) (66.61,74.8) (68.68,76.3) (70.65,77.76) (72.56,79.17) (74.37,80.54) (75.8,81.87) (56.06,59.39) (60.79,62.43) (64.62,65.42) (67.78,68.38) (70.46,71.29) (72.85,74.16) (74.92,76.99) (76.41,79.78) (56.99,57.97) (60.92,60.83) (64.15,63.66) (67.08,66.44) (69.84,69.18) (72.36,71.88) (74.51,74.53) (76.07,77.15) (58.51,57.5) (61.96,60.19) (64.75,62.84) (67.25,65.45) (69.62,68.02) (71.91,70.54) (74.04,73.02) (75.63,75.46) (60,57.97) (63.28,60.49) (65.78,62.97) (68.04,65.4) (70.2,67.8) (72.29,70.15) (74.2,72.45) (75.65,74.72) (61.4,59.38) (64.6,61.73) (66.91,64.04) (69,66.3) (71.02,68.52) (73.01,70.7) (74.8,72.83) (76.25,74.93) (62.8,61.75) (65.88,63.92) (68.01,66.05) (69.96,68.14) (71.88,70.19) (73.78,72.19) (75.48,74.16) (76.93,76.08) (64.25,65.05) (67.04,67.05) (68.94,69.01) (70.7,70.93) (72.46,72.8) (74.23,74.63) (75.84,76.42) (77.34,78.17) (65.6,69.31) (67.83,71.13) (69.48,72.92) (71.14,74.66) (72.83,76.36) (74.55,78.02) (76.1,79.64) (77.67,81.21) (59.66,64.9) (64.32,68.15) (67.98,71.35) (71.07,74.51) (73.73,77.62) (75.83,80.7) (78.4,83.73) (52.62,59.67) (58.87,62.78) (63.49,65.85) (67.13,68.88) (70.28,71.86) (73,74.81) (75.19,77.71) (77.82,80.57) (53.06,58.66) (59.22,61.6) (63.35,64.5) (66.64,67.36) (69.61,70.17) (55.22,58.61) (61.09,61.37) (64.62,64.1) (67.47,66.78) (70.03,69.42) (72.38,72.01) (74.37,74.57) (58.78,59.49) (63.92,62.08) (66.75,64.63) (69.17,67.14) (71.46,69.61) (73.6,72.03) (75.38,74.41) (77.9,76.75) (63.61,61.32) (67.03,63.74) (69.14,66.12) (71.26,68.45) (73.19,70.75) (75.03,73) (76.49,75.2) (79.06,77.37) (67.91,64.1) (69.63,66.34) (70.99,68.55) (72.52,70.71) (74.19,72.83) (75.85,74.91) (77.1,76.94) (79.9,78.93) (70.18,67.82) (71.14,69.89) (72.19,71.92) (73.54,73.91) (75.08,75.86) (76.62,77.76) (77.93,79.62) (80.71,81.44) (77.35,85.29) (80.94,88.57) (51.16,60.13) (58.99,63.15) (64.2,66.11) (67.67,69.04) (53,60.49) (62.85,63.33) (66.72,66.12) (69.72,68.88) (72.79,71.59) (75.08,74.26) (77.05,76.88) (80.23,79.47) (56.46,61.79) (67.66,64.45) (70.35,67.07) (72.34,69.66) (74.32,72.19) (76.15,74.69) (78.19,77.14) (81.28,79.55) (71.04,64.03) (71.33,66.52) (72.45,68.97) (73.99,71.38) (75.71,73.74) (77.31,76.07) (79.56,78.35) (82.58,80.58) (74.39,67.22) (74.06,69.54) (74.6,71.82) (75.79,74.05) (77.27,76.24) (78.63,78.39) (81.12,80.5) (84.1,82.56) (51.78,62.55) (55.1,65.81) (52.2,62.37) (67.1,65.46) (69.89,68.5) (72.31,71.5) (74.63,74.46) (76.57,77.37) (54.35,63.14) (68.71,66.05) (70.92,68.92) (73.13,71.75) (75.3,74.53) (77.12,77.27) (80.28,79.97) (83.6,82.63) (71.42,64.85) (71.78,67.59) (73.25,70.29) (75.08,72.94) (77,75.55) (79,78.12) (82.11,80.64) (85.38,83.13) (78.01,67.51) (76.09,70.08) (76.51,72.6) (77.81,75.08) (79.33,77.51) (81.53,79.91) (84.51,82.26) (87.69,84.57) (66.3,65.15) (68.91,68.48) (71.46,71.77) (73.97,75.01) (66.37,65.39) (68.9,68.54) (71.45,71.66) (73.97,74.73) (76.31,77.76) (79.1,80.74) (82.53,83.69) (86.13,86.59) (69.82,66.57) (71.48,69.55) (73.62,72.49) (75.91,75.39) (78.02,78.25) (81.24,81.06) (84.67,83.83) (88.24,86.56) (78.86,68.69) (76.86,71.5) (77.8,74.27) (79.47,77) (81.59,79.68) (84.66,82.32) (87.97,84.92) (91.49,87.47) (66.81,68.52) (69.6,71.92) (72.4,75.28) (75.14,78.6) (77.8,81.87) (81.41,85.11) (68.46,69.17) (71.11,72.4) (73.91,75.59) (76.66,78.73) (79.79,81.83) (83.48,84.89) (87.28,87.91) (91.24,90.88) (76.08,70.77) (76.52,73.82) (78.49,76.83) (80.81,79.8) (84.13,82.73) (87.69,85.62) (91.43,88.46) (95.35,91.26) (68.18,72.55) (71.1,76.19) (74.07,79.8) (67.88,72.66) (71.13,76.14) (74.42,79.57) (77.71,82.96) (81.7,86.31) (85.79,89.61) (90.02,92.87) (94.37,96.09) (73.02,73.73) (75.61,77.03) (78.71,80.29) (82.39,83.5) (86.38,86.68) (90.52,89.81) (94.78,92.9) (99.18,95.94)};\addlegendentry{transition extrados : Michel}
\addplot[only marks,mark=*,mark size=1.1pt,corange,fill opacity=0.6,draw opacity=0.6] coordinates {(53.81,61.62) (72.31,72.94) (75.28,75.66) (78.85,78.35) (76.45,77.08) (55.45,64.74) (60.94,68.27) (67.16,71.76) (72.15,75.21) (76.45,78.61) (80.16,81.97) (52.97,62.26) (58.8,65.62) (65.12,68.93) (69.85,72.21) (73.88,75.44) (77.41,78.63) (80.88,81.77) (85.57,84.87) (50.53,60.72) (57.23,63.91) (63.54,67.05) (67.98,70.15) (71.76,73.21) (75.13,76.22) (78.28,79.2) (82.79,82.13) (69.91,71.93) (73.1,74.77) (75.94,77.57) (80.17,80.32) (76.74,68.75) (80.15,72.53) (83.58,76.26) (74.61,79.95) (79.39,83.6) (83.22,87.21) (88.67,90.78) (94.19,94.3) (73.14,65.74) (76.13,69.34) (66.65,72.9) (72.39,76.42) (76.86,79.9) (80.4,83.33) (85.36,86.73) (90.5,90.07) (53.59,63.67) (57,67.1) (65.4,70.49) (70.68,73.84) (74.75,77.14) (78.05,80.4) (82.45,83.62) (87.22,86.79) (65,69.02) (69.47,72.2) (73.16,75.33) (76.24,78.41) (80.15,81.46) (84.57,84.46) (78.86,80.24) (82.94,83.07) (79.65,73.65) (83.29,77.68) (87.03,81.66) (90.81,85.59) (94.53,89.49) (87.78,93.34) (93.76,97.15) (99.85,100.9) (76.11,70.11) (79.47,73.96) (82.89,77.77) (86.34,81.53) (89.51,85.25) (84.71,88.93) (90.18,92.57) (95.8,96.16) (72.62,67.51) (75.7,71.19) (78.8,74.82) (81.92,78.41) (84.8,81.96) (82.24,85.47) (87.21,88.93) (92.3,92.35) (68.97,65.86) (71.86,69.36) (74.78,72.82) (77.71,76.24) (80.53,79.61) (83.52,82.95) (87.56,86.24) (91.71,89.49) (76.47,78.21) (79.46,81.37) (83.49,84.49) (87.74,87.57) (82.81,79.45) (86.66,83.71) (90.69,87.94) (94.82,92.12) (98.79,96.26) (104.5,100.4) (110.2,104.4) (106,108.4) (79.07,75.37) (82.68,79.47) (86.37,83.52) (90.12,87.53) (93.6,91.49) (98.69,95.42) (103.7,99.3) (101.7,103.1) (75.68,72.24) (78.99,76.16) (82.37,80.04) (85.77,83.88) (88.9,87.67) (93.47,91.42) (98.16,95.13) (98.07,98.8) (72.2,70.06) (75.28,73.8) (78.41,77.51) (81.56,81.17) (84.45,84.79) (88.72,88.37) (93.16,91.91) (97.78,95.4) (68.25,68.82) (71.27,72.39) (74.32,75.92) (77.35,79.41) (80.18,82.86) (84.27,86.27) (88.57,89.63) (93.08,92.95) (84.16,88.3) (88.8,91.44) (86.36,86.13) (90.41,90.64) (94.65,95.11) (98.91,99.54) (104,103.9) (110.2,108.3) (116.8,112.6) (123.7,116.8) (82.6,81.52) (86.43,85.86) (90.39,90.16) (94.34,94.41) (98.99,98.63) (104.5,102.8) (110.2,106.9) (116.1,111) (78.97,77.86) (82.54,82.03) (86.16,86.15) (89.77,90.23) (94.02,94.27) (99.06,98.27) (104.3,102.2) (109.7,106.1) (75.55,75.14) (78.85,79.14) (82.24,83.09) (85.56,87) (89.52,90.86) (94.17,94.69) (99.01,98.47) (104.1,102.2) (71.72,73.37) (74.91,77.19) (78.15,80.97) (81.31,84.71) (85.14,88.4) (89.53,92.05) (94.14,95.66) (98.99,99.22) (76.61,83.36) (80.49,86.88) (84.82,90.36) (89.35,93.79) (94.21,97.19)};\addlegendentry{transition extrados : décollement laminaire}
\addplot[cmuted,dashed] coordinates {(50.5,50.5) (123.7,123.7)};
\node[font=\small,anchor=south east] at (rel axis cs:0.98,0.03) {$R^2=0.908$};
\end{axis}
\end{tikzpicture}

\caption{Modèle quadratique à 10 coefficients ajusté sur les 512 profils (mode $\alpha=2^\circ$), contre la vraie valeur. Les couleurs distinguent le mode de transition sur l'extrados. Le $R^2$ de 0,91 est limité par la bascule entre transition de Michel et transition par décollement laminaire.}\label{fig:quad}
\end{figure}

\begin{center}
\begin{tabular}{lcc}
\toprule
Sortie & quadratique en $(m,p,t)$ (10 coef.) & QUBO libre sur 9 bits (46 coef.) \\
\midrule
$C_L$ & $R^2=0{,}9998$ & $0{,}9998$ \\
$C_D$ & $R^2=0{,}908$ & $0{,}920$ \\
\bottomrule
\end{tabular}
\end{center}
\begin{itemize}
\item $C_L$ est quasi affine en $(m,\alpha)$ par construction (panneaux linéaires en $\alpha$, théorie des profils minces linéaire en $m$) : un quadratique le représente parfaitement.
\item Les 9~\% de variance de $C_D$ non expliqués sont structurels : c'est la bascule de mode de transition du chapitre~\ref{ch:transition}. Interpoler $x_{tr}$ entre stations n'a fait passer le $R^2$ que de 0,899 à 0,908.
\item Avec 20 profils tirés au hasard, le quadratique classe bien les profils (corrélation de rang de Spearman d'environ 0,89), mais son minimum n'est dans le vrai top 5~\% qu'une fois sur trois environ.
\end{itemize}

\begin{projet}
Conséquence pour la phase~4 : un QUBO global ajusté une seule fois ne suffit pas. Il faut la \textbf{boucle itérative} de la note d'intention, qui réajuste le modèle autour des meilleurs candidats. C'est ce que montre le cours compagnon : la boucle retrouve le vrai minimum de $C_D$ en 25 simulations (médiane sur 50 essais), contre environ 256 en tirant les profils au hasard.
\end{projet}

\exercices
\begin{exercice}
(a) Donner la chaîne des 9 bits du profil d'indices $(i_m,i_p,i_t)=(4,3,1)$, et les valeurs physiques correspondantes. (b) Combien de chaînes de bits distinctes représentent un profil symétrique ($m=0$) ?
\end{exercice}
\begin{exercice}
Un collègue propose de comparer NSGA-II et la boucle QUBO en temps de calcul plutôt qu'en nombre de simulations distinctes. Donner deux raisons pour lesquelles ce serait une mauvaise métrique dans ce hackathon.
\end{exercice}

\chapter{Validation, limites et pistes}\label{ch:validation}
\iterbanner{V0}

\section{Ce qui a été vérifié}
\code{validation.py} enchaîne cinq tests. Les quatre premiers sont des vérifications du code contre des solutions connues ; le cinquième est une validation contre XFOIL, réalisée en V1 (§\ref{sec:v5}).
\begin{center}
\renewcommand{\arraystretch}{1.2}
\begin{tabularx}{\linewidth}{lXXl}
\toprule
Test & Grandeur & Référence & Résultat \\
\midrule
V1 & $C_L$, NACA 0012, $5^\circ$ & convergence en $N$ & 0,6026 à 200 panneaux, 0,06~\% de la limite \\
V1 & $C_{L,C_p}$ contre $C_{L,KJ}$ & égalité théorique & écart $<0{,}4$~\% à 200 panneaux, en $1/N$ \\
V1 & $C_{D,p}$ & 0 (d'Alembert) & $2\cdot10^{-4}$ à 200 panneaux, en $1/N$ \\
V1 & pente $\dd C_L/\dd\alpha$ & $2\pi(1+0{,}77t)=6{,}86$~/rad & 6,89 (0012), 6,88 (2412) \\
V1 & $C_{M,c/4}$, NACA 2412 & $\approx-0{,}053$ (profil mince) & $-0{,}055$ \\
V2 & $\theta$ laminaire, plaque plane & Blasius & 1,0~\% ; $H=2{,}61$ (Blasius 2,59) \\
V3 & $\theta$ turbulent, plaque plane & loi en $1/7$ & $<5$~\% ; $H\to1{,}34$ \\
V4 & $C_D$ selon \code{n\_half} & --- & $\pm0{,}1$~\% de 60 à 200 \\
V4 & $C_D$ selon \code{x\_stop} & --- & au plus $\pm1{,}5$~\% de 0,97 à 0,99 \\
V5 & $C_D$, $C_L$ sur les 512 profils & XFOIL, $N_{crit}=9$ & rang 0,91 ($C_D$), 0,94 ($C_L$) ; voir §\ref{sec:v5} \\
\bottomrule
\end{tabularx}
\end{center}

\section{Le test V5 : XFOIL sur toute la grille\texorpdfstring{\ \protect\iterbadge{V1}}{ (V1)}}\label{sec:v5}
\begin{evolution}{V1}{section ajoutée}Le test V5 était « à faire » en V0, faute d'XFOIL. Il a été réalisé en V1 sur toute la grille ; ses conclusions motivent l'itération V2 (XFOIL comme simulateur de référence, voir la fiche projet).\end{evolution}
XFOIL est installé localement (paquet Debian extrait dans \code{Code/.tools/}, sans droits administrateur). Cette compilation ne résout le visqueux que si son affichage est actif : \code{xfoil\_runner.py} démarre donc un écran virtuel Xvfb quand aucun écran n'est disponible. Le script \code{validation\_xfoil\_grille.py} calcule les 512 profils à $\alpha=2^\circ$, $\Re=3\cdot10^6$, $N_{crit}=9$, en une dizaine de minutes sur 8 cœurs.

\begin{center}
\begin{tabular}{lcc}
\toprule
 & $C_D$ & $C_L$ \\
\midrule
profils où XFOIL converge & \multicolumn{2}{c}{498 / 512} \\
écart moyen (notre modèle / XFOIL $-1$) & $+27$~\% & $+25$~\% \\
corrélation de rang de Spearman & 0,91 & 0,94 \\
NACA 2412 & 69 contre 50 counts & 0,498 contre 0,436 \\
\bottomrule
\end{tabular}
\end{center}

\paragraph{Lecture.} Les deux biais vont dans le sens prévu, mais ils sont plus forts qu'estimé :
\begin{itemize}
\item la \textbf{traînée} est surestimée de 27~\% en moyenne, parce que Michel fait transitionner trop tôt. Sur le NACA 2412, la transition de l'extrados est en $x/c=0{,}28$ chez nous et $0{,}43$ dans XFOIL ;
\item la \textbf{portance} est surestimée de 25~\% en moyenne, et non de 5 à 10~\% : à $\alpha=2^\circ$, le décambrage visqueux pèse lourd, surtout pour les profils très cambrés vers l'arrière. La contrainte $C_L\ge0{,}3$ est donc trop optimiste dans notre modèle (359 profils faisables avec XFOIL, contre 392).
\end{itemize}
Le classement est en revanche bien conservé (rang 0,91 et 0,94), ce qui est l'essentiel pour l'optimisation.

\paragraph{Et le front de Pareto ?} En combinant la traînée d'XFOIL avec la même SER, on obtient le front « de référence ». Deux conclusions :
\begin{itemize}
\item \textbf{tous les profils du front restent à $t=8$~\%} : le résultat structurant (compromis dégénéré en épaisseur) ne dépend pas du modèle aérodynamique ;
\item \textbf{la position de cambrure optimale, elle, change.} XFOIL préfère des cambrures plus en arrière ($p=0{,}40$ à $0{,}45$) pour les profils de faible traînée, parce que sa transition tardive récompense les écoulements laminaires étendus. Dans le scénario S1, un seul profil sur 5 est commun aux deux fronts ; dans S2, 13 sur 18.
\end{itemize}
C'est la piste d'amélioration prioritaire : caler la transition sur XFOIL (facteur sur le seuil de Michel, ou critère dépendant du gradient), puis refaire la comparaison des fronts.

\section{Couplage faible contre couplage fort}
Notre modèle et XFOIL résolvent la même physique, mais pas de la même façon. Dans le \textbf{couplage faible}, on calcule l'écoulement extérieur une fois, puis la couche limite sur la vitesse obtenue, sans retour. Dans le \textbf{couplage fort}, l'écoulement extérieur est déplacé par l'épaisseur $\delta^*$ et tout est résolu simultanément par une méthode de Newton.

\begin{center}
\renewcommand{\arraystretch}{1.2}
\begin{tabularx}{\linewidth}{lXX}
\toprule
Effet & Notre modèle (faible) & XFOIL (fort) \\
\midrule
Décambrage visqueux, $C_L$ réel plus bas & absent : $C_L$ surestimé de 5 à 10~\% & capturé \\
Point d'arrêt du bord de fuite fermé & artefact, contourné par \code{x\_stop} & effacé par $\delta^*$ et le sillage \\
Décollement, bulbe laminaire & seulement détecté (mode direct singulier) & traité en mode inverse \\
Transition & Michel (corrélation) & $e^N$ (amplification des ondes) \\
Sillage & absent (Squire-Young au bord de fuite) & calculé \\
Robustesse & jamais de divergence & échecs possibles, silencieux \\
Coût & environ 8~ms & 0,1 à 1~s par incidence \\
\bottomrule
\end{tabularx}
\end{center}

\paragraph{Pourquoi le mode direct est singulier au décollement.} En mode direct, on impose $U_e(s)$ et on calcule $\theta$ et $H$. Près de $c_f=0$, l'équation intégrale fait apparaître $\dd H/\dd s\to\infty$ : c'est la singularité de Goldstein. En mode inverse, on impose $\delta^*$ ou $H$ et on en déduit $U_e$, et le problème redevient régulier. XFOIL passe de l'un à l'autre automatiquement en résolvant tout ensemble. Notre modèle ne peut que détecter le décollement, et c'est pourquoi il le signale au lieu de le calculer.

\paragraph{Une amélioration intermédiaire : la vitesse de transpiration.} On peut faire « sentir » la couche limite à l'écoulement extérieur sans changer la géométrie, en remplaçant la condition de glissement par une vitesse normale de soufflage
\[
V_n=\frac{\dd(U_e\,\delta^*)}{\dd s},
\]
puis en recalculant les panneaux et la couche limite, et en itérant. Cela converge tant qu'il n'y a pas de décollement, et cela corrigerait l'essentiel du décambrage visqueux de $C_L$. Ce n'est pas implémenté.

\section{Le biais connu, et comment le mesurer}
Sur le NACA 0012 à $\alpha=0$ et $\Re=10^6$, le modèle donne $C_D=75$~counts avec une transition en $x/c=0{,}46$. XFOIL est couramment cité vers 54~counts, valeur à reconfirmer. Le test à transition imposée isole la cause : en imposant la transition au décollement laminaire ($x\approx0{,}62$), on obtient 62~counts. L'essentiel de l'écart vient donc de Michel, trop précoce à ce Reynolds ; l'écart restant, environ 14~\%, vient de la couche limite et de Squire-Young. À $\Re=3\cdot10^6$, Michel est plus réaliste et l'écart se réduit.

\paragraph{Le protocole V5.} Il est écrit dans \code{validation.py} et s'exécute dès qu'XFOIL est dans le \code{PATH} :
\begin{enumerate}
\item comparer $C_L$, $C_D$ et $x_{tr}$ en transition libre, sur une polaire de trois profils ;
\item recalculer notre $C_D$ avec la transition \emph{imposée aux $x_{tr}$ d'XFOIL} (\code{free\_transition=False}). L'écart restant mesure l'erreur « couche limite + Squire-Young » seule, séparée de l'erreur de transition.
\end{enumerate}

\begin{resultat}
Pour l'optimisation, ce qui compte est moins l'erreur absolue que la \textbf{conservation du classement} des profils. La métrique de validation à présenter au jury est la corrélation de rang de Spearman entre notre $C_D$ et celui d'XFOIL sur un échantillon de la grille, en plus de l'erreur relative. Le test qui compte vraiment est : les deux modèles donnent-ils le même front de Pareto aérodynamique ?
\end{resultat}

\section{Ce que le modèle ne fait pas}
\begin{itemize}
\item \textbf{Pas de correction visqueuse de $C_L$} : le $C_L$ réel est 5 à 10~\% plus bas. Cela compte pour la contrainte $C_L\ge C_{L,min}$, et les profils à cambrure arrière sont les plus surestimés.
\item \textbf{Pas de décrochage} : la portance reste linéaire. Au-delà de 8 à $10^\circ$, le modèle sort de son domaine, ce que les drapeaux signalent.
\item \textbf{Incompressible} : pas d'effet de Mach. Une correction de Prandtl-Glauert ou de Kármán-Tsien sur $C_p$ serait possible en dessous du Mach critique.
\item \textbf{Transition par Michel} et non par $e^N$ : trop précoce à bas Reynolds.
\item \textbf{Bulbe laminaire} réduit à une transition instantanée.
\item \textbf{Méthode de Head} d'équilibre, sans effet de retard.
\end{itemize}

\section{Pistes, par rapport utilité / effort}
\begin{enumerate}
\item \textbf{Fait} : XFOIL installé, V5 lancé sur toute la grille (§\ref{sec:v5}). Prochaine étape : caler la transition sur XFOIL, la cause principale des écarts.
\item \textbf{Caler la transition} sur XFOIL (facteur sur le seuil de Michel, ou critère d'Arnal ou de Granville qui tient compte du gradient), ou proposer une option « transition forcée à 5~\% » : cas tout turbulent, plus régulier, donc plus favorable au modèle quadratique.
\item \textbf{Corriger $C_L$} par la vitesse de transpiration, une ou deux itérations.
\item \textbf{Kármán-Tsien} pour montrer l'effet de $M=0{,}5$ à $0{,}6$.
\end{enumerate}

\begin{projet}
Le message à porter devant le jury tient en trois phrases. Le modèle est volontairement basse fidélité, parce que le hackathon compare des méthodes d'optimisation en nombre de simulations, et que le front exact doit être calculable. Ses tendances sont physiques et ses erreurs sont identifiées et chiffrées, la principale étant la transition. Pour un cas industriel, seul le simulateur change : l'interface \code{AeroSimulator}, la grille, l'encodage et la boucle restent les mêmes.
\end{projet}

\exercices
\begin{exercice}
Deux modèles de traînée A et B sont comparés à XFOIL sur 30 profils. A a une erreur relative moyenne de 25~\%, mais une corrélation de rang de 0,97. B a une erreur moyenne de 8~\%, mais une corrélation de rang de 0,75. Lequel préférer pour l'optimisation du hackathon, et pourquoi ?
\end{exercice}
\begin{exercice}
Expliquer qualitativement pourquoi la vitesse de transpiration $V_n=\dd(U_e\delta^*)/\dd s$ diminue la portance calculée d'un profil cambré, en raisonnant sur l'épaisseur de déplacement de l'extrados et de l'intrados près du bord de fuite.
\end{exercice}

\chapter{V2 : rayon de bord d'attaque libre et XFOIL de référence}\label{ch:v2}
\iterbanner{V2}

\begin{evolution}{V2}{pourquoi ce chapitre}La V1 a montré deux limites du modèle de la note d'intention. Avec la famille NACA à 4 chiffres, l'épaisseur pilote à la fois la traînée et le rayon de bord d'attaque, donc la SER : il n'y a pas de vrai compromis. Et notre modèle rapide surestime $C_D$ de 27~\% par rapport à XFOIL, au point de déplacer l'optimum. La V2 change donc la géométrie et le simulateur de référence. Les chapitres~\ref{ch:intro} à \ref{ch:validation} restent valables pour le modèle rapide, conservé comme outil de développement.\end{evolution}

\section{La série NACA à 4 chiffres modifiée}
La série modifiée (Abbott et von Doenhoff, §6.4) garde la ligne moyenne à deux paraboles du chapitre~\ref{ch:naca}, mais remplace la loi d'épaisseur par deux polynômes raccordés à l'abscisse $T$ de l'épaisseur maximale, avec un nouveau paramètre : l'indice de rayon de bord d'attaque $I$. Pour un profil de 20~\% (demi-épaisseur 0,1 en $x=T$) :
\[
y_{20}(x)=\begin{cases}a_0\sqrt x+a_1x+a_2x^2+a_3x^3 & x\le T,\\ d_1(1-x)+d_2(1-x)^2+d_3(1-x)^3 & x\ge T,\end{cases}
\qquad y_t=\frac{t}{0{,}2}\,y_{20}.
\]
\begin{itemize}
\item $a_0=0{,}2969\,I/6$ fixe le nez. Le rayon de bord d'attaque vaut $r_{BA}=1{,}1019\,(tI/6)^2$ : à épaisseur donnée, il varie comme $I^2$. $I=6$ redonne le rayon du NACA standard.
\item $d_1=\dfrac{2{,}24-5{,}42T+12{,}3T^2}{10(1-0{,}878T)}$ fixe la pente au bord de fuite (formule de Ladson et Brooks, 1975), et $d_0=0$ ferme le bord de fuite.
\item $d_2$, $d_3$ sont choisis pour que $y_{20}(T)=0{,}1$ et $y_{20}'(T)=0$ ; puis $a_1$, $a_2$, $a_3$ pour que l'arc avant ait aussi $y=0{,}1$, $y'=0$ \emph{et la même courbure} que l'arc arrière en $x=T$. Deux petits systèmes linéaires suffisent.
\end{itemize}
Avec $I=6$ et $T=0{,}3$, on retrouve la loi du NACA standard à 0,1~\% de corde près. En V2, $T=0{,}3$ est fixé, et $I$ devient la quatrième variable de conception.

\begin{figure}[htbp]
\centering
% figure générée par tools/make_figures_v2.py : v2_formes
\begin{tikzpicture}
\begin{groupplot}[group style={group size=2 by 1,horizontal sep=1.4cm},title style={font=\small},tick label style={font=\scriptsize}]
\nextgroupplot[width=0.62\linewidth,height=4.2cm,axis equal image,xmin=-0.02,xmax=1.02,ymin=-0.08,ymax=0.12,
  xlabel={$x/c$},ylabel={$y/c$},title={NACA 2412 modifié, quatre rayons de bord d'attaque},
  legend style={draw=none,font=\scriptsize,at={(0.5,-0.45)},anchor=north},legend columns=2]
\addplot[corange,very thick] coordinates {(1,0) (0.99989,2.2648e-05) (0.99957,9.0569e-05) (0.99902,0.0002037) (0.99826,0.00036191) (0.99729,0.00056507) (0.99609,0.00081295) (0.99469,0.0011053) (0.99306,0.0014419) (0.99122,0.0018223) (0.98917,0.0022462) (0.98691,0.0027131) (0.98443,0.0032226) (0.98175,0.0037741) (0.97885,0.0043671) (0.97575,0.005001) (0.97243,0.0056751) (0.96892,0.0063887) (0.96519,0.0071412) (0.96127,0.0079316) (0.95714,0.0087592) (0.95282,0.0096231) (0.94829,0.010522) (0.94358,0.011456) (0.93866,0.012423) (0.93356,0.013423) (0.92826,0.014454) (0.92278,0.015515) (0.91712,0.016605) (0.91127,0.017723) (0.90524,0.018867) (0.89903,0.020037) (0.89265,0.021231) (0.88609,0.022448) (0.87937,0.023685) (0.87248,0.024943) (0.86542,0.026219) (0.85821,0.027512) (0.85083,0.02882) (0.84331,0.030143) (0.83563,0.031477) (0.8278,0.032823) (0.81982,0.034178) (0.81171,0.03554) (0.80346,0.036908) (0.79507,0.03828) (0.78655,0.039655) (0.77791,0.041031) (0.76914,0.042406) (0.76025,0.043778) (0.75124,0.045146) (0.74213,0.046509) (0.7329,0.047864) (0.72357,0.049209) (0.71414,0.050544) (0.70461,0.051865) (0.695,0.053173) (0.68529,0.054465) (0.6755,0.055739) (0.66563,0.056993) (0.65569,0.058227) (0.64568,0.059439) (0.6356,0.060626) (0.62546,0.061789) (0.61526,0.062924) (0.60501,0.064031) (0.59471,0.065108) (0.58436,0.066155) (0.57398,0.067169) (0.56356,0.06815) (0.55312,0.069096) (0.54264,0.070006) (0.53215,0.070879) (0.52164,0.071714) (0.51112,0.072511) (0.5006,0.073268) (0.49007,0.073984) (0.47955,0.074659) (0.46903,0.075293) (0.45852,0.075884) (0.44804,0.076431) (0.43757,0.076936) (0.42713,0.077397) (0.41672,0.077813) (0.40635,0.078186) (0.39599,0.078513) (0.38562,0.078784) (0.3753,0.078996) (0.36503,0.07915) (0.35482,0.079245) (0.34468,0.079283) (0.3346,0.079265) (0.3246,0.07919) (0.31468,0.079061) (0.30484,0.078878) (0.29508,0.078642) (0.28542,0.078353) (0.27585,0.078004) (0.26639,0.077593) (0.25702,0.077115) (0.24777,0.076566) (0.23864,0.075944) (0.22962,0.075246) (0.22073,0.074469) (0.21196,0.073614) (0.20333,0.072678) (0.19483,0.071661) (0.18647,0.070562) (0.17826,0.069384) (0.1702,0.068125) (0.16229,0.066788) (0.15453,0.065375) (0.14694,0.063887) (0.13951,0.062328) (0.13225,0.060701) (0.12515,0.059009) (0.11823,0.057255) (0.11149,0.055446) (0.10493,0.053583) (0.098544,0.051674) (0.092346,0.049722) (0.086336,0.047734) (0.080515,0.045714) (0.074885,0.043669) (0.069449,0.041604) (0.064208,0.039526) (0.059164,0.03744) (0.054319,0.035353) (0.049675,0.033271) (0.045232,0.0312) (0.040992,0.029148) (0.036957,0.027118) (0.033127,0.025119) (0.029503,0.023156) (0.026086,0.021234) (0.022877,0.019359) (0.019877,0.017538) (0.017086,0.015775) (0.014504,0.014074) (0.012133,0.012442) (0.0099725,0.010882) (0.008023,0.0093986) (0.0062846,0.0079955) (0.0047577,0.0066763) (0.0034425,0.0054439) (0.002339,0.0043012) (0.0014474,0.0032504) (0.00076767,0.0022936) (0.00029989,0.0014323) (4.4022e-05,0.00066751) (0,0) (0.00017529,-0.00064558) (0.00057728,-0.0013446) (0.0012056,-0.0020965) (0.0020598,-0.0029005) (0.0031391,-0.0037552) (0.0044428,-0.0046592) (0.0059699,-0.0056107) (0.0077193,-0.0066074) (0.0096898,-0.007647) (0.01188,-0.0087267) (0.014288,-0.0098435) (0.016913,-0.010994) (0.019752,-0.012176) (0.022804,-0.013384) (0.026067,-0.014615) (0.029538,-0.015865) (0.033215,-0.01713) (0.037097,-0.018405) (0.04118,-0.019686) (0.045462,-0.020969) (0.049941,-0.022249) (0.054614,-0.023522) (0.059477,-0.024783) (0.064529,-0.026027) (0.069767,-0.02725) (0.075187,-0.028448) (0.080787,-0.029616) (0.086564,-0.030751) (0.092515,-0.031848) (0.098637,-0.032904) (0.10493,-0.033915) (0.11138,-0.034877) (0.118,-0.035788) (0.12477,-0.036646) (0.1317,-0.037446) (0.13879,-0.038189) (0.14602,-0.038871) (0.1534,-0.039491) (0.16092,-0.040049) (0.16858,-0.040543) (0.17638,-0.040974) (0.18432,-0.041342) (0.19238,-0.041648) (0.20057,-0.041892) (0.20889,-0.042076) (0.21733,-0.042203) (0.22588,-0.042275) (0.23455,-0.042294) (0.24334,-0.042265) (0.25223,-0.042191) (0.26122,-0.042076) (0.27032,-0.041926) (0.27951,-0.041745) (0.2888,-0.041539) (0.29818,-0.041314) (0.30765,-0.041076) (0.3172,-0.040828) (0.32683,-0.04057) (0.33653,-0.040303) (0.3463,-0.040026) (0.35615,-0.039741) (0.36605,-0.039446) (0.37601,-0.039144) (0.38603,-0.038834) (0.3961,-0.038517) (0.40627,-0.03819) (0.41651,-0.037844) (0.42679,-0.037477) (0.4371,-0.037091) (0.44744,-0.036685) (0.4578,-0.036259) (0.46818,-0.035816) (0.47858,-0.035354) (0.48899,-0.034875) (0.4994,-0.034379) (0.50982,-0.033867) (0.52023,-0.03334) (0.53064,-0.032797) (0.54103,-0.032241) (0.55141,-0.031672) (0.56177,-0.03109) (0.5721,-0.030496) (0.58241,-0.029892) (0.59268,-0.029277) (0.60291,-0.028653) (0.61309,-0.028021) (0.62323,-0.027381) (0.63332,-0.026734) (0.64335,-0.026082) (0.65333,-0.025424) (0.66323,-0.024763) (0.67307,-0.024098) (0.68283,-0.02343) (0.69252,-0.022761) (0.70212,-0.022091) (0.71164,-0.021421) (0.72106,-0.020752) (0.7304,-0.020085) (0.73963,-0.01942) (0.74876,-0.018757) (0.75778,-0.018099) (0.76669,-0.017446) (0.77549,-0.016797) (0.78416,-0.016155) (0.79272,-0.015519) (0.80114,-0.01489) (0.80944,-0.01427) (0.8176,-0.013658) (0.82562,-0.013054) (0.8335,-0.01246) (0.84124,-0.011877) (0.84883,-0.011304) (0.85626,-0.010742) (0.86354,-0.010191) (0.87067,-0.009652) (0.87763,-0.0091254) (0.88442,-0.0086114) (0.89104,-0.0081104) (0.8975,-0.0076226) (0.90378,-0.0071484) (0.90988,-0.0066881) (0.91581,-0.0062418) (0.92155,-0.0058099) (0.9271,-0.0053925) (0.93247,-0.0049899) (0.93764,-0.0046021) (0.94263,-0.0042295) (0.94742,-0.003872) (0.95201,-0.0035299) (0.9564,-0.0032032) (0.96059,-0.0028921) (0.96458,-0.0025966) (0.96837,-0.0023168) (0.97194,-0.0020528) (0.97531,-0.0018046) (0.97847,-0.0015723) (0.98142,-0.0013559) (0.98415,-0.0011555) (0.98667,-0.00097106) (0.98897,-0.0008026) (0.99106,-0.00065015) (0.99293,-0.00051373) (0.99459,-0.00039334) (0.99602,-0.00028899) (0.99724,-0.00020069) (0.99823,-0.00012845) (0.999,-7.2252e-05) (0.99956,-3.2112e-05) (0.99989,-8.0281e-06) (1,0)};\addlegendentry{$I=2$, $r_{BA}=0.18$~\%}
\addplot[cgreen,very thick] coordinates {(1,0) (0.99989,2.2648e-05) (0.99957,9.0569e-05) (0.99902,0.0002037) (0.99826,0.00036191) (0.99729,0.00056507) (0.99609,0.00081295) (0.99469,0.0011053) (0.99306,0.0014419) (0.99122,0.0018223) (0.98917,0.0022462) (0.98691,0.0027131) (0.98443,0.0032226) (0.98175,0.0037741) (0.97885,0.0043671) (0.97575,0.005001) (0.97243,0.0056751) (0.96892,0.0063887) (0.96519,0.0071412) (0.96127,0.0079316) (0.95714,0.0087592) (0.95282,0.0096231) (0.94829,0.010522) (0.94358,0.011456) (0.93866,0.012423) (0.93356,0.013423) (0.92826,0.014454) (0.92278,0.015515) (0.91712,0.016605) (0.91127,0.017723) (0.90524,0.018867) (0.89903,0.020037) (0.89265,0.021231) (0.88609,0.022448) (0.87937,0.023685) (0.87248,0.024943) (0.86542,0.026219) (0.85821,0.027512) (0.85083,0.02882) (0.84331,0.030143) (0.83563,0.031477) (0.8278,0.032823) (0.81982,0.034178) (0.81171,0.03554) (0.80346,0.036908) (0.79507,0.03828) (0.78655,0.039655) (0.77791,0.041031) (0.76914,0.042406) (0.76025,0.043778) (0.75124,0.045146) (0.74213,0.046509) (0.7329,0.047864) (0.72357,0.049209) (0.71414,0.050544) (0.70461,0.051865) (0.695,0.053173) (0.68529,0.054465) (0.6755,0.055739) (0.66563,0.056993) (0.65569,0.058227) (0.64568,0.059439) (0.6356,0.060626) (0.62546,0.061789) (0.61526,0.062924) (0.60501,0.064031) (0.59471,0.065108) (0.58436,0.066155) (0.57398,0.067169) (0.56356,0.06815) (0.55312,0.069096) (0.54264,0.070006) (0.53215,0.070879) (0.52164,0.071714) (0.51112,0.072511) (0.5006,0.073268) (0.49007,0.073984) (0.47955,0.074659) (0.46903,0.075293) (0.45852,0.075884) (0.44804,0.076431) (0.43757,0.076936) (0.42713,0.077397) (0.41672,0.077813) (0.40635,0.078186) (0.39599,0.078513) (0.38562,0.078784) (0.3753,0.078996) (0.36503,0.07915) (0.35482,0.079245) (0.34468,0.079283) (0.3346,0.079265) (0.3246,0.07919) (0.31468,0.079061) (0.30484,0.078878) (0.29508,0.078642) (0.28542,0.078353) (0.27585,0.078008) (0.26638,0.077605) (0.25702,0.07714) (0.24777,0.076612) (0.23863,0.076019) (0.22961,0.075359) (0.22072,0.074633) (0.21195,0.073837) (0.20331,0.072973) (0.19481,0.072041) (0.18645,0.071039) (0.17823,0.069969) (0.17016,0.068832) (0.16224,0.067628) (0.15447,0.066359) (0.14687,0.065027) (0.13943,0.063634) (0.13215,0.062181) (0.12504,0.060672) (0.1181,0.059109) (0.11134,0.057495) (0.10476,0.055833) (0.098361,0.054127) (0.092144,0.052379) (0.086114,0.050594) (0.080272,0.048775) (0.074622,0.046927) (0.069166,0.045052) (0.063905,0.043156) (0.058842,0.041243) (0.053979,0.039316) (0.049317,0.03738) (0.044859,0.035439) (0.040604,0.033497) (0.036556,0.031558) (0.032715,0.029627) (0.029083,0.027707) (0.025661,0.025803) (0.022449,0.023918) (0.019449,0.022056) (0.016662,0.02022) (0.014087,0.018413) (0.011727,0.01664) (0.0095815,0.014902) (0.0076509,0.013204) (0.0059358,0.011546) (0.0044365,0.009932) (0.0031534,0.0083637) (0.0020865,0.006843) (0.0012362,0.0053715) (0.00060239,0.0039505) (0.00018514,0.0025811) (-1.5618e-05,0.0012641) (0,0) (0.00023493,-0.0012421) (0.00069203,-0.0024934) (0.0013709,-0.0037534) (0.0022709,-0.0050215) (0.0033916,-0.0062971) (0.0047319,-0.0075791) (0.0062911,-0.0088664) (0.0080682,-0.010158) (0.010062,-0.011452) (0.012271,-0.012747) (0.014694,-0.014041) (0.017329,-0.015333) (0.020176,-0.016621) (0.023232,-0.017901) (0.026494,-0.019173) (0.029963,-0.020434) (0.033635,-0.021681) (0.037508,-0.022913) (0.041581,-0.024126) (0.04585,-0.025318) (0.050314,-0.026487) (0.054971,-0.027631) (0.059817,-0.028745) (0.064851,-0.02983) (0.070069,-0.030881) (0.07547,-0.031896) (0.08105,-0.032874) (0.086806,-0.033812) (0.092737,-0.034708) (0.098839,-0.03556) (0.10511,-0.036367) (0.11154,-0.037127) (0.11814,-0.037838) (0.1249,-0.0385) (0.13181,-0.03911) (0.13888,-0.039669) (0.1461,-0.040177) (0.15347,-0.040631) (0.16098,-0.041033) (0.16863,-0.041383) (0.17642,-0.041681) (0.18435,-0.041928) (0.1924,-0.042124) (0.20059,-0.042272) (0.2089,-0.042372) (0.21734,-0.042427) (0.22589,-0.042438) (0.23456,-0.042408) (0.24334,-0.04234) (0.25223,-0.042237) (0.26122,-0.042101) (0.27032,-0.041937) (0.27951,-0.041749) (0.2888,-0.041539) (0.29818,-0.041314) (0.30765,-0.041076) (0.3172,-0.040828) (0.32683,-0.04057) (0.33653,-0.040303) (0.3463,-0.040026) (0.35615,-0.039741) (0.36605,-0.039446) (0.37601,-0.039144) (0.38603,-0.038834) (0.3961,-0.038517) (0.40627,-0.03819) (0.41651,-0.037844) (0.42679,-0.037477) (0.4371,-0.037091) (0.44744,-0.036685) (0.4578,-0.036259) (0.46818,-0.035816) (0.47858,-0.035354) (0.48899,-0.034875) (0.4994,-0.034379) (0.50982,-0.033867) (0.52023,-0.03334) (0.53064,-0.032797) (0.54103,-0.032241) (0.55141,-0.031672) (0.56177,-0.03109) (0.5721,-0.030496) (0.58241,-0.029892) (0.59268,-0.029277) (0.60291,-0.028653) (0.61309,-0.028021) (0.62323,-0.027381) (0.63332,-0.026734) (0.64335,-0.026082) (0.65333,-0.025424) (0.66323,-0.024763) (0.67307,-0.024098) (0.68283,-0.02343) (0.69252,-0.022761) (0.70212,-0.022091) (0.71164,-0.021421) (0.72106,-0.020752) (0.7304,-0.020085) (0.73963,-0.01942) (0.74876,-0.018757) (0.75778,-0.018099) (0.76669,-0.017446) (0.77549,-0.016797) (0.78416,-0.016155) (0.79272,-0.015519) (0.80114,-0.01489) (0.80944,-0.01427) (0.8176,-0.013658) (0.82562,-0.013054) (0.8335,-0.01246) (0.84124,-0.011877) (0.84883,-0.011304) (0.85626,-0.010742) (0.86354,-0.010191) (0.87067,-0.009652) (0.87763,-0.0091254) (0.88442,-0.0086114) (0.89104,-0.0081104) (0.8975,-0.0076226) (0.90378,-0.0071484) (0.90988,-0.0066881) (0.91581,-0.0062418) (0.92155,-0.0058099) (0.9271,-0.0053925) (0.93247,-0.0049899) (0.93764,-0.0046021) (0.94263,-0.0042295) (0.94742,-0.003872) (0.95201,-0.0035299) (0.9564,-0.0032032) (0.96059,-0.0028921) (0.96458,-0.0025966) (0.96837,-0.0023168) (0.97194,-0.0020528) (0.97531,-0.0018046) (0.97847,-0.0015723) (0.98142,-0.0013559) (0.98415,-0.0011555) (0.98667,-0.00097106) (0.98897,-0.0008026) (0.99106,-0.00065015) (0.99293,-0.00051373) (0.99459,-0.00039334) (0.99602,-0.00028899) (0.99724,-0.00020069) (0.99823,-0.00012845) (0.999,-7.2252e-05) (0.99956,-3.2112e-05) (0.99989,-8.0281e-06) (1,0)};\addlegendentry{$I=4$, $r_{BA}=0.71$~\%}
\addplot[cblue,very thick] coordinates {(1,0) (0.99989,2.2648e-05) (0.99957,9.0569e-05) (0.99902,0.0002037) (0.99826,0.00036191) (0.99729,0.00056507) (0.99609,0.00081295) (0.99469,0.0011053) (0.99306,0.0014419) (0.99122,0.0018223) (0.98917,0.0022462) (0.98691,0.0027131) (0.98443,0.0032226) (0.98175,0.0037741) (0.97885,0.0043671) (0.97575,0.005001) (0.97243,0.0056751) (0.96892,0.0063887) (0.96519,0.0071412) (0.96127,0.0079316) (0.95714,0.0087592) (0.95282,0.0096231) (0.94829,0.010522) (0.94358,0.011456) (0.93866,0.012423) (0.93356,0.013423) (0.92826,0.014454) (0.92278,0.015515) (0.91712,0.016605) (0.91127,0.017723) (0.90524,0.018867) (0.89903,0.020037) (0.89265,0.021231) (0.88609,0.022448) (0.87937,0.023685) (0.87248,0.024943) (0.86542,0.026219) (0.85821,0.027512) (0.85083,0.02882) (0.84331,0.030143) (0.83563,0.031477) (0.8278,0.032823) (0.81982,0.034178) (0.81171,0.03554) (0.80346,0.036908) (0.79507,0.03828) (0.78655,0.039655) (0.77791,0.041031) (0.76914,0.042406) (0.76025,0.043778) (0.75124,0.045146) (0.74213,0.046509) (0.7329,0.047864) (0.72357,0.049209) (0.71414,0.050544) (0.70461,0.051865) (0.695,0.053173) (0.68529,0.054465) (0.6755,0.055739) (0.66563,0.056993) (0.65569,0.058227) (0.64568,0.059439) (0.6356,0.060626) (0.62546,0.061789) (0.61526,0.062924) (0.60501,0.064031) (0.59471,0.065108) (0.58436,0.066155) (0.57398,0.067169) (0.56356,0.06815) (0.55312,0.069096) (0.54264,0.070006) (0.53215,0.070879) (0.52164,0.071714) (0.51112,0.072511) (0.5006,0.073268) (0.49007,0.073984) (0.47955,0.074659) (0.46903,0.075293) (0.45852,0.075884) (0.44804,0.076431) (0.43757,0.076936) (0.42713,0.077397) (0.41672,0.077813) (0.40635,0.078186) (0.39599,0.078513) (0.38562,0.078784) (0.3753,0.078996) (0.36503,0.07915) (0.35482,0.079245) (0.34468,0.079283) (0.3346,0.079265) (0.3246,0.07919) (0.31468,0.079061) (0.30484,0.078878) (0.29508,0.078642) (0.28542,0.078354) (0.27585,0.078013) (0.26638,0.077616) (0.25702,0.077165) (0.24777,0.076658) (0.23863,0.076094) (0.22961,0.075473) (0.22071,0.074796) (0.21194,0.074061) (0.2033,0.073269) (0.19479,0.07242) (0.18642,0.071515) (0.1782,0.070554) (0.17012,0.069538) (0.16219,0.068468) (0.15441,0.067344) (0.1468,0.066167) (0.13934,0.06494) (0.13205,0.063662) (0.12493,0.062336) (0.11798,0.060963) (0.1112,0.059545) (0.1046,0.058083) (0.098178,0.056579) (0.091942,0.055036) (0.085892,0.053454) (0.08003,0.051836) (0.07436,0.050184) (0.068884,0.0485) (0.063603,0.046787) (0.058521,0.045046) (0.053639,0.043279) (0.04896,0.041488) (0.044485,0.039677) (0.040216,0.037846) (0.036156,0.035998) (0.032304,0.034135) (0.028664,0.032259) (0.025236,0.030372) (0.022021,0.028476) (0.019021,0.026573) (0.016237,0.024665) (0.01367,0.022752) (0.011321,0.020838) (0.0091904,0.018923) (0.0072788,0.017008) (0.005587,0.015096) (0.0041153,0.013188) (0.0028643,0.011284) (0.0018341,0.0093849) (0.001025,0.0074925) (0.00043712,0.0056073) (7.0386e-05,0.0037299) (-7.5258e-05,0.0018606) (0,0) (0.00029457,-0.0018387) (0.00080678,-0.0036422) (0.0015362,-0.0054103) (0.0024821,-0.0071426) (0.003644,-0.0088389) (0.005021,-0.010499) (0.0066124,-0.012122) (0.008417,-0.013708) (0.010434,-0.015257) (0.012662,-0.016767) (0.0151,-0.018239) (0.017746,-0.019672) (0.0206,-0.021066) (0.023659,-0.022419) (0.026922,-0.023732) (0.030388,-0.025003) (0.034054,-0.026233) (0.037919,-0.027421) (0.041981,-0.028566) (0.046238,-0.029668) (0.050688,-0.030726) (0.055328,-0.031739) (0.060157,-0.032708) (0.065172,-0.033632) (0.070371,-0.034511) (0.075752,-0.035344) (0.081312,-0.036132) (0.087048,-0.036873) (0.092959,-0.037568) (0.099041,-0.038217) (0.10529,-0.03882) (0.11171,-0.039377) (0.11829,-0.039888) (0.12503,-0.040353) (0.13193,-0.040774) (0.13898,-0.04115) (0.14619,-0.041482) (0.15354,-0.041771) (0.16104,-0.042018) (0.16868,-0.042223) (0.17646,-0.042388) (0.18438,-0.042513) (0.19243,-0.042601) (0.20061,-0.042652) (0.20892,-0.042668) (0.21735,-0.04265) (0.2259,-0.042601) (0.23456,-0.042522) (0.24334,-0.042415) (0.25223,-0.042283) (0.26122,-0.042127) (0.27032,-0.041949) (0.27951,-0.041753) (0.2888,-0.04154) (0.29818,-0.041314) (0.30765,-0.041076) (0.3172,-0.040828) (0.32683,-0.04057) (0.33653,-0.040303) (0.3463,-0.040026) (0.35615,-0.039741) (0.36605,-0.039446) (0.37601,-0.039144) (0.38603,-0.038834) (0.3961,-0.038517) (0.40627,-0.03819) (0.41651,-0.037844) (0.42679,-0.037477) (0.4371,-0.037091) (0.44744,-0.036685) (0.4578,-0.036259) (0.46818,-0.035816) (0.47858,-0.035354) (0.48899,-0.034875) (0.4994,-0.034379) (0.50982,-0.033867) (0.52023,-0.03334) (0.53064,-0.032797) (0.54103,-0.032241) (0.55141,-0.031672) (0.56177,-0.03109) (0.5721,-0.030496) (0.58241,-0.029892) (0.59268,-0.029277) (0.60291,-0.028653) (0.61309,-0.028021) (0.62323,-0.027381) (0.63332,-0.026734) (0.64335,-0.026082) (0.65333,-0.025424) (0.66323,-0.024763) (0.67307,-0.024098) (0.68283,-0.02343) (0.69252,-0.022761) (0.70212,-0.022091) (0.71164,-0.021421) (0.72106,-0.020752) (0.7304,-0.020085) (0.73963,-0.01942) (0.74876,-0.018757) (0.75778,-0.018099) (0.76669,-0.017446) (0.77549,-0.016797) (0.78416,-0.016155) (0.79272,-0.015519) (0.80114,-0.01489) (0.80944,-0.01427) (0.8176,-0.013658) (0.82562,-0.013054) (0.8335,-0.01246) (0.84124,-0.011877) (0.84883,-0.011304) (0.85626,-0.010742) (0.86354,-0.010191) (0.87067,-0.009652) (0.87763,-0.0091254) (0.88442,-0.0086114) (0.89104,-0.0081104) (0.8975,-0.0076226) (0.90378,-0.0071484) (0.90988,-0.0066881) (0.91581,-0.0062418) (0.92155,-0.0058099) (0.9271,-0.0053925) (0.93247,-0.0049899) (0.93764,-0.0046021) (0.94263,-0.0042295) (0.94742,-0.003872) (0.95201,-0.0035299) (0.9564,-0.0032032) (0.96059,-0.0028921) (0.96458,-0.0025966) (0.96837,-0.0023168) (0.97194,-0.0020528) (0.97531,-0.0018046) (0.97847,-0.0015723) (0.98142,-0.0013559) (0.98415,-0.0011555) (0.98667,-0.00097106) (0.98897,-0.0008026) (0.99106,-0.00065015) (0.99293,-0.00051373) (0.99459,-0.00039334) (0.99602,-0.00028899) (0.99724,-0.00020069) (0.99823,-0.00012845) (0.999,-7.2252e-05) (0.99956,-3.2112e-05) (0.99989,-8.0281e-06) (1,0)};\addlegendentry{$I=6$, $r_{BA}=1.59$~\%}
\addplot[cmuted,very thick] coordinates {(1,0) (0.99989,2.2648e-05) (0.99957,9.0569e-05) (0.99902,0.0002037) (0.99826,0.00036191) (0.99729,0.00056507) (0.99609,0.00081295) (0.99469,0.0011053) (0.99306,0.0014419) (0.99122,0.0018223) (0.98917,0.0022462) (0.98691,0.0027131) (0.98443,0.0032226) (0.98175,0.0037741) (0.97885,0.0043671) (0.97575,0.005001) (0.97243,0.0056751) (0.96892,0.0063887) (0.96519,0.0071412) (0.96127,0.0079316) (0.95714,0.0087592) (0.95282,0.0096231) (0.94829,0.010522) (0.94358,0.011456) (0.93866,0.012423) (0.93356,0.013423) (0.92826,0.014454) (0.92278,0.015515) (0.91712,0.016605) (0.91127,0.017723) (0.90524,0.018867) (0.89903,0.020037) (0.89265,0.021231) (0.88609,0.022448) (0.87937,0.023685) (0.87248,0.024943) (0.86542,0.026219) (0.85821,0.027512) (0.85083,0.02882) (0.84331,0.030143) (0.83563,0.031477) (0.8278,0.032823) (0.81982,0.034178) (0.81171,0.03554) (0.80346,0.036908) (0.79507,0.03828) (0.78655,0.039655) (0.77791,0.041031) (0.76914,0.042406) (0.76025,0.043778) (0.75124,0.045146) (0.74213,0.046509) (0.7329,0.047864) (0.72357,0.049209) (0.71414,0.050544) (0.70461,0.051865) (0.695,0.053173) (0.68529,0.054465) (0.6755,0.055739) (0.66563,0.056993) (0.65569,0.058227) (0.64568,0.059439) (0.6356,0.060626) (0.62546,0.061789) (0.61526,0.062924) (0.60501,0.064031) (0.59471,0.065108) (0.58436,0.066155) (0.57398,0.067169) (0.56356,0.06815) (0.55312,0.069096) (0.54264,0.070006) (0.53215,0.070879) (0.52164,0.071714) (0.51112,0.072511) (0.5006,0.073268) (0.49007,0.073984) (0.47955,0.074659) (0.46903,0.075293) (0.45852,0.075884) (0.44804,0.076431) (0.43757,0.076936) (0.42713,0.077397) (0.41672,0.077813) (0.40635,0.078186) (0.39599,0.078513) (0.38562,0.078784) (0.3753,0.078996) (0.36503,0.07915) (0.35482,0.079245) (0.34468,0.079283) (0.3346,0.079265) (0.3246,0.07919) (0.31468,0.079061) (0.30484,0.078878) (0.29508,0.078642) (0.28542,0.078355) (0.27585,0.078017) (0.26638,0.077628) (0.25702,0.07719) (0.24777,0.076704) (0.23863,0.076169) (0.22961,0.075587) (0.2207,0.074959) (0.21193,0.074284) (0.20328,0.073564) (0.19477,0.0728) (0.1864,0.071992) (0.17816,0.07114) (0.17008,0.070245) (0.16214,0.069307) (0.15435,0.068328) (0.14673,0.067307) (0.13926,0.066245) (0.13195,0.065143) (0.12481,0.064) (0.11785,0.062817) (0.11105,0.061595) (0.10443,0.060333) (0.097995,0.059032) (0.09174,0.057692) (0.08567,0.056314) (0.079788,0.054897) (0.074098,0.053442) (0.068601,0.051949) (0.063301,0.050418) (0.0582,0.048848) (0.053299,0.047242) (0.048603,0.045597) (0.044112,0.043915) (0.039828,0.042195) (0.035755,0.040438) (0.031893,0.038643) (0.028244,0.036811) (0.02481,0.034942) (0.021593,0.033035) (0.018594,0.031091) (0.015813,0.02911) (0.013254,0.027091) (0.010915,0.025036) (0.0087994,0.022943) (0.0069068,0.020813) (0.0052381,0.018647) (0.0037941,0.016443) (0.0025752,0.014203) (0.0015817,0.011927) (0.00081385,0.0096136) (0.00027184,0.0072642) (-4.4368e-05,0.0048787) (-0.0001349,0.0024572) (0,0) (0.00035421,-0.0024353) (0.00092154,-0.004791) (0.0017014,-0.0070671) (0.0026933,-0.0092636) (0.0038964,-0.011381) (0.0053101,-0.013419) (0.0069336,-0.015378) (0.0087658,-0.017259) (0.010806,-0.019062) (0.013053,-0.020788) (0.015506,-0.022437) (0.018163,-0.024011) (0.021024,-0.025511) (0.024087,-0.026937) (0.02735,-0.02829) (0.030813,-0.029573) (0.034474,-0.030785) (0.038331,-0.031929) (0.042382,-0.033006) (0.046626,-0.034017) (0.051061,-0.034964) (0.055686,-0.035848) (0.060497,-0.036671) (0.065494,-0.037435) (0.070674,-0.038142) (0.076035,-0.038793) (0.081574,-0.039389) (0.087291,-0.039934) (0.093181,-0.040428) (0.099243,-0.040874) (0.10547,-0.041272) (0.11187,-0.041626) (0.11844,-0.041937) (0.12516,-0.042207) (0.13204,-0.042438) (0.13908,-0.042631) (0.14627,-0.042788) (0.15361,-0.042911) (0.1611,-0.043002) (0.16873,-0.043063) (0.1765,-0.043094) (0.18441,-0.043098) (0.19245,-0.043077) (0.20063,-0.043031) (0.20893,-0.042963) (0.21736,-0.042873) (0.2259,-0.042764) (0.23457,-0.042636) (0.24335,-0.04249) (0.25223,-0.042329) (0.26122,-0.042152) (0.27032,-0.041961) (0.27951,-0.041757) (0.2888,-0.041541) (0.29818,-0.041314) (0.30765,-0.041076) (0.3172,-0.040828) (0.32683,-0.04057) (0.33653,-0.040303) (0.3463,-0.040026) (0.35615,-0.039741) (0.36605,-0.039446) (0.37601,-0.039144) (0.38603,-0.038834) (0.3961,-0.038517) (0.40627,-0.03819) (0.41651,-0.037844) (0.42679,-0.037477) (0.4371,-0.037091) (0.44744,-0.036685) (0.4578,-0.036259) (0.46818,-0.035816) (0.47858,-0.035354) (0.48899,-0.034875) (0.4994,-0.034379) (0.50982,-0.033867) (0.52023,-0.03334) (0.53064,-0.032797) (0.54103,-0.032241) (0.55141,-0.031672) (0.56177,-0.03109) (0.5721,-0.030496) (0.58241,-0.029892) (0.59268,-0.029277) (0.60291,-0.028653) (0.61309,-0.028021) (0.62323,-0.027381) (0.63332,-0.026734) (0.64335,-0.026082) (0.65333,-0.025424) (0.66323,-0.024763) (0.67307,-0.024098) (0.68283,-0.02343) (0.69252,-0.022761) (0.70212,-0.022091) (0.71164,-0.021421) (0.72106,-0.020752) (0.7304,-0.020085) (0.73963,-0.01942) (0.74876,-0.018757) (0.75778,-0.018099) (0.76669,-0.017446) (0.77549,-0.016797) (0.78416,-0.016155) (0.79272,-0.015519) (0.80114,-0.01489) (0.80944,-0.01427) (0.8176,-0.013658) (0.82562,-0.013054) (0.8335,-0.01246) (0.84124,-0.011877) (0.84883,-0.011304) (0.85626,-0.010742) (0.86354,-0.010191) (0.87067,-0.009652) (0.87763,-0.0091254) (0.88442,-0.0086114) (0.89104,-0.0081104) (0.8975,-0.0076226) (0.90378,-0.0071484) (0.90988,-0.0066881) (0.91581,-0.0062418) (0.92155,-0.0058099) (0.9271,-0.0053925) (0.93247,-0.0049899) (0.93764,-0.0046021) (0.94263,-0.0042295) (0.94742,-0.003872) (0.95201,-0.0035299) (0.9564,-0.0032032) (0.96059,-0.0028921) (0.96458,-0.0025966) (0.96837,-0.0023168) (0.97194,-0.0020528) (0.97531,-0.0018046) (0.97847,-0.0015723) (0.98142,-0.0013559) (0.98415,-0.0011555) (0.98667,-0.00097106) (0.98897,-0.0008026) (0.99106,-0.00065015) (0.99293,-0.00051373) (0.99459,-0.00039334) (0.99602,-0.00028899) (0.99724,-0.00020069) (0.99823,-0.00012845) (0.999,-7.2252e-05) (0.99956,-3.2112e-05) (0.99989,-8.0281e-06) (1,0)};\addlegendentry{$I=8$, $r_{BA}=2.82$~\%}
\nextgroupplot[width=0.36\linewidth,height=5cm,axis equal image,xmin=-0.002,xmax=0.05,ymin=-0.035,ymax=0.04,
  xlabel={$x/c$},title={zoom sur le nez},xticklabel style={/pgf/number format/fixed}]
\addplot[corange,very thick] coordinates {(0.059164,0.03744) (0.054319,0.035353) (0.049675,0.033271) (0.045232,0.0312) (0.040992,0.029148) (0.036957,0.027118) (0.033127,0.025119) (0.029503,0.023156) (0.026086,0.021234) (0.022877,0.019359) (0.019877,0.017538) (0.017086,0.015775) (0.014504,0.014074) (0.012133,0.012442) (0.0099725,0.010882) (0.008023,0.0093986) (0.0062846,0.0079955) (0.0047577,0.0066763) (0.0034425,0.0054439) (0.002339,0.0043012) (0.0014474,0.0032504) (0.00076767,0.0022936) (0.00029989,0.0014323) (4.4022e-05,0.00066751) (0,0) (0.00017529,-0.00064558) (0.00057728,-0.0013446) (0.0012056,-0.0020965) (0.0020598,-0.0029005) (0.0031391,-0.0037552) (0.0044428,-0.0046592) (0.0059699,-0.0056107) (0.0077193,-0.0066074) (0.0096898,-0.007647) (0.01188,-0.0087267) (0.014288,-0.0098435) (0.016913,-0.010994) (0.019752,-0.012176) (0.022804,-0.013384) (0.026067,-0.014615) (0.029538,-0.015865) (0.033215,-0.01713) (0.037097,-0.018405) (0.04118,-0.019686) (0.045462,-0.020969) (0.049941,-0.022249) (0.054614,-0.023522) (0.059477,-0.024783)};
\addplot[cgreen,very thick] coordinates {(0.058842,0.041243) (0.053979,0.039316) (0.049317,0.03738) (0.044859,0.035439) (0.040604,0.033497) (0.036556,0.031558) (0.032715,0.029627) (0.029083,0.027707) (0.025661,0.025803) (0.022449,0.023918) (0.019449,0.022056) (0.016662,0.02022) (0.014087,0.018413) (0.011727,0.01664) (0.0095815,0.014902) (0.0076509,0.013204) (0.0059358,0.011546) (0.0044365,0.009932) (0.0031534,0.0083637) (0.0020865,0.006843) (0.0012362,0.0053715) (0.00060239,0.0039505) (0.00018514,0.0025811) (-1.5618e-05,0.0012641) (0,0) (0.00023493,-0.0012421) (0.00069203,-0.0024934) (0.0013709,-0.0037534) (0.0022709,-0.0050215) (0.0033916,-0.0062971) (0.0047319,-0.0075791) (0.0062911,-0.0088664) (0.0080682,-0.010158) (0.010062,-0.011452) (0.012271,-0.012747) (0.014694,-0.014041) (0.017329,-0.015333) (0.020176,-0.016621) (0.023232,-0.017901) (0.026494,-0.019173) (0.029963,-0.020434) (0.033635,-0.021681) (0.037508,-0.022913) (0.041581,-0.024126) (0.04585,-0.025318) (0.050314,-0.026487) (0.054971,-0.027631) (0.059817,-0.028745)};
\addplot[cblue,very thick] coordinates {(0.058521,0.045046) (0.053639,0.043279) (0.04896,0.041488) (0.044485,0.039677) (0.040216,0.037846) (0.036156,0.035998) (0.032304,0.034135) (0.028664,0.032259) (0.025236,0.030372) (0.022021,0.028476) (0.019021,0.026573) (0.016237,0.024665) (0.01367,0.022752) (0.011321,0.020838) (0.0091904,0.018923) (0.0072788,0.017008) (0.005587,0.015096) (0.0041153,0.013188) (0.0028643,0.011284) (0.0018341,0.0093849) (0.001025,0.0074925) (0.00043712,0.0056073) (7.0386e-05,0.0037299) (-7.5258e-05,0.0018606) (0,0) (0.00029457,-0.0018387) (0.00080678,-0.0036422) (0.0015362,-0.0054103) (0.0024821,-0.0071426) (0.003644,-0.0088389) (0.005021,-0.010499) (0.0066124,-0.012122) (0.008417,-0.013708) (0.010434,-0.015257) (0.012662,-0.016767) (0.0151,-0.018239) (0.017746,-0.019672) (0.0206,-0.021066) (0.023659,-0.022419) (0.026922,-0.023732) (0.030388,-0.025003) (0.034054,-0.026233) (0.037919,-0.027421) (0.041981,-0.028566) (0.046238,-0.029668) (0.050688,-0.030726) (0.055328,-0.031739)};
\addplot[cmuted,very thick] coordinates {(0.0582,0.048848) (0.053299,0.047242) (0.048603,0.045597) (0.044112,0.043915) (0.039828,0.042195) (0.035755,0.040438) (0.031893,0.038643) (0.028244,0.036811) (0.02481,0.034942) (0.021593,0.033035) (0.018594,0.031091) (0.015813,0.02911) (0.013254,0.027091) (0.010915,0.025036) (0.0087994,0.022943) (0.0069068,0.020813) (0.0052381,0.018647) (0.0037941,0.016443) (0.0025752,0.014203) (0.0015817,0.011927) (0.00081385,0.0096136) (0.00027184,0.0072642) (-4.4368e-05,0.0048787) (-0.0001349,0.0024572) (0,0) (0.00035421,-0.0024353) (0.00092154,-0.004791) (0.0017014,-0.0070671) (0.0026933,-0.0092636) (0.0038964,-0.011381) (0.0053101,-0.013419) (0.0069336,-0.015378) (0.0087658,-0.017259) (0.010806,-0.019062) (0.013053,-0.020788) (0.015506,-0.022437) (0.018163,-0.024011) (0.021024,-0.025511) (0.024087,-0.026937) (0.02735,-0.02829) (0.030813,-0.029573) (0.034474,-0.030785) (0.038331,-0.031929) (0.042382,-0.033006) (0.046626,-0.034017) (0.051061,-0.034964) (0.055686,-0.035848)};
\end{groupplot}
\end{tikzpicture}

\caption{NACA 2412 modifié pour $I=2$, 4, 6, 8 (calculé par \code{naca4\_modified}). L'épaisseur maximale et sa position ne changent pas ; seul le nez, et la façon dont l'épaisseur y est atteinte, changent.}\label{fig:v2-formes}
\end{figure}

\begin{codebox}[ : \code{aero/naca4.py}, V2]
\begin{lstlisting}
def _modified_coeffs(I, T):
    d1 = (2.24 - 5.42*T + 12.3*T**2) / (10*(1 - 0.878*T))          # pente au bord de fuite
    u = 1 - T
    d2, d3 = np.linalg.solve([[u**2, u**3], [2*u, 3*u**2]], [0.1 - d1*u, -d1])
    ypp = 2*d2 + 6*d3*u                                           # courbure de l'arc arrière en T
    a0 = 0.296904 * I / 6
    M = [[T, T**2, T**3], [1, 2*T, 3*T**2], [0, 2, 6*T]]
    rhs = [0.1 - a0*np.sqrt(T), -0.5*a0/np.sqrt(T), ypp + 0.25*a0*T**-1.5]
    a1, a2, a3 = np.linalg.solve(M, rhs)
\end{lstlisting}
\end{codebox}

\section{XFOIL comme simulateur de référence}
\paragraph{Pourquoi.} Le test V5 (§\ref{sec:v5}) a montré que le modèle rapide classe bien les profils (corrélation de rang 0,91) mais surestime $C_D$ de 27~\% et déplace l'optimum. Pour un livrable, il faut le simulateur le plus juste disponible ; XFOIL coûte environ une seconde par point, et la grille V2 (1\,024 profils) se calcule en 6~minutes sur 10 cœurs. Le modèle rapide reste utile pour développer et tester, là où sa vitesse compte.

\paragraph{Le point de croisière, à portance imposée.} On impose $C_L=0{,}4$ (commande \texttt{CL} d'XFOIL) : on compare ainsi les profils à mission égale, et XFOIL renvoie l'incidence de croisière, qui est aussi l'attitude vue par le radar. Le $C_D$ correspondant est le premier objectif.

\paragraph{La contrainte de manœuvre.} Un avion doit pouvoir manœuvrer : on exige une portance maximale $C_{L,\max}\ge1{,}2$, estimée par une rampe d'incidence de 4 à $14^\circ$. Un nez aigu décroche tôt : la couche limite décolle juste après le pic d'aspiration, très intense sur un petit rayon. Dans XFOIL, cela se voit de deux façons : une portance qui plafonne puis chute, ou une non-convergence au-delà d'une certaine incidence. Dans les deux cas, $C_{L,\max}$ est pris comme la plus grande portance convergée.

\begin{figure}[htbp]
\centering
% figure générée par tools/make_figures_v2.py : v2_xfoil
\begin{tikzpicture}
\begin{groupplot}[group style={group size=2 by 1,horizontal sep=1.6cm},width=0.48\linewidth,height=5.4cm,
  xlabel={épaisseur $t$ (\%)},title style={font=\small},tick label style={font=\scriptsize},legend style={draw=none,font=\scriptsize}]
\nextgroupplot[title={$C_D$ de croisière ($C_L=0{,}4$), médiane sur $m,p$},ylabel={counts},legend pos=north west]
\addplot[corange,very thick,mark=*,mark size=1.2pt] coordinates {(8,48) (10,50.8) (12,54.1) (14,59.6) (16,61.6) (18,65.3) (20,70.3)};\addlegendentry{$I=2$}
\addplot[cgreen,very thick,mark=*,mark size=1.2pt] coordinates {(8,53.5) (10,52.9) (12,56.1) (14,58) (16,60.6) (18,64.7) (20,69.3)};\addlegendentry{$I=4$}
\addplot[cblue,very thick,mark=*,mark size=1.2pt] coordinates {(8,56) (10,55.9) (12,57.7) (14,59.7) (16,62.45) (18,65.3) (20,70.15)};\addlegendentry{$I=6$}
\addplot[cmuted,very thick,mark=*,mark size=1.2pt] coordinates {(8,55.95) (10,59.2) (12,61.25) (14,63.3) (16,67.5) (18,70.05) (20,76.35)};\addlegendentry{$I=8$}
\nextgroupplot[title={$C_{L,\max}$ (rampe 4--14°), médiane sur $m,p$},ymin=0.6]
\addplot[corange,very thick,mark=*,mark size=1.2pt] coordinates {(8,0.981) (10,1.0245) (12,1.0112) (14,1.1883) (16,1.328) (18,1.3478) (20,1.2788)};
\addplot[cgreen,very thick,mark=*,mark size=1.2pt] coordinates {(8,1.2839) (10,1.5288) (12,1.5613) (14,1.5372) (16,1.5132) (18,1.4736) (20,1.5252)};
\addplot[cblue,very thick,mark=*,mark size=1.2pt] coordinates {(8,1.6178) (10,1.6147) (12,1.609) (14,1.5857) (16,1.5264) (18,1.5132) (20,1.5011)};
\addplot[cmuted,very thick,mark=*,mark size=1.2pt] coordinates {(8,1.6741) (10,1.6447) (12,1.6326) (14,1.6006) (16,1.5582) (18,1.4953) (20,1.4599)};
\addplot[cink,dashed] coordinates {(8,1.2) (20,1.2)};
\node[font=\scriptsize,anchor=south east] at (axis cs:20,1.2) {exigence de manœuvre};
\end{groupplot}
\end{tikzpicture}

\caption{Effet du rayon de bord d'attaque sur les 1\,024 profils V2 (XFOIL, $\Re=3\cdot10^6$). Un nez aigu ($I=2$) réduit un peu la traînée de croisière, mais fait chuter la portance maximale sous l'exigence de manœuvre.}\label{fig:v2-xfoil}
\end{figure}

\paragraph{Ce que montre la grille.}
\begin{itemize}
\item À $t=12$~\%, en médiane sur la cambrure, passer de $I=8$ à $I=2$ \emph{réduit} la traînée de croisière de 61 à 54~counts : un nez plus fin limite la survitesse et le gradient adverse qui suit.
\item Mais $C_{L,\max}$ tombe de 1,63 à 1,01. Parmi les profils assez épais, 94 profils à $I=2$ ne tiennent pas l'exigence de manœuvre, contre 22 à $I=4$ et moins de 10 au-delà.
\item XFOIL ne converge pas en croisière sur 104 profils, dont 70 à $I=2$ : ils sont déclarés infaisables. Au total, 702 profils sur 1\,024 sont faisables.
\end{itemize}

\begin{projet}
C'est exactement le mécanisme qui manquait en V1. La SER veut un nez aigu (cours SER, chapitre V2), la croisière aussi, mais la manœuvre l'interdit. Le compromis traînée/SER se joue désormais sur l'épaisseur, la cambrure \emph{et} le rayon de bord d'attaque, et le front de Pareto n'est plus confiné à $t=t_{\min}$ (cours d'optimisation, chapitre V2).
\end{projet}

\begin{piege}
La non-convergence d'XFOIL n'est pas une preuve de décrochage : elle peut venir d'un bulbe laminaire mal résolu ou d'un manque d'itérations. Considérer « non convergé » comme « décroché » est une hypothèse prudente mais grossière. Un critère plus propre serait une polaire complète avec détection du maximum de $C_L$, ou un calcul RANS sur quelques profils de contrôle.
\end{piege}

\exercices
\begin{exercice}
Pour $t=12$~\%, calculer le rayon de bord d'attaque pour $I=2$, 4, 6 et 8. De combien de décibels la SER de nez (loi $\pi r_{BA}$) diminue-t-elle entre $I=8$ et $I=2$ ?
\end{exercice}

\begin{pistes}
\begin{itemize}
\item Libérer aussi la position de l'épaisseur maximale $T$ (cinquième variable), qui agit sur la position de la transition, donc sur la traînée.
\item Remplacer « non convergé = décroché » par une détection propre du maximum de $C_L$ sur une polaire fine, et valider sur quelques profils par RANS.
\item Calibrer le modèle rapide sur XFOIL (seuil de Michel, décambrage visqueux) pour disposer d'un simulateur basse fidélité fiable, utile à une approche multi-fidélité.
\end{itemize}
\end{pistes}


\appendix
\chapter{Corrigés des exercices}\label{annexe:corriges}
\iterbanner{V0}

\section*{Chapitre~\ref{ch:intro} — Rôle du modèle}
\textbf{Ex.~1.1.} (a) $40\times50\times10=2\cdot10^4$ appels au plus (un peu plus avec la population initiale). (b) À 8~ms : 160~s, moins de 3~minutes. Avec un RANS à 30~min : $10^4$~heures, plus d'un an de calcul séquentiel. (c) La grille ne contient que 512 profils, donc au plus 512 simulations distinctes ; NSGA-II réévalue en permanence les mêmes individus (élitisme, croisements qui redonnent des parents), et le cache d'\code{AeroSimulator} ne compte chaque profil qu'une fois.

\textbf{Ex.~1.2.} (1) Une énumération qui contient des trous n'est plus exhaustive : le front « exact » ne l'est plus, et il faut décider quoi faire des profils manquants. (2) Les échecs ne sont pas aléatoires : ils se concentrent sur les formes extrêmes, donc biaisent l'échantillon et l'ajustement quadratique. (3) Un échec silencieux, ou une convergence partielle, injecte des valeurs fausses ou bruitées que les moindres carrés ne savent pas distinguer des bonnes.

\section*{Chapitre~\ref{ch:naca} — Géométrie}
\textbf{Ex.~2.1.} (a) Avant $p$ : $y_c(p)=\frac{m}{p^2}(2p^2-p^2)=m=0{,}02$. Après $p$ : $\frac{m}{(1-p)^2}(1-2p+2p^2-p^2)=\frac{m}{(1-p)^2}(1-p)^2=m$. Les deux pentes valent $\frac{2m}{\cdot}(p-p)=0$ : raccord $C^1$. (b) $r_{BA}=1{,}1019\times0{,}12^2=0{,}0159\,c$. (c) $y_c''=-2m/p^2=-0{,}25$ avant, $-2m/(1-p)^2=-0{,}111$ après : la courbure saute en $x=p$.

\textbf{Ex.~2.2.} Cosinus : $x_1=\tfrac12(1-\cos\tfrac{\pi}{100})=2{,}47\cdot10^{-4}$. Uniforme : $10^{-2}$, quarante fois plus grand. Les points du contour sont des \emph{extrémités} de panneaux ; les collocations sont les milieux. Le point de bord d'attaque est l'extrémité commune de deux panneaux, il n'est donc jamais un point de collocation (de même pour le bord de fuite). C'est pourquoi le point d'arrêt est interpolé entre deux collocations.

\section*{Chapitre~\ref{ch:potentiel} — Écoulement potentiel}
\textbf{Ex.~3.1.} $\Gamma=\tfrac12C_LU_\infty c=\tfrac12\times0{,}5\times50\times1=12{,}5$~m\textsuperscript{2}/s. $L'=\rho U_\infty\Gamma=1{,}2\times50\times12{,}5=750$~N/m, ce qu'on retrouve par $\tfrac12\rho U_\infty^2cC_L$.

\textbf{Ex.~3.2.} Le rectangle a pour longueur $\Delta s$ le long de la nappe et une hauteur qui tend vers 0. Le débit sortant par les côtés courts tend vers 0 ; par les grands côtés, il vaut $(v_+-v_-)\Delta s$. Il doit égaler le débit des sources contenues, $\sigma\Delta s$ : $v_+-v_-=\sigma$. Pour une nappe plane isolée, le problème est symétrique par réflexion à travers la nappe, donc $v_-=-v_+$, d'où $\pm\sigma/2$.

\section*{Chapitre~\ref{ch:panneaux} — Hess-Smith}
\textbf{Ex.~4.1.} Une primitive de $\eta/((\xi'-\xi)^2+\eta^2)$ est $\arctan\frac{\xi'-\xi}{\eta}$, d'où $\int_0^L=\arctan\frac{L-\xi}{\eta}+\arctan\frac{\xi}{\eta}$. Pour $\eta>0$, $\theta_1=\tfrac\pi2-\arctan\frac\xi\eta$ et $\theta_2=\tfrac\pi2+\arctan\frac{L-\xi}{\eta}$, donc $\theta_2-\theta_1$ est bien cette somme : c'est l'angle sous lequel on voit le panneau depuis $P$. Très loin, cet angle vaut environ $L\,\eta/r^2$, et $v_\sigma\approx\frac{\sigma L}{2\pi r}\frac{\eta}{r}$ : la composante normale d'une source ponctuelle de débit $\sigma L$. Vu de loin, le panneau est un point.

\textbf{Ex.~4.2.} (a) Inversion : de l'ordre de $N^3\approx8\cdot10^6$ opérations, environ une milliseconde. Nouvelle incidence : un produit matrice-vecteur, $N^2\approx4\cdot10^4$ opérations, deux cents fois moins. La construction des matrices d'influence, en $O(N^2)$ avec des fonctions transcendantes, coûte plus que ce produit. (b) Avec le cache, seule la première évaluation construit et inverse la matrice ; les 5 à 8 suivantes de la recherche de Brent ne coûtent qu'un produit matrice-vecteur et l'intégration des efforts.

\section*{Chapitre~\ref{ch:cl} — Couche limite intégrale}
\textbf{Ex.~5.1.} Profil linéaire : $\delta^*/\delta=\int_0^1(1-\eta)\dd\eta=\tfrac12$, $\theta/\delta=\int_0^1\eta(1-\eta)\dd\eta=\tfrac16$, $H=3$. Profil en $1/7$ : $\delta^*/\delta=1-\tfrac78=\tfrac18$, $\theta/\delta=\tfrac78-\tfrac79=\tfrac7{72}$, $H=\tfrac97\approx1{,}29$. Le premier, peu rempli, ressemble à un profil laminaire (Blasius : 2,59) ; le second, très plein près de la paroi, à un profil turbulent.

\textbf{Ex.~5.2.} À incidence positive, l'écoulement amont arrive par en dessous du profil, et la circulation horaire accélère l'extrados et ralentit l'intrados. La ligne de courant qui se partage entre les deux faces vient donc frapper le profil sous le bord d'attaque géométrique, sur l'intrados. Quand $\alpha$ augmente, le point d'arrêt recule le long de l'intrados ($x\approx0{,}001$ pour le 2412 à $2^\circ$).

\section*{Chapitre~\ref{ch:thwaites} — Thwaites}
\textbf{Ex.~6.1.} $\int_0^sk^5s'^5\dd s'=k^5s^6/6$, donc $\theta^2=\frac{0{,}45\,\nu}{k^6s^6}\frac{k^5s^6}{6}=0{,}075\,\frac{\nu}{k}$, constant. $\lambda=\theta^2k/\nu=0{,}075$. $H(0{,}075)=2{,}61-3{,}75\times0{,}075+5{,}24\times0{,}075^2=2{,}36$.

\textbf{Ex.~6.2.} $\frac{U_e}{\nu}\frac{\dd\theta^2}{\dd s}=a-b\frac{\theta^2U_e'}{\nu}$ s'écrit $U_e(\theta^2)'+b\,\theta^2U_e'=a\nu$. En multipliant par $U_e^{b-1}$, le membre de gauche devient $\frac{\dd}{\dd s}(\theta^2U_e^b)$, d'où $\theta^2U_e^b=a\nu\int_0^sU_e^{b-1}$. Avec $b=6$, on retrouve la puissance 5.

\section*{Chapitre~\ref{ch:transition} — Transition}
\textbf{Ex.~7.1.} L'équation est $0{,}671\sqrt{\Re_x}=1{,}174\,(1+22\,400/\Re_x)\,\Re_x^{0{,}46}$. Une méthode de dichotomie ou de Newton donne $\Re_x\approx1{,}66\cdot10^6$ et $\Re_\theta\approx864$. À $\Re=3\cdot10^6$, $x_{tr}\approx0{,}55$.

\textbf{Ex.~7.2.} Sans interpolation, $x_{tr}$ ne prend que les abscisses des stations : quand $t$ varie, il reste constant puis saute d'une station à la suivante. $C_D(t)$ est alors une fonction en escalier, avec des marches d'environ 1~\%. Pour les moindres carrés, ces marches sont un bruit qui dégrade le $R^2$ et peut déplacer le minimum du modèle. Pour un optimiseur par gradient, la dérivée est nulle presque partout et infinie aux marches, et les différences finies deviennent erratiques.

\section*{Chapitre~\ref{ch:head} — Head}
\textbf{Ex.~8.1.} $H_1(1{,}3)=9{,}83$, $H_1(1{,}4)=7{,}18$, $H_1(2{,}4)=3{,}59$. $F(7{,}18)=0{,}0306\times4{,}18^{-0{,}6169}\approx0{,}0127$. Pour $\theta=10^{-3}c$ et $H=1{,}4$ : $\delta-\delta^*=H_1\theta\approx7{,}2\cdot10^{-3}c$.

\textbf{Ex.~8.2.} $\frac{\dd\ln H_1}{\dd s}=\frac{\dd\ln E}{\dd s}-\frac{\dd\ln(U_e\theta)}{\dd s}$. Avec le système de Head : $\frac{\dd\ln E}{\dd s}=\frac{F}{\theta H_1}$ et $\frac{\dd(U_e\theta)}{\dd s}=U_e\frac{c_f}{2}-(H+1)\theta U_e'$. D'où
\[
\frac{\dd\ln H_1}{\dd s}=\frac{F}{\theta H_1}-\frac{c_f}{2\theta}+(H+1)\frac{U_e'}{U_e}.
\]
En gradient adverse, le dernier terme est négatif et domine : $H_1$ diminue. Comme $H_1$ est une fonction décroissante de $H$, $H$ augmente.

\section*{Chapitre~\ref{ch:trainee} — Squire-Young}
\textbf{Ex.~9.1.} Extrados : exposant $(1{,}755+5)/2=3{,}377$, $0{,}905^{3{,}377}=0{,}714$, contribution $2\times3{,}362\cdot10^{-3}\times0{,}714=48{,}0$~counts. Intrados : exposant 3,252, $0{,}873^{3{,}252}=0{,}642$, contribution $21{,}2$~counts. Total : $69{,}2$~counts, la valeur du code.

\textbf{Ex.~9.2.} Avec $H$ constant, $\ln(\theta_\infty/\theta_{TE})=-(H_{TE}+2)\int_{u_{TE}}^0\dd u=(H_{TE}+2)\,u_{TE}$ : l'exposant devient $H_{TE}+2$ (3,75 au lieu de 3,38 pour $H_{TE}=1{,}75$), et la traînée est sous-estimée. L'hypothèse ignore que le sillage se remplit ($H\to1$) : le facteur $(H+2)$ qui fait décroître $\theta$ diminue en réalité vers l'aval.

\section*{Chapitre~\ref{ch:modele} — Assemblage}
\textbf{Ex.~10.1.} $\alpha=\alpha_{L0}+C_L/(\dd C_L/\dd\alpha)=-2{,}13^\circ+\frac{0{,}4}{6{,}88}\,\text{rad}=-2{,}13^\circ+3{,}33^\circ=1{,}20^\circ$, contre $1{,}19^\circ$ pour \code{alpha\_for\_cl} : $C_L(\alpha)$ est bien quasi linéaire.

\textbf{Ex.~10.2.} \code{up:turbulent\_separation} : sur l'extrados, $H$ a atteint 2,4 dans la méthode de Head, en $x_{sep}$. \code{te\_separation} : ce décollement est en aval de $x_{sep,crit}=0{,}9$, donc toléré comme artefact probable du bord de fuite fermé. Le résultat devient invalide si le décollement remonte en amont de $x/c=0{,}9$ (incidence plus forte, gradient adverse plus marqué), ou si l'intégration échoue.

\section*{Chapitre~\ref{ch:espace} — Espace de conception}
\textbf{Ex.~11.1.} (a) $4\to(0,0,1)$, $3\to(1,1,0)$, $1\to(1,0,0)$, soit \texttt{001110100}. Valeurs : $m=4$~\%, $p=0{,}20+3\times0{,}05=0{,}35$, $t=6+2=8$~\%. (b) $m=0$ impose les trois premiers bits à 0 ; restent $8\times8=64$ chaînes, pour seulement 8 profils distincts (un par épaisseur) : chaque profil symétrique est représenté par 8 chaînes, une par valeur de $p$.

\textbf{Ex.~11.2.} (1) Le temps dépend de l'implémentation et du matériel : QAOA simulé sur un ordinateur classique est lent par construction, et une machine réelle ajoute des files d'attente. (2) Dans un cas industriel, une simulation RANS coûte des heures et domine tout le reste ; le nombre de simulations distinctes est la grandeur qui se transpose à l'échelle, pas le temps du démonstrateur.

\section*{Chapitre~\ref{ch:validation} — Validation}
\textbf{Ex.~12.1.} A. Un optimiseur ne compare que des profils entre eux : un biais commun ne change pas le classement, donc pas le front. Une corrélation de rang de 0,97 garantit que les meilleurs profils restent les meilleurs. B, plus juste en moyenne, se trompe souvent sur l'ordre, et mènerait à de mauvais choix.

\textbf{Ex.~12.2.} Près du bord de fuite, la couche limite de l'extrados, sous gradient adverse, est bien plus épaisse que celle de l'intrados : $\delta^*_{up}>\delta^*_{lo}$. Le profil « vu » par l'écoulement extérieur, déplacé de $\delta^*$, a donc un bord de fuite relevé : c'est comme un volet braqué vers le haut, qui décambre le profil. La condition de Kutta, appliquée à ce bord de fuite effectif, donne moins de circulation. La vitesse de transpiration, positive (soufflage) là où $U_e\delta^*$ croît, c'est-à-dire surtout en fin d'extrados, reproduit cet effet.

\section*{Chapitre~\ref{ch:v2} \iterbadge{V2}}
\textbf{Ex.~13.1.} $r_{BA}=1{,}1019\,(0{,}12\,I/6)^2$ : 1,8~mm ($I=2$), 7,1~mm ($I=4$), 15,9~mm ($I=6$), 28,2~mm ($I=8$) pour une corde de 1~m. La loi $\pi r_{BA}$ prédit $10\log_{10}(r_8/r_2)=10\log_{10}16=12$~dB de moins. Le modèle d'optique physique donne un peu moins (environ 7~dB entre $I=6$ et $I=2$ au lieu des 9,5~dB prévus) : à 10~GHz, $k\,r_{BA}=0{,}37$ pour $I=2$, bien en dessous du régime optique, et la loi $\pi\rho$ n'est plus qu'indicative (cours SER).

\chapter{Aide-mémoire}\label{annexe:memo}
\iterbanner{V0}

\begin{center}
\small
\renewcommand{\arraystretch}{1.45}
\begin{tabularx}{\linewidth}{>{\bfseries}l X}
\toprule
NACA & $y_c=\frac{m}{p^2}(2px-x^2)$ ($x<p$), $\frac{m}{(1-p)^2}(1-2p+2px-x^2)$ ($x\ge p$) ; $y_t=5t(0{,}2969\sqrt x-0{,}126x-0{,}3516x^2+0{,}2843x^3-0{,}1036x^4)$ ; $r_{BA}=1{,}1019\,t^2$ ; $x_k=\frac12(1-\cos\frac{k\pi}{n})$\\
Panneau (local) & source : $u=\frac{\sigma}{2\pi}\ln\frac{r_1}{r_2}$, $v=\frac{\sigma}{2\pi}(\theta_2-\theta_1)$ ; tourbillon : $u=-\frac{\gamma}{2\pi}(\theta_2-\theta_1)$, $v=\frac{\gamma}{2\pi}\ln\frac{r_1}{r_2}$ ; auto-influence $\theta_2-\theta_1=-\pi$\\
Hess-Smith & $N$ sources $\sigma_j$ + 1 tourbillon $\gamma$ ; $\vect V\cdot\vect n_i=0$ ; Kutta $V_{t,1}+V_{t,N}=0$ ; $\vect n=(t_y,-t_x)$ sortante\\
Sorties & $C_p=1-V_t^2$ ; $\vect F=-\sum C_p\vect nL$ ; $C_L=F_y\cos\alpha-F_x\sin\alpha$ ; $C_{L,KJ}=-2\gamma\,\ell_{tot}$ ; $C_{D,p}\to0$\\
Von Kármán & $\theta'+(2+H)\frac{\theta}{U_e}U_e'=\frac{c_f}{2}$ ; $H=\delta^*/\theta$\\
Thwaites & $\theta^2=\frac{0{,}45\,\nu}{U_e^6}\int_0^sU_e^5$ ; $\lambda=\frac{\theta^2}{\nu}U_e'$ ; $c_f=2\ell(\lambda)/\Re_\theta$ ; décollement $\lambda=-0{,}09$\\
Transition & premier de : Michel $\Re_\theta>1{,}174(1+\frac{22\,400}{\Re_x})\Re_x^{0{,}46}$ ; $\lambda<-0{,}09$ ; $x_{trip}$ ; position interpolée\\
Head & $\theta'=\frac{c_f}{2}-(H+2)\frac{\theta}{U_e}U_e'$ ; $(U_e\theta H_1)'=U_e\,0{,}0306(H_1-3)^{-0{,}6169}$ ; $c_f=0{,}246\cdot10^{-0{,}678H}\Re_\theta^{-0{,}268}$ ; départ $H=1{,}4$ ; décollement $H=2{,}4$\\
Squire-Young & $C_D=\sum_{\text{faces}}2\,\theta_{TE}\,U_{e,TE}^{(H_{TE}+5)/2}$, appliqué en $x_{stop}=0{,}985$\\
Validité & invalide si décollement turbulent en amont de $x/c=0{,}9$, point d'arrêt ambigu, ou échec d'intégration\\
\midrule
API & \code{evaluate(m, p, t, alpha, Re)} $\to$ \code{AeroResult} ; \code{evaluate\_at\_cl(m, p, t, CL)} ; \code{alpha\_for\_cl} ; \code{naca4(m, p, t, n\_half)} ; \code{solve(X, Y, alpha, cache)}\\
Grille & \code{DesignSpace} : $8^3$ niveaux, 9 bits, \code{encode}/\code{decode} (poids faible en premier) ; \code{AeroSimulator(space, case)} : cache, \code{n\_sims}, \code{constraints}\\
Scripts & \code{demo\_profil.py} ; \code{validation.py} (V1--V5) ; \code{generer\_base.py [cl]} ; \code{analyse\_quadratique.py}\\
Ordres de grandeur & 8~ms/profil ; NACA 2412, $2^\circ$, $\Re=3\cdot10^6$ : $C_L=0{,}498$, $C_D=69$~counts (55 frottement + 14 forme), $C_M=-0{,}057$, $x_{tr}=0{,}28/0{,}48$\\
\bottomrule
\end{tabularx}
\end{center}

\chapter{Pistes d'amélioration, chapitre par chapitre}\label{annexe:pistes}
\iterbanner{V1}

Le chapitre~\ref{ch:validation} donne les pistes globales. Cette annexe les détaille partie par partie, de la plus utile à la moins utile pour le hackathon.

\section*{Géométrie (ch.~\ref{ch:naca})}
\begin{itemize}
\item Ajouter une paramétrisation qui découple le rayon de bord d'attaque de l'épaisseur (NACA 4 chiffres modifiée, CST de Kulfan) : c'est la condition pour que traînée et SER entrent vraiment en conflit (cours d'optimisation, chapitre sur le front exact).
\item Prévoir un bord de fuite d'épaisseur finie en option, plus réaliste pour la fabrication.
\end{itemize}

\section*{Panneaux (ch.~\ref{ch:potentiel} et \ref{ch:panneaux})}
\begin{itemize}
\item Passer à des tourbillons linéaires par panneau (formulation d'XFOIL) : erreur en $1/N^2$ au lieu de $1/N$.
\item Ajouter une correction de compressibilité de Kármán-Tsien sur $C_p$, pour un drone à Mach 0,5 à 0,6.
\end{itemize}

\section*{Couche limite (ch.~\ref{ch:cl} à \ref{ch:head})}
\begin{itemize}
\item Remplacer Michel par un critère qui tient compte du gradient de pression, ou par une méthode $e^N$ simplifiée (enveloppe de Drela et Giles), et caler sur XFOIL.
\item Remplacer Head par la méthode à décalage de Green (\emph{lag-entrainment}), plus juste en gradient adverse.
\item Option « transition forcée à 5~\% » : cas tout turbulent, plus régulier, donc plus favorable au modèle quadratique de la phase~4.
\end{itemize}

\section*{Traînée (ch.~\ref{ch:trainee})}
\begin{itemize}
\item Remplacer la coupure \code{x\_stop} par une itération de couplage avec vitesse de transpiration, qui efface physiquement le point d'arrêt du bord de fuite.
\end{itemize}

\section*{Assemblage, espace de conception (ch.~\ref{ch:modele} et \ref{ch:espace})}
\begin{itemize}
\item \iterbadge{V1}~\textbf{Fait} : intégration avec un modèle SER simplifié et une chaîne d'optimisation complète (\code{Code/optim/}).
\item Paralléliser l'énumération (512 profils indépendants) pour un passage à des grilles plus fines.
\end{itemize}

\section*{Validation (ch.~\ref{ch:validation})}
\begin{itemize}
\item \iterbadge{V1}~\textbf{Fait} : XFOIL installé localement (\code{Code/.tools/}), test V5 étendu à toute la grille (§\ref{sec:v5}). Reste à caler la transition sur XFOIL.
\item Corriger aussi la portance (décambrage visqueux) : notre $C_L$ est surestimé de 25~\% en moyenne à $2^\circ$, ce qui rend la contrainte $C_L\ge0{,}3$ trop permissive (392 profils faisables contre 359 avec XFOIL).
\end{itemize}


\end{document}
